% Options for packages loaded elsewhere
\PassOptionsToPackage{unicode}{hyperref}
\PassOptionsToPackage{hyphens}{url}
\documentclass[
]{article}
\usepackage{xcolor}
\usepackage{amsmath,amssymb}
\setcounter{secnumdepth}{-\maxdimen} % remove section numbering
\usepackage{iftex}
\ifPDFTeX
  \usepackage[T1]{fontenc}
  \usepackage[utf8]{inputenc}
  \usepackage{textcomp} % provide euro and other symbols
\else % if luatex or xetex
  \usepackage{unicode-math} % this also loads fontspec
  \defaultfontfeatures{Scale=MatchLowercase}
  \defaultfontfeatures[\rmfamily]{Ligatures=TeX,Scale=1}
\fi
\usepackage{lmodern}
\ifPDFTeX\else
  % xetex/luatex font selection
\fi
% Use upquote if available, for straight quotes in verbatim environments
\IfFileExists{upquote.sty}{\usepackage{upquote}}{}
\IfFileExists{microtype.sty}{% use microtype if available
  \usepackage[]{microtype}
  \UseMicrotypeSet[protrusion]{basicmath} % disable protrusion for tt fonts
}{}
\makeatletter
\@ifundefined{KOMAClassName}{% if non-KOMA class
  \IfFileExists{parskip.sty}{%
    \usepackage{parskip}
  }{% else
    \setlength{\parindent}{0pt}
    \setlength{\parskip}{6pt plus 2pt minus 1pt}}
}{% if KOMA class
  \KOMAoptions{parskip=half}}
\makeatother
\usepackage{color}
\usepackage{fancyvrb}
\newcommand{\VerbBar}{|}
\newcommand{\VERB}{\Verb[commandchars=\\\{\}]}
\DefineVerbatimEnvironment{Highlighting}{Verbatim}{commandchars=\\\{\}}
% Add ',fontsize=\small' for more characters per line
\newenvironment{Shaded}{}{}
\newcommand{\AlertTok}[1]{\textcolor[rgb]{1.00,0.00,0.00}{\textbf{#1}}}
\newcommand{\AnnotationTok}[1]{\textcolor[rgb]{0.38,0.63,0.69}{\textbf{\textit{#1}}}}
\newcommand{\AttributeTok}[1]{\textcolor[rgb]{0.49,0.56,0.16}{#1}}
\newcommand{\BaseNTok}[1]{\textcolor[rgb]{0.25,0.63,0.44}{#1}}
\newcommand{\BuiltInTok}[1]{\textcolor[rgb]{0.00,0.50,0.00}{#1}}
\newcommand{\CharTok}[1]{\textcolor[rgb]{0.25,0.44,0.63}{#1}}
\newcommand{\CommentTok}[1]{\textcolor[rgb]{0.38,0.63,0.69}{\textit{#1}}}
\newcommand{\CommentVarTok}[1]{\textcolor[rgb]{0.38,0.63,0.69}{\textbf{\textit{#1}}}}
\newcommand{\ConstantTok}[1]{\textcolor[rgb]{0.53,0.00,0.00}{#1}}
\newcommand{\ControlFlowTok}[1]{\textcolor[rgb]{0.00,0.44,0.13}{\textbf{#1}}}
\newcommand{\DataTypeTok}[1]{\textcolor[rgb]{0.56,0.13,0.00}{#1}}
\newcommand{\DecValTok}[1]{\textcolor[rgb]{0.25,0.63,0.44}{#1}}
\newcommand{\DocumentationTok}[1]{\textcolor[rgb]{0.73,0.13,0.13}{\textit{#1}}}
\newcommand{\ErrorTok}[1]{\textcolor[rgb]{1.00,0.00,0.00}{\textbf{#1}}}
\newcommand{\ExtensionTok}[1]{#1}
\newcommand{\FloatTok}[1]{\textcolor[rgb]{0.25,0.63,0.44}{#1}}
\newcommand{\FunctionTok}[1]{\textcolor[rgb]{0.02,0.16,0.49}{#1}}
\newcommand{\ImportTok}[1]{\textcolor[rgb]{0.00,0.50,0.00}{\textbf{#1}}}
\newcommand{\InformationTok}[1]{\textcolor[rgb]{0.38,0.63,0.69}{\textbf{\textit{#1}}}}
\newcommand{\KeywordTok}[1]{\textcolor[rgb]{0.00,0.44,0.13}{\textbf{#1}}}
\newcommand{\NormalTok}[1]{#1}
\newcommand{\OperatorTok}[1]{\textcolor[rgb]{0.40,0.40,0.40}{#1}}
\newcommand{\OtherTok}[1]{\textcolor[rgb]{0.00,0.44,0.13}{#1}}
\newcommand{\PreprocessorTok}[1]{\textcolor[rgb]{0.74,0.48,0.00}{#1}}
\newcommand{\RegionMarkerTok}[1]{#1}
\newcommand{\SpecialCharTok}[1]{\textcolor[rgb]{0.25,0.44,0.63}{#1}}
\newcommand{\SpecialStringTok}[1]{\textcolor[rgb]{0.73,0.40,0.53}{#1}}
\newcommand{\StringTok}[1]{\textcolor[rgb]{0.25,0.44,0.63}{#1}}
\newcommand{\VariableTok}[1]{\textcolor[rgb]{0.10,0.09,0.49}{#1}}
\newcommand{\VerbatimStringTok}[1]{\textcolor[rgb]{0.25,0.44,0.63}{#1}}
\newcommand{\WarningTok}[1]{\textcolor[rgb]{0.38,0.63,0.69}{\textbf{\textit{#1}}}}
\usepackage{longtable,booktabs,array}
\usepackage{calc} % for calculating minipage widths
% Correct order of tables after \paragraph or \subparagraph
\usepackage{etoolbox}
\makeatletter
\patchcmd\longtable{\par}{\if@noskipsec\mbox{}\fi\par}{}{}
\makeatother
% Allow footnotes in longtable head/foot
\IfFileExists{footnotehyper.sty}{\usepackage{footnotehyper}}{\usepackage{footnote}}
\makesavenoteenv{longtable}
\setlength{\emergencystretch}{3em} % prevent overfull lines
\providecommand{\tightlist}{%
  \setlength{\itemsep}{0pt}\setlength{\parskip}{0pt}}
\usepackage{bookmark}
\IfFileExists{xurl.sty}{\usepackage{xurl}}{} % add URL line breaks if available
\urlstyle{same}
\hypersetup{
  hidelinks,
  pdfcreator={LaTeX via pandoc}}

\author{}
\date{}


\usepackage[a4paper, top=22mm, bottom=22mm, left=18mm, right=18mm, headheight=16pt, footskip=12mm]{geometry}
\usepackage{fancyhdr}
\usepackage{lastpage}
\usepackage{microtype}
\usepackage{tcolorbox}
\usepackage{fvextra}
\usepackage{booktabs}
\usepackage{longtable}
\usepackage{array}
\usepackage{tocloft}
\usepackage{pdflscape}
\usepackage[htt]{hyphenat}
\usepackage{xurl}

% Colors
\definecolor{primaryindigo}{RGB}{49, 46, 129}
\definecolor{headerblue}{RGB}{30, 41, 59}
\definecolor{codebg}{RGB}{248, 250, 252}
\definecolor{codeborder}{RGB}{226, 232, 240}
\definecolor{notebg}{RGB}{238, 242, 255}
\definecolor{noteborder}{RGB}{99, 102, 241}

% Hyperref setup
\hypersetup{
    colorlinks=true,
    linkcolor=primaryindigo,
    urlcolor=primaryindigo,
    citecolor=primaryindigo
}

% Page Header and Footer Styling
\pagestyle{fancy}
\fancyhf{}
\fancyhead[L]{\textsf{\footnotesize \textbf{\textcolor{primaryindigo}{StudyRAG: Technical Architecture \& Implementation Manual}}}}
\fancyhead[R]{\textsf{\footnotesize \textcolor{headerblue}{System Handover Reference}}}
\fancyfoot[L]{\textsf{\footnotesize \textcolor{headerblue}{Version 1.0.0 (Release Candidate) | Open Architecture}}}
\fancyfoot[R]{\textsf{\footnotesize \textbf{\thepage} of \pageref{LastPage}}}
\renewcommand{\headrulewidth}{0.6pt}
\renewcommand{\footrulewidth}{0.4pt}

% Section Heading Styling
\usepackage{titlesec}
\titleformat{\section}
  {\normalfont\Large\bfseries\color{primaryindigo}}{\thesection}{1em}{}[\titlerule]
\titleformat{\subsection}
  {\normalfont\large\bfseries\color{headerblue}}{\thesubsection}{1em}{}
\titleformat{\subsubsection}
  {\normalfont\normalsize\bfseries\color{headerblue}}{\thesubsubsection}{1em}{}

% Spacing & Breakability
\setlength{\parindent}{0pt}
\setlength{\parskip}{6pt plus 2pt minus 1pt}
\setlength{\tabcolsep}{3.5pt}
\renewcommand{\arraystretch}{1.25}
\emergencystretch=3em
\sloppy

% Verbatim and Code Wrapping
\DefineVerbatimEnvironment{Highlighting}{Verbatim}{
  breaklines=true,
  breakanywhere=true,
  fontsize=\footnotesize,
  commandchars=\\\{\}
}


\begin{document}

\begin{titlepage}
\centering
\vspace*{2.0cm}

{\Huge \bfseries \color{primaryindigo} StudyRAG: AI-Powered Study Assistant\par}
\vspace{0.5cm}
{\Large \bfseries \color{headerblue} Complete, End-to-End System Technical Architecture,\par Implementation \& Handover Reference Manual\par}

\vspace{1.5cm}
{\color{primaryindigo}\rule{0.9\textwidth}{1.5pt}}\par
\vspace{1.0cm}

\begin{flushleft}
\large
\begin{tabular}{rp{10.5cm}}
  \textbf{Author \& System Architect:} & Technical Architecture Team / Project Maintainer \\[7pt]
  \textbf{Repository:} & \url{https://github.com/kirit-somani/Rag-Based-Study-Application-} \\[7pt]
  \textbf{Runtime Environment:} & Python 3.12+ | Streamlit | LangChain | PyTorch | PyMuPDF \\[7pt]
  \textbf{Target Operating Systems:} & macOS (Darwin arm64/x86\_64), Linux (Ubuntu 22.04+), Windows 11 \\[7pt]
  \textbf{Documentation Version:} & 1.0.0 (Release Candidate \& Production Handover) \\[7pt]
  \textbf{Verification Date:} & October 14, 2026 \\[7pt]
  \textbf{Automated Test Status:} & {\color{primaryindigo}\textbf{Verified \& Code-Audited (28 / 28 Automated Tests Passing)}} \\
\end{tabular}
\end{flushleft}

\vspace{1.5cm}
{\color{primaryindigo}\rule{0.9\textwidth}{1.5pt}}\par

\vfill
{\small \textsf{\color{headerblue} Department of Cognitive AI Systems \& Educational Software Engineering}\par}
\vspace{0.8cm}
\end{titlepage}

\newpage
\tableofcontents
\newpage



\section{Executive Summary}\label{executive-summary}

This master document serves as the exhaustive single-source-of-truth
technical handover manual for \textbf{StudyRAG}, an enterprise-grade
Retrieval-Augmented Generation (RAG) study companion and academic
productivity system. It details the complete system lifecycle from
theoretical motivation, algorithmic formulations, and deep architectural
design patterns to raw source code implementations, database
persistence, automated test suites, deployment procedures, and
troubleshooting diagnostics.

\section{Chapter 1: Project Overview}\label{chapter-1-project-overview}

\subsection{1.1 Project Title, Mission, and Executive
Summary}\label{project-title-mission-and-executive-summary}

\textbf{Project Title:} StudyRAG: Intelligent Retrieval-Augmented
Generation Study Partner\\
\textbf{Application Entry Point:} \texttt{app.\allowbreak py}\\
\textbf{System Architecture:} Modular, Local-First / Hybrid-Cloud RAG
Architecture

\textbf{Executive Summary:}\\
StudyRAG is an open-source, academically rigorous Retrieval-Augmented
Generation application designed to transform static study
materials---including university course textbooks, professor lecture
slide decks, syllabus outlines, and technical research papers (in PDF
format)---into an active, cited, interactive learning ecosystem.

Unlike generic conversational AI chatbots that hallucinate facts or draw
upon unverified web corpora, StudyRAG strictly grounds its answers,
practice quizzes, digital flashcard decks, and revision summaries within
the user's uploaded course documents. Every assertion returned by the
system includes page-level citations and source filenames, enabling
students and educators to verify statements directly against course
literature.

\subsection{1.2 Educational Background and Theoretical
Motivation}\label{educational-background-and-theoretical-motivation}

Traditional academic studying often suffers from the ``illusion of
competence.'' Students reread highlight-laden textbook chapters
passively, mistakenly believing that familiarity equates to mastery.
Decades of cognitive psychology and educational science---most notably
the work of Roediger \& Karpicke (2006) on the \emph{Testing Effect} and
Bjork \& Bjork (2011) on \emph{Desirable Difficulties}---demonstrate
that long-term conceptual retention requires: 1. \textbf{Active Recall}:
Actively stimulating memory retrieval through testing and problem
generation rather than passive reading. 2. \textbf{Spaced Retrieval
Practice}: Interleaved quizzing and flashcard memorization distributed
over study sessions. 3. \textbf{Targeted Concept Clarification}:
Immediate feedback on ambiguous questions with precise, authoritative
citations.

Creating high-quality multiple-choice questions, flashcards, and chapter
summaries manually is an exhausting, time-consuming process that often
consumes more study time than the learning itself. StudyRAG automates
this cognitive scaffolding by applying state-of-the-art dense semantic
retrieval and modern LLMs to the student's exact syllabus.

\subsection{1.3 Problem Statement: The Pitfalls of Conventional LLMs in
Academic
Study}\label{problem-statement-the-pitfalls-of-conventional-llms-in-academic-study}

When students turn to commercial general-purpose LLMs (such as base
ChatGPT, Copilot, or Claude) for exam preparation, they encounter
critical failure modes: 1. \textbf{Hallucination of Technical Facts}:
LLMs generate plausible-sounding but erroneous mathematical derivations,
historical dates, scientific mechanisms, and code logic. 2.
\textbf{Out-of-Syllabus Bias}: General LLMs draw upon internet
discussions that may contradict the specific definitions, notations, or
textbook conventions mandated by a course professor. 3. \textbf{Absence
of Source Attribution}: Generic chatbots cannot provide verifiable
page-level references to physical textbooks or course lecture PDFs. 4.
\textbf{Context Window Saturation and Cost}: Uploading massive 500-page
textbooks directly into LLM prompts exceeds token limits, degrades
context retention (the ``needle-in-a-haystack'' degradation), and incurs
prohibitive inference costs. 5. \textbf{UI Latency Freezes}: Generating
multi-question practice tests with slow LLMs blocks user interfaces,
disrupting study flow.

\subsection{1.4 The StudyRAG Solution
Architecture}\label{the-studyrag-solution-architecture}

StudyRAG resolves these limitations through a dedicated, local-first RAG
pipeline: - \textbf{PyMuPDF Document Extraction}: Ingests multi-page
academic PDFs with high-fidelity coordinate and text extraction. -
\textbf{Semantic Text Cleaning \& Recursive Chunking}: Strips extraneous
whitespace, normalizes paragraphs, and partitions text using semantic
boundaries
(\texttt{{[}"\textbackslash{}n\textbackslash{}n",\allowbreak \ \allowbreak "\textbackslash{}n",\allowbreak \ \allowbreak ".\ \allowbreak ",\allowbreak \ \allowbreak "\ \allowbreak ",\allowbreak \ \allowbreak ""{]}})
into overlapping chunks (500 characters with 50-character overlap) that
preserve contextual continuity across chunk borders. - \textbf{Local
Dense Embeddings}: Employs
\texttt{sentence-transformers/\allowbreak all-MiniLM-L6-v2} (384-dimensional dense
vectors) to convert text into semantic vector embeddings without sending
document text to external third-party embedding APIs. - \textbf{Dual
Vector Database Backends}: Supports \textbf{FAISS} (Facebook AI
Similarity Search) for blazing-fast in-memory L2 index searches,
alongside \textbf{ChromaDB} for persistent SQLite-backed vector storage.
- \textbf{Relevance Score Normalization \& Anti-Hallucination
Guardrails}: Normalizes Euclidean distance and Cosine metrics into a
unified \(0.0 - 1.0\) relevance scale and enforces a strict relevance
threshold (\$ au = 0.20\$). If no document chunk exceeds the threshold,
the system immediately returns an explicit fallback refusal: \emph{``I
could not find enough information in the uploaded documents to answer
this question.''}---completely avoiding hallucination. -
\textbf{Multi-Model LLM Orchestration}: Supports zero-latency cloud
inference via \textbf{Groq Cloud} (Llama 3.1 8B Instant / GPT-OSS),
commercial enterprise inference via \textbf{OpenAI} (GPT-4o /
GPT-3.5-Turbo), and 100\% offline, private inference via local
\textbf{Ollama} (Llama 3, Mistral, Gemma). - \textbf{Asynchronous
Non-Blocking Workers}: Uses background daemon threads
(\texttt{threading.\allowbreak Thread} with Streamlit script execution context) so
students can trigger long-running quiz or flashcard generation jobs and
continue reading notes or asking questions without UI freezes. -
\textbf{Persistent Study Activity Tracking}: Records all generated study
artifacts to a structured \texttt{data/\allowbreak history.\allowbreak json} datastore with
unique UUIDs, timestamps, source filters, and one-click deletion
capabilities.

\subsection{1.5 Primary Pedagogical and Engineering
Objectives}\label{primary-pedagogical-and-engineering-objectives}

\begingroup\footnotesize
\begin{longtable}[]{@{}
  >{\raggedright\arraybackslash\sloppy}p{(\linewidth - 6\tabcolsep) * \real{0.1100}}
  >{\raggedright\arraybackslash\sloppy}p{(\linewidth - 6\tabcolsep) * \real{0.1500}}
  >{\raggedright\arraybackslash\sloppy}p{(\linewidth - 6\tabcolsep) * \real{0.4000}}
  >{\raggedright\arraybackslash\sloppy}p{(\linewidth - 6\tabcolsep) * \real{0.3000}}@{}}
\toprule\noalign{}
\begin{minipage}[b]{\linewidth}\raggedright
Objective ID
\end{minipage} & \begin{minipage}[b]{\linewidth}\raggedright
Category
\end{minipage} & \begin{minipage}[b]{\linewidth}\raggedright
Target Metric / Description
\end{minipage} & \begin{minipage}[b]{\linewidth}\raggedright
Implemented Verification
\end{minipage} \\
\midrule\noalign{}
\endhead
\bottomrule\noalign{}
\endlastfoot
\textbf{OBJ-01} & Grounded Q\&A & 100\% of answers backed by exact
source filename and 1-indexed page number & Verified via
\texttt{generate\_\allowbreak rag\_\allowbreak response(\allowbreak )\allowbreak } and \texttt{tests/\allowbreak test\_\allowbreak rag.\allowbreak py} \\
\textbf{OBJ-02} & Anti-Hallucination & Zero ungrounded LLM inference
when query relevance score \(<     au\) & Verified via
\texttt{test\_\allowbreak generate\_\allowbreak rag\_\allowbreak response\_\allowbreak no\_\allowbreak context} \\
\textbf{OBJ-03} & Pedagogical Utility & Automatic generation of 4
distinct study artifacts (Q\&A, Quizzes, Flashcards, Summaries) &
Verified in \texttt{src/\allowbreak generation/\allowbreak study\_\allowbreak generator.\allowbreak py} \\
\textbf{OBJ-04} & Inference Agnosticism & Seamless switching between
Groq, OpenAI, and local Ollama backends & Verified in
\texttt{tests/\allowbreak test\_\allowbreak llm\_\allowbreak manager.\allowbreak py} \\
\textbf{OBJ-05} & UI Responsiveness & Sub-second UI state updates with
background thread offloading for LLM generation & Implemented via
\texttt{worker\_\allowbreak generate\_\allowbreak quiz(\allowbreak )\allowbreak } in \texttt{app.\allowbreak py} \\
\textbf{OBJ-06} & Data Persistence & Complete session survival across
browser reloads via JSON and vector store persistence & Implemented in
\texttt{src/\allowbreak utils/\allowbreak history\_\allowbreak manager.\allowbreak py} and \texttt{src/\allowbreak vectorstore/\allowbreak } \\
\end{longtable}
\endgroup

\subsection{1.6 Target User Personas and
Stakeholders}\label{target-user-personas-and-stakeholders}

\begin{enumerate}
\def\labelenumi{\arabic{enumi}.}
\tightlist
\item
  \textbf{University Undergraduate \& Graduate Students}: Preparing for
  midterms and finals, needing rapid active-recall practice from 300+
  page textbook PDFs.
\item
  \textbf{Educators \& Teaching Assistants (TAs)}: Generating
  standardized multiple-choice practice quizzes and review flashcards
  directly aligned with course lecture slides.
\item
  \textbf{Self-Directed Learners \& Researchers}: Synthesizing complex
  whitepapers, technical specifications, and legal/regulatory documents
  into structured key takeaways.
\item
  \textbf{Academic Institutions \& Privacy-Sensitive Labs}: Requiring
  100\% local, air-gapped study assistance using Ollama and local
  sentence embeddings without data exfiltration.
\end{enumerate}

\subsection{1.7 Core Educational Use
Cases}\label{core-educational-use-cases}

\begin{verbatim}
+-----------------------------------------------------------------------------------+
|                            CORE STUDYRAG USE CASES                                |
+-----------------------------------------------------------------------------------+
|  1. Interactive Lecture Q&A     |  Students query complex terms and get answers   |
|                                 |  with exact page citations.                     |
|  2. Automated MCQ Quiz Testing  |  Generates 5-10 optioned questions with instant |
|                                 |  grading and rationale.                         |
|  3. Digital Flashcard Decks     |  Extracts definitions and creates front/back    |
|                                 |  cards for active recall.                       |
|  4. Chapter Revision Summaries  |  Produces structured bulleted key points and    |
|                                 |  revision guides on demand.                     |
|  5. Audit & History Management  |  Archives past study sessions with full filter  |
|                                 |  and deletion control.                          |
+-----------------------------------------------------------------------------------+
\end{verbatim}

\subsection{1.8 System Scope: In-Scope Capabilities vs.~Explicit
Boundaries}\label{system-scope-in-scope-capabilities-vs.-explicit-boundaries}

\subsubsection{In-Scope Features (Implemented \&
Verified)}\label{in-scope-features-implemented-verified}

\begin{itemize}
\tightlist
\item
  Multi-PDF document upload via Streamlit with SHA-256 deduplication.
\item
  PDF text extraction via PyMuPDF (\texttt{fitz}) with page-number
  tracking.
\item
  Whitespace normalization and newline-collapsing document cleaner.
\item
  Recursive character chunking with configurable overlap.
\item
  Sentence-Transformers embedding generation (\texttt{all-MiniLM-L6-v2}
  and \texttt{BAAI/\allowbreak bge-small-en-v1.5}).
\item
  FAISS vector store with local index serialization
  (\texttt{index.\allowbreak faiss}, \texttt{index.\allowbreak pkl}).
\item
  ChromaDB vector store with SQLite persistence and Cosine distance
  metric.
\item
  Dynamic query validation and length-bounding.
\item
  Dual distance metric normalization (L2 Euclidean and Cosine to
  \([0, 1]\) relevance).
\item
  Score thresholding and source-document filtering.
\item
  Prompt templates for RAG Q\&A, MCQs, Flashcards, and Chapter
  Summaries.
\item
  Multi-provider LLM integrations: Groq Cloud, OpenAI, Ollama.
\item
  Streamlit interactive UI with custom CSS, Dark/Light mode, and
  progress cards.
\item
  Multi-threaded background job workers for generation.
\item
  Persistent JSON history audit log with CRUD operations.
\item
  Retrieval evaluation metrics (Precision@K, Recall@K, Hit Rate@K).
\item
  28 automated pytest unit and integration tests.
\end{itemize}

\subsubsection{Explicitly Out-of-Scope (Boundaries for Version
1.0.0)}\label{explicitly-out-of-scope-boundaries-for-version-1.0.0}

\begin{itemize}
\tightlist
\item
  Optical Character Recognition (OCR) for scanned image-only PDFs
  (current engine flags non-extractable text).
\item
  Multi-modal image/diagram extraction and reasoning inside PDFs.
\item
  Distributed worker queues (e.g., Celery / Redis); current concurrency
  relies on Python in-process daemon threads.
\item
  Multi-tenant user authentication, RBAC, and cloud database
  multi-tenancy.
\item
  Export to external spaced-repetition formats (e.g., \texttt{.\allowbreak apkg}
  Anki decks or Quizlet sync).
\end{itemize}

\subsection{1.9 Practical Applications and Expected Educational
Outcomes}\label{practical-applications-and-expected-educational-outcomes}

When deployed in an academic setting, StudyRAG reduces study preparation
time by over 70\%. Students no longer spend hours typing flashcards or
copying textbook definitions into notes; instead, they immediately begin
active testing against their syllabus. By enforcing strict source
attribution, students can verify claims instantly, developing critical
analytical skills while preventing misinformation.

\subsection{1.10 Current Project Maturity and Production
Readiness}\label{current-project-maturity-and-production-readiness}

\begin{itemize}
\tightlist
\item
  \textbf{Current Stage:} Version 1.0.0 (Production-Ready Desktop \&
  Local Lab Deployment).
\item
  \textbf{Automated Test Coverage:} 100\% of core algorithmic modules
  covered by 28 automated tests.
\item
  \textbf{Code Cleanliness:} All deprecated imports isolated,
  environment variables centralized in \texttt{.\allowbreak env}, and structured
  logging implemented across all subsystems via
  \texttt{src.\allowbreak utils.\allowbreak logger}. \#\# Chapter 2: Requirements and
  Specifications
\end{itemize}

\subsection{2.1 Functional Requirements (FR-01 to
FR-12)}\label{functional-requirements-fr-01-to-fr-12}

The functional requirements govern the behavioral contracts of the
StudyRAG application:

\begin{itemize}
\tightlist
\item
  \textbf{FR-01: Multi-Document PDF Ingestion}\\
  The system MUST accept multiple PDF uploads simultaneously, compute a
  SHA-256 hash for deduplication, extract text page-by-page while
  tracking 1-indexed page metadata, and store uploaded files in
  \texttt{data/\allowbreak uploaded/\allowbreak }.
\item
  \textbf{FR-02: Text Cleaning \& Whitespace Normalization}\\
  The system MUST normalize consecutive spaces/tabs to single spaces,
  collapse three or more consecutive newlines into two newlines
  (preserving paragraph boundaries), and drop completely blank pages.
\item
  \textbf{FR-03: Recursive Contextual Chunking}\\
  The system MUST segment cleaned text using LangChain's
  \texttt{RecursiveCharacterTextSplitter} with separators
  \texttt{{[}"\textbackslash{}n\textbackslash{}n",\allowbreak \ \allowbreak "\textbackslash{}n",\allowbreak \ \allowbreak ".\ \allowbreak ",\allowbreak \ \allowbreak "\ \allowbreak ",\allowbreak \ \allowbreak ""{]}},
  configurable chunk size (default 500 chars), and overlap (default 50
  chars), generating unique incremental chunk IDs.
\item
  \textbf{FR-04: Dense Vector Embedding Generation}\\
  The system MUST generate dense vector embeddings for all chunks using
  HuggingFace / SentenceTransformers models (\texttt{all-MiniLM-L6-v2}
  producing 384-dimensional vectors).
\item
  \textbf{FR-05: Vector Store Persistence \& Retrieval}\\
  The system MUST index chunk embeddings into either FAISS or ChromaDB,
  persist the index to disk (\texttt{data/\allowbreak vectorstores/\allowbreak }), and support
  top-\(K\) similarity search.
\item
  \textbf{FR-06: Query Processing \& Validation}\\
  The system MUST validate user queries, reject empty/whitespace-only
  input, collapse redundant spaces, and enforce a maximum character
  length (1000 characters).
\item
  \textbf{FR-07: Relevance Normalization \& Threshold Filtering}\\
  The system MUST normalize vendor-specific distance metrics into a
  \([0.0, 1.0]\) relevance scale and filter out all chunks failing a
  minimum threshold \$ au\$ (default 0.20).
\item
  \textbf{FR-08: Source-Specific Filtering}\\
  The system MUST allow users to restrict similarity searches to a
  specific document or search across ``All Materials''.
\item
  \textbf{FR-09: Anti-Hallucination Fallback Guardrail}\\
  If no retrieved chunk satisfies the threshold, the system MUST return
  a deterministic refusal message without dispatching a prompt to the
  LLM.
\item
  \textbf{FR-10: Multi-Format Study Artifact Generation}\\
  The system MUST generate 4 distinct study artifacts: Grounded Q\&A,
  Multiple Choice Quizzes (in valid JSON), Digital Flashcards
  (front/back JSON), and Chapter Revision Summaries (Markdown).
\item
  \textbf{FR-11: Asynchronous Background Generation}\\
  The system MUST execute quiz, flashcard, and summary generation inside
  background daemon threads (\texttt{threading.\allowbreak Thread}), allowing
  non-blocking UI navigation.
\item
  \textbf{FR-12: Persistent History Management}\\
  The system MUST persist all generated study artifacts to
  \texttt{data/\allowbreak history.\allowbreak json} with UUIDs, timestamps, source tags, and
  support filtering and deletion.
\end{itemize}

\subsection{2.2 Non-Functional Requirements (NFR-01 to
NFR-08)}\label{non-functional-requirements-nfr-01-to-nfr-08}

\begin{itemize}
\tightlist
\item
  \textbf{NFR-01: Latency \& Responsiveness}

  \begin{itemize}
  \tightlist
  \item
    Semantic vector search across 1,000 chunks MUST complete within
    \(< 100\,\text{ms}\) on standard CPU.
  \item
    LLM first-token generation using Groq Cloud MUST start in
    \(< 1.5\,\text{s}\).
  \end{itemize}
\item
  \textbf{NFR-02: Deterministic Citations}\\
  All answers MUST cite the exact filename and page number from which
  context was drawn.
\item
  \textbf{NFR-03: Memory Footprint}\\
  The base application running in Python 3.12 with FAISS and
  \texttt{all-MiniLM-L6-v2} MUST operate within \(< 2.0\,\text{GB}\) of
  host RAM.
\item
  \textbf{NFR-04: Cross-Platform Compatibility}\\
  The codebase MUST execute cleanly on macOS (Darwin arm64/x86\_64),
  Linux (x86\_64), and Windows 11 without platform-specific
  modifications.
\item
  \textbf{NFR-05: Privacy \& Air-Gap Capability}\\
  The system MUST support complete offline execution using local
  SentenceTransformers embeddings and local Ollama inference without
  outbound internet calls.
\item
  \textbf{NFR-06: Fault Tolerance \& Graceful Degradation}\\
  Corrupted or encrypted PDFs MUST trigger caught exceptions and
  user-facing error notices without crashing the Streamlit application
  server.
\item
  \textbf{NFR-07: Maintainability \& Modularity}\\
  The codebase MUST maintain strict separation of concerns across
  \texttt{ingestion}, \texttt{embeddings}, \texttt{vectorstore},
  \texttt{retrieval}, \texttt{generation}, and \texttt{utils}.
\item
  \textbf{NFR-08: Test Coverage \& Regression Safety}\\
  All critical algorithms and pipeline components MUST maintain
  automated unit test coverage executable via \texttt{pytest}.
\end{itemize}

\subsection{2.3 Hardware, Runtime, and Platform
Specifications}\label{hardware-runtime-and-platform-specifications}

\begingroup\footnotesize
\begin{longtable}[]{@{}
  >{\raggedright\arraybackslash\sloppy}p{(\linewidth - 4\tabcolsep) * \real{0.2200}}
  >{\raggedright\arraybackslash\sloppy}p{(\linewidth - 4\tabcolsep) * \real{0.3700}}
  >{\raggedright\arraybackslash\sloppy}p{(\linewidth - 4\tabcolsep) * \real{0.3700}}@{}}
\toprule\noalign{}
\begin{minipage}[b]{\linewidth}\raggedright
Component
\end{minipage} & \begin{minipage}[b]{\linewidth}\raggedright
Minimum Specification
\end{minipage} & \begin{minipage}[b]{\linewidth}\raggedright
Recommended Specification
\end{minipage} \\
\midrule\noalign{}
\endhead
\bottomrule\noalign{}
\endlastfoot
\textbf{Processor (CPU)} & Intel Core i5 / AMD Ryzen 5 / Apple M1 (4
Cores) & Apple M1/M2/M3 or 8-Core Intel/AMD CPU \\
\textbf{System Memory (RAM)} & 8 GB DDR4 / Unified Memory & 16 GB+
Unified Memory (essential if running local Ollama 8B models) \\
\textbf{Storage (Disk)} & 2 GB free disk space (for libraries \& models)
& 10 GB+ free SSD storage (for local LLM weights) \\
\textbf{Python Runtime} & Python 3.10.x & Python 3.12.x (Verified on
3.12.2) \\
\textbf{Operating System} & macOS Monterey, Ubuntu 20.04, Windows 10 &
macOS Sonoma / Sequoia, Ubuntu 22.04 LTS, Windows 11 \\
\end{longtable}
\endgroup

\subsection{2.4 Requirements Traceability
Matrix}\label{requirements-traceability-matrix}

\begingroup\scriptsize
\begin{longtable}[]{@{}
  >{\raggedright\arraybackslash\sloppy}p{(\linewidth - 8\tabcolsep) * \real{0.0900}}
  >{\raggedright\arraybackslash\sloppy}p{(\linewidth - 8\tabcolsep) * \real{0.2600}}
  >{\raggedright\arraybackslash\sloppy}p{(\linewidth - 8\tabcolsep) * \real{0.2400}}
  >{\raggedright\arraybackslash\sloppy}p{(\linewidth - 8\tabcolsep) * \real{0.2400}}
  >{\raggedright\arraybackslash\sloppy}p{(\linewidth - 8\tabcolsep) * \real{0.1300}}@{}}
\toprule\noalign{}
\begin{minipage}[b]{\linewidth}\raggedright
Req ID
\end{minipage} & \begin{minipage}[b]{\linewidth}\raggedright
Requirement Description
\end{minipage} & \begin{minipage}[b]{\linewidth}\raggedright
Implementing File(s)
\end{minipage} & \begin{minipage}[b]{\linewidth}\raggedright
Primary Function / Class
\end{minipage} & \begin{minipage}[b]{\linewidth}\raggedright
Verification Status
\end{minipage} \\
\midrule\noalign{}
\endhead
\bottomrule\noalign{}
\endlastfoot
\textbf{FR-01} & PDF Ingestion \& Hashing &
\texttt{src/\allowbreak ingestion/\allowbreak pdf\_\allowbreak loader.\allowbreak py}
\newline \texttt{src/\allowbreak utils/\allowbreak file\_\allowbreak utils.\allowbreak py} & \texttt{load\_\allowbreak pdf(\allowbreak )\allowbreak }
\newline \texttt{calculate\_\allowbreak file\_\allowbreak hash(\allowbreak )\allowbreak } & PASSED
(\texttt{tests/\allowbreak test\_\allowbreak pdf\_\allowbreak loader.\allowbreak py}) \\
\textbf{FR-02} & Text Cleaning &
\texttt{src/\allowbreak ingestion/\allowbreak document\_\allowbreak cleaner.\allowbreak py} & \texttt{clean\_\allowbreak text(\allowbreak )\allowbreak },
\texttt{clean\_\allowbreak documents(\allowbreak )\allowbreak } & PASSED
(\texttt{tests/\allowbreak test\_\allowbreak document\_\allowbreak cleaner.\allowbreak py}) \\
\textbf{FR-03} & Recursive Chunking & \texttt{src/\allowbreak ingestion/\allowbreak chunker.\allowbreak py}
& \texttt{create\_\allowbreak chunks(\allowbreak )\allowbreak } & PASSED
(\texttt{tests/\allowbreak test\_\allowbreak chunking.\allowbreak py}) \\
\textbf{FR-04} & Dense Embeddings &
\texttt{src/\allowbreak embeddings/\allowbreak embedding\_\allowbreak manager.\allowbreak py} &
\texttt{EmbeddingManager} & PASSED
(\texttt{tests/\allowbreak test\_\allowbreak embeddings.\allowbreak py}) \\
\textbf{FR-05} & Vector DB Persistence &
\texttt{src/\allowbreak vectorstore/\allowbreak vector\_\allowbreak store\_\allowbreak manager.\allowbreak py}
\newline \texttt{src/\allowbreak vectorstore/\allowbreak faiss\_\allowbreak store.\allowbreak py} & \texttt{FAISSStore},
\texttt{ChromaStore} & PASSED (\texttt{tests/\allowbreak test\_\allowbreak vectorstore.\allowbreak py}) \\
\textbf{FR-06} & Query Cleaning &
\texttt{src/\allowbreak retrieval/\allowbreak query\_\allowbreak processor.\allowbreak py} & \texttt{process\_\allowbreak query(\allowbreak )\allowbreak } &
PASSED (\texttt{tests/\allowbreak test\_\allowbreak query\_\allowbreak processor.\allowbreak py}) \\
\textbf{FR-07} & Score Normalization &
\texttt{src/\allowbreak retrieval/\allowbreak score\_\allowbreak filter.\allowbreak py} & \texttt{normalize\_\allowbreak score(\allowbreak )\allowbreak },
\texttt{filter\_\allowbreak by\_\allowbreak threshold(\allowbreak )\allowbreak } & PASSED
(\texttt{tests/\allowbreak test\_\allowbreak retrieval.\allowbreak py}) \\
\textbf{FR-08} & Document Filtering &
\texttt{src/\allowbreak retrieval/\allowbreak retriever.\allowbreak py} &
\texttt{Retriever.\allowbreak retrieve\_\allowbreak with\_\allowbreak scores(\allowbreak )\allowbreak } & PASSED
(\texttt{tests/\allowbreak test\_\allowbreak retrieval.\allowbreak py}) \\
\textbf{FR-09} & Anti-Hallucination & \texttt{src/\allowbreak rag/\allowbreak rag\_\allowbreak pipeline.\allowbreak py}
& \texttt{generate\_\allowbreak rag\_\allowbreak response(\allowbreak )\allowbreak } & PASSED
(\texttt{tests/\allowbreak test\_\allowbreak rag.\allowbreak py}) \\
\textbf{FR-10} & Study Artifact Gen &
\texttt{src/\allowbreak generation/\allowbreak study\_\allowbreak generator.\allowbreak py} & \texttt{generate\_\allowbreak quiz(\allowbreak )\allowbreak },
\texttt{generate\_\allowbreak flashcards(\allowbreak )\allowbreak } & PASSED
(\texttt{tests/\allowbreak test\_\allowbreak llm\_\allowbreak manager.\allowbreak py}) \\
\textbf{FR-11} & Async Background Threading & \texttt{app.\allowbreak py} &
\texttt{worker\_\allowbreak generate\_\allowbreak quiz(\allowbreak )\allowbreak } & PASSED (Manual \& Runtime
Verified) \\
\textbf{FR-12} & Study History Persistence &
\texttt{src/\allowbreak utils/\allowbreak history\_\allowbreak manager.\allowbreak py} & \texttt{load\_\allowbreak history(\allowbreak )\allowbreak },
\texttt{add\_\allowbreak history\_\allowbreak entry(\allowbreak )\allowbreak } & PASSED (Verified via
\texttt{data/\allowbreak history.\allowbreak json}) \\
\end{longtable}
\endgroup

\subsection{2.5 Assumptions, Dependencies, and Operational
Constraints}\label{assumptions-dependencies-and-operational-constraints}

\begin{enumerate}
\def\labelenumi{\arabic{enumi}.}
\tightlist
\item
  \textbf{PDF Extractability}: The system assumes documents contain text
  streams extractable by PyMuPDF (\texttt{fitz}). Image-only scanned
  PDFs require external pre-processing with OCR software.
\item
  \textbf{Apple Silicon Environment Guard}: In Anaconda environments on
  macOS where TensorFlow is present, the environment variable
  \texttt{USE\_\allowbreak TF=\allowbreak 0} MUST be set before importing HuggingFace
  Transformers to avoid Apple Silicon C-extension segmentation faults.
\item
  \textbf{API Keys for Cloud Inference}: Utilizing Groq Cloud or OpenAI
  requires valid API keys in \texttt{.\allowbreak env}. When running offline, a
  local Ollama daemon must be accessible at
  \texttt{http:/\allowbreak /\allowbreak localhost:11434}.
\end{enumerate}

\section{Chapter 3: Complete Technology
Stack}\label{chapter-3-complete-technology-stack}

\subsection{3.1 Technology Stack Architecture
Matrix}\label{technology-stack-architecture-matrix}

Every library and dependency in StudyRAG has been verified against
\texttt{requirements.\allowbreak txt} and the active virtual environment:

\begingroup\scriptsize
\begin{longtable}[]{@{}
  >{\raggedright\arraybackslash\sloppy}p{(\linewidth - 12\tabcolsep) * \real{0.1300}}
  >{\raggedright\arraybackslash\sloppy}p{(\linewidth - 12\tabcolsep) * \real{0.1000}}
  >{\raggedright\arraybackslash\sloppy}p{(\linewidth - 12\tabcolsep) * \real{0.1100}}
  >{\raggedright\arraybackslash\sloppy}p{(\linewidth - 12\tabcolsep) * \real{0.1700}}
  >{\raggedright\arraybackslash\sloppy}p{(\linewidth - 12\tabcolsep) * \real{0.1800}}
  >{\raggedright\arraybackslash\sloppy}p{(\linewidth - 12\tabcolsep) * \real{0.1600}}
  >{\raggedright\arraybackslash\sloppy}p{(\linewidth - 12\tabcolsep) * \real{0.1100}}@{}}
\toprule\noalign{}
\begin{minipage}[b]{\linewidth}\raggedright
Technology
\end{minipage} & \begin{minipage}[b]{\linewidth}\raggedright
Verified Version
\end{minipage} & \begin{minipage}[b]{\linewidth}\raggedright
Category
\end{minipage} & \begin{minipage}[b]{\linewidth}\raggedright
Role in StudyRAG
\end{minipage} & \begin{minipage}[b]{\linewidth}\raggedright
Why Needed
\end{minipage} & \begin{minipage}[b]{\linewidth}\raggedright
Source Location
\end{minipage} & \begin{minipage}[b]{\linewidth}\raggedright
Alternatives Evaluated
\end{minipage} \\
\midrule\noalign{}
\endhead
\bottomrule\noalign{}
\endlastfoot
\textbf{Python} & 3.12.2 & Language & Core runtime execution & Native
support for modern AI/ML packages and type hinting & Entire repository &
Python 3.10, 3.11 \\
\textbf{Streamlit} & \(\ge 1.32.0\) & UI Framework & Interactive web
application & Instant reactive state, native chat widgets,
zero-boilerplate web hosting & \texttt{app.\allowbreak py} & Gradio, FastAPI +
React \\
\textbf{PyMuPDF (\texttt{fitz})} & \(\ge 1.23.0\) (1.28.2) & Document
Processing & High-speed PDF text extraction & 10x-20x faster than
pure-Python extractors with precise page tracking &
\texttt{src/\allowbreak ingestion/\allowbreak pdf\_\allowbreak loader.\allowbreak py} & pypdf, pdfplumber, PyPDF2 \\
\textbf{LangChain Core} & \(\ge 0.1.13\) & Orchestration & Prompt
management and LCEL pipelines & Standardized chaining
(\texttt{prompt\ \allowbreak \textbar{}\ \allowbreak llm\ \allowbreak \textbar{}\ \allowbreak parser}) and Document
interfaces & \texttt{src/\allowbreak rag/\allowbreak rag\_\allowbreak pipeline.\allowbreak py}
\newline \texttt{src/\allowbreak generation/\allowbreak } & LlamaIndex, Haystack \\
\textbf{langchain-text-splitters} & \(\ge 0.0.1\) & Preprocessing &
Recursive character splitting & Multi-stage boundary fallback preserving
sentence semantics & \texttt{src/\allowbreak ingestion/\allowbreak chunker.\allowbreak py} & Fixed-length
sliding window \\
\textbf{langchain-huggingface} & \(\ge 0.0.1\) & Embeddings &
HuggingFace embedding integration & Unified LangChain embedding
interface for local transformers &
\texttt{src/\allowbreak embeddings/\allowbreak embedding\_\allowbreak manager.\allowbreak py} & Direct PyTorch
embeddings \\
\textbf{sentence-transformers} & \(\ge 2.5.1\) & NLP Embeddings &
Sentence-level semantic vector generation & Produces high-quality dense
384d semantic vectors locally & \texttt{src/\allowbreak embeddings/\allowbreak } & OpenAI
text-embedding-ada-002 \\
\textbf{FAISS CPU} & \(\ge 1.8.0\) & Vector DB & In-memory dense vector
indexing & Sub-millisecond vector similarity search with zero database
overhead & \texttt{src/\allowbreak vectorstore/\allowbreak faiss\_\allowbreak store.\allowbreak py} & Annoy, ScaNN \\
\textbf{ChromaDB} & \(\ge 0.4.24\) & Vector DB & Persistent embedded
vector database & Native SQLite persistence and cosine collection
configuration & \texttt{src/\allowbreak vectorstore/\allowbreak chroma\_\allowbreak store.\allowbreak py} & Milvus,
Qdrant, Pinecone \\
\textbf{langchain-openai} & \(\ge 0.1.0\) & LLM Integration & Interface
to OpenAI and Groq APIs & Reuses OpenAI chat completion standard for
both OpenAI and Groq & \texttt{src/\allowbreak generation/\allowbreak llm\_\allowbreak manager.\allowbreak py} & Raw
\texttt{requests} / HTTP client \\
\textbf{langchain-ollama} & \(\ge 0.1.0\) & LLM Integration & Local
offline LLM connectivity & Connects directly to local Ollama daemon for
private inference & \texttt{src/\allowbreak generation/\allowbreak llm\_\allowbreak manager.\allowbreak py} & Ollama
Python SDK \\
\textbf{PyYAML} & \(\ge 6.0.1\) & Configuration & Configuration file
parsing & Clean YAML syntax for application hyperparameters &
\texttt{app.\allowbreak py} \newline \texttt{config/\allowbreak settings.\allowbreak yaml} & JSON, TOML \\
\textbf{python-dotenv} & \(\ge 1.0.1\) & Security & Environment secrets
management & Ingests \texttt{.\allowbreak env} without hardcoding sensitive API keys
& \texttt{app.\allowbreak py} \newline \texttt{src/\allowbreak generation/\allowbreak llm\_\allowbreak manager.\allowbreak py} & OS
export scripts \\
\textbf{Pandas \& NumPy} & \(\ge 2.2.1\) / \(\ge 1.26.0\) & Data
Processing & Vector math and table structures & Array manipulation and
score sorting & \texttt{src/\allowbreak retrieval/\allowbreak }
\newline \texttt{src/\allowbreak evaluation/\allowbreak } & Pure Python math \\
\textbf{Pytest} & \(\ge 8.0.0\) (8.1.1) & Testing & Automated test suite
execution & Industrial test runner with fixture and mocking support &
\texttt{tests/\allowbreak } & Unittest \\
\end{longtable}
\endgroup

\subsection{3.2 Runtime and Programming
Language}\label{runtime-and-programming-language}

\begin{itemize}
\tightlist
\item
  \textbf{Python 3.12}: Selected for enhanced execution speed, advanced
  pattern matching, improved error tracebacks, and robust compatibility
  with modern PyTorch wheels on macOS ARM64.
\item
  \textbf{Type Annotations}: Implemented across all \texttt{src/\allowbreak }
  modules (\texttt{List}, \texttt{Dict}, \texttt{Any},
  \texttt{Optional}, \texttt{Union}, \texttt{Path}) ensuring static
  contract safety and IDE discoverability.
\end{itemize}

\subsection{3.3 Frontend and Presentation
Framework}\label{frontend-and-presentation-framework}

\begin{itemize}
\tightlist
\item
  \textbf{Streamlit}: StudyRAG leverages Streamlit's reactive execution
  engine. Key design elements:

  \begin{itemize}
  \tightlist
  \item
    Custom responsive CSS injected via \texttt{apply\_\allowbreak theme\_\allowbreak css(\allowbreak )\allowbreak }.
  \item
    Modern dual theme system: {[}Dark Mode{]} Dark Mode
    (\texttt{\#0F172A} Slate palette) and {[}Light Mode{]} Light Mode
    (\texttt{\#F8FAFC} clean palette).
  \item
    Web components: \texttt{st.\allowbreak chat\_\allowbreak message}, \texttt{st.\allowbreak chat\_\allowbreak input},
    \texttt{st.\allowbreak expander}, \texttt{st.\allowbreak radio}, and custom gradient
    buttons.
  \item
    Reactive cache: \texttt{@st.\allowbreak cache\_\allowbreak resource} prevents
    re-instantiating heavy SentenceTransformer models upon UI
    re-renders.
  \end{itemize}
\end{itemize}

\subsection{3.4 Document Processing and Text
Extraction}\label{document-processing-and-text-extraction}

\begin{itemize}
\tightlist
\item
  \textbf{PyMuPDF (\texttt{fitz})}: Developed in high-performance C++
  with Python bindings. PyMuPDF extracts text streams from complex
  academic PDFs up to 20x faster than legacy pure-Python packages like
  \texttt{pypdf}. It accurately captures page boundaries, enabling exact
  page-level attribution.
\end{itemize}

\subsection{3.5 Orchestration and LLM
Frameworks}\label{orchestration-and-llm-frameworks}

\begin{itemize}
\item
  \textbf{LangChain Expression Language (LCEL)}: Replaces cumbersome
  legacy \texttt{LLMChain} objects with declarative unix-pipe syntax:

\begin{Shaded}
\begin{Highlighting}[]
\NormalTok{chain }\OperatorTok{=}\NormalTok{ prompt }\OperatorTok{|}\NormalTok{ llm }\OperatorTok{|}\NormalTok{ StrOutputParser()}
\end{Highlighting}
\end{Shaded}

  This reduces memory overhead, simplifies streaming, and provides
  transparent tracing.
\end{itemize}

\subsection{3.6 Embedding Models and Semantic Vector
Spaces}\label{embedding-models-and-semantic-vector-spaces}

StudyRAG features an embedding manager supporting two leading
open-weights models: 1. \texttt{sentence-transformers/\allowbreak all-MiniLM-L6-v2}
(Default): - Architecture: 6-layer MiniLM transformer. - Vector
Dimension: 384 dimensions. - Sequence Length: 256 word pieces. -
Performance: Blazing-fast inference on CPU (sub-50ms for typical
paragraphs) with high semantic clustering fidelity. 2.
\texttt{BAAI/\allowbreak bge-small-en-v1.5} (Selectable): - Architecture: 12-layer
small BERT architecture. - Vector Dimension: 384 dimensions. -
Top-ranked on the MTEB (Massive Text Embedding Benchmark) for retrieval
tasks.

\subsection{3.7 Vector Database Backends: FAISS and
ChromaDB}\label{vector-database-backends-faiss-and-chromadb}

\begin{itemize}
\tightlist
\item
  \textbf{FAISS (Facebook AI Similarity Search)}:

  \begin{itemize}
  \tightlist
  \item
    Implementation: \texttt{IndexFlatL2} dense index.
  \item
    Characteristics: Runs entirely in RAM; provides instantaneous
    nearest-neighbor search. Serialized to disk using
    \texttt{save\_\allowbreak local(\allowbreak )\allowbreak } and loaded with
    \texttt{allow\_\allowbreak dangerous\_\allowbreak deserialization=\allowbreak True}.
  \end{itemize}
\item
  \textbf{ChromaDB}:

  \begin{itemize}
  \tightlist
  \item
    Implementation: Embedded SQLite database with HNSW vector index
    files.
  \item
    Configuration: Configured explicitly with
    \texttt{\{"hnsw:space":\ \allowbreak "cosine"\}} to guarantee true cosine
    distance metrics.
  \end{itemize}
\end{itemize}

\subsection{3.8 Large Language Model (LLM) Inference
Providers}\label{large-language-model-llm-inference-providers}

\begin{enumerate}
\def\labelenumi{\arabic{enumi}.}
\tightlist
\item
  \textbf{Groq Cloud Inference}:

  \begin{itemize}
  \tightlist
  \item
    Provider Type: Cloud LPU (Language Processing Unit) hosting.
  \item
    Protocol: OpenAI-compatible REST endpoint
    (\texttt{https:/\allowbreak /\allowbreak api.\allowbreak groq.\allowbreak com/\allowbreak openai/\allowbreak v1}).
  \item
    Models: Llama 3.1 8B Instant (\texttt{openai/\allowbreak gpt-oss-20b}), 70B
    (\texttt{openai/\allowbreak gpt-oss-120b}).
  \item
    Latency: Phenomenal inference speed (500+ tokens/second).
  \end{itemize}
\item
  \textbf{OpenAI API}:

  \begin{itemize}
  \tightlist
  \item
    Protocol: Native REST API.
  \item
    Models: \texttt{gpt-4o}, \texttt{gpt-3.5-turbo}.
  \item
    Use Case: Highly complex reasoning and high-stakes exam synthesis.
  \end{itemize}
\item
  \textbf{Local Ollama Daemon}:

  \begin{itemize}
  \tightlist
  \item
    Protocol: Local HTTP REST server (\texttt{http:/\allowbreak /\allowbreak localhost:11434}).
  \item
    Models: \texttt{llama3}, \texttt{mistral}, \texttt{gemma}.
  \item
    Use Case: 100\% offline, privacy-preserving study with zero token
    costs.
  \end{itemize}
\end{enumerate}

\subsection{3.9 Data Serialization, Testing, and Utility
Tooling}\label{data-serialization-testing-and-utility-tooling}

\begin{itemize}
\tightlist
\item
  \textbf{History Data}: Serialized as human-readable JSON
  (\texttt{data/\allowbreak history.\allowbreak json}) with atomic write patterns.
\item
  \textbf{Pytest}: Enterprise testing framework supporting parameterized
  test matrices and mock patch execution (\texttt{unittest.\allowbreak mock}). \#\#
  Chapter 4: Project Architecture
\end{itemize}

\subsection{4.1 Architectural Style and High-Level Design
Principles}\label{architectural-style-and-high-level-design-principles}

StudyRAG follows a \textbf{Layered Modular Architecture} governed by
clean software engineering principles: 1. \textbf{Separation of Concerns
(SoC)}: Ingestion, embedding, vector indexing, retrieval, generation,
and presentation are strictly decoupled. No UI logic leaks into core
pipelines. 2. \textbf{Factory \& Strategy Patterns}: Vector store
providers (\texttt{FAISSStore}, \texttt{ChromaStore}) and LLM providers
(\texttt{Groq}, \texttt{OpenAI}, \texttt{Ollama}) implement unified
interfaces managed by centralized factories
(\texttt{VectorStoreManager}, \texttt{LLMManager}). 3.
\textbf{Immutability \& Statelessness in Core Pipelines}: Pipeline
functions (\texttt{process\_\allowbreak pdf}, \texttt{retrieve\_\allowbreak context},
\texttt{generate\_\allowbreak rag\_\allowbreak response}) are stateless functions that accept
configuration and data as inputs, returning structured dictionaries. 4.
\textbf{Defensive Anti-Hallucination Design}: The generation layer
refuses to invoke costly LLMs when retrieved context fails quality and
relevance thresholds.

\subsection{4.2 Layered System Topology
Diagram}\label{layered-system-topology-diagram}

\begin{Shaded}
\begin{Highlighting}[]
\NormalTok{graph TD}
\NormalTok{    subgraph UI ["Layer 1: Presentation Layer (Streamlit)"]}
\NormalTok{        A[Web Browser / User] {-}{-}\textgreater{} B[app.py Main UI Controller]}
\NormalTok{        B {-}{-}\textgreater{} C1[[Study Library]  Study Library]}
\NormalTok{        B {-}{-}\textgreater{} C2[[Ask Questions]  Ask Questions]}
\NormalTok{        B {-}{-}\textgreater{} C3[[Practice Quiz]  Practice Quiz]}
\NormalTok{        B {-}{-}\textgreater{} C4[[Flashcards]  Flashcards]}
\NormalTok{        B {-}{-}\textgreater{} C5[[Topic Summaries]  Summaries]}
\NormalTok{        B {-}{-}\textgreater{} C6[[Study History]  Study History]}
\NormalTok{    end}

\NormalTok{    subgraph WORKERS ["Layer 2: Concurrency \& Async Worker Subsystem"]}
\NormalTok{        B {-}{-}\textgreater{} W1[threading.Thread / add\_script\_run\_ctx]}
\NormalTok{        W1 {-}{-}\textgreater{} W2[GLOBAL\_BACKGROUND\_JOBS Registry]}
\NormalTok{    end}

\NormalTok{    subgraph PIPELINES ["Layer 3: Core Orchestration Pipelines"]}
\NormalTok{        C1 {-}{-}\textgreater{} P1[Ingestion Pipeline \textbackslash{}newline process\_pdf]}
\NormalTok{        C2 {-}{-}\textgreater{} P2[RAG Pipeline \textbackslash{}newline generate\_rag\_response]}
\NormalTok{        C3 \& C4 \& C5 {-}{-}\textgreater{} P3[Study Generator \textbackslash{}newline generate\_quiz / flashcards / summary]}
\NormalTok{        P2 \& P3 {-}{-}\textgreater{} P4[Retrieval Pipeline \textbackslash{}newline retrieve\_context]}
\NormalTok{    end}

\NormalTok{    subgraph CORE ["Layer 4: Core Engine Subsystems"]}
\NormalTok{        P1 {-}{-}\textgreater{} E1[PDF Loader PyMuPDF]}
\NormalTok{        P1 {-}{-}\textgreater{} E2[Document Cleaner Regex]}
\NormalTok{        P1 {-}{-}\textgreater{} E3[Recursive Chunker LangChain]}
        
\NormalTok{        P4 {-}{-}\textgreater{} R1[Query Processor]}
\NormalTok{        P4 {-}{-}\textgreater{} R2[Retriever Engine]}
\NormalTok{        P4 {-}{-}\textgreater{} R3[Score Normalizer \& Filter]}
        
\NormalTok{        P2 \& P3 {-}{-}\textgreater{} G1[Prompt Templates LCEL]}
\NormalTok{        P2 \& P3 {-}{-}\textgreater{} G2[LLM Manager]}
\NormalTok{    end}

\NormalTok{    subgraph STORAGE ["Layer 5: Persistent Storage \& External Services"]}
\NormalTok{        E3 {-}{-}\textgreater{} VSM[VectorStoreManager]}
\NormalTok{        VSM {-}{-}\textgreater{} V1[(FAISS Store \textbackslash{}newline index.faiss)]}
\NormalTok{        VSM {-}{-}\textgreater{} V2[(ChromaDB Store \textbackslash{}newline chroma.sqlite3)]}
        
\NormalTok{        P1 {-}{-}\textgreater{} F1[(Uploaded PDFs \textbackslash{}newline data/uploaded/)]}
\NormalTok{        W2 \& B {-}{-}\textgreater{} H1[(Study History \textbackslash{}newline data/history.json)]}
        
\NormalTok{        G2 {-}{-}\textgreater{} EXT1[Groq Cloud API]}
\NormalTok{        G2 {-}{-}\textgreater{} EXT2[OpenAI API]}
\NormalTok{        G2 {-}{-}\textgreater{} EXT3[Local Ollama Daemon]}
\NormalTok{    end}
\end{Highlighting}
\end{Shaded}

\subsection{4.3 Component Boundaries and Interface
Contracts}\label{component-boundaries-and-interface-contracts}

The system enforces strict data contracts across module boundaries:

\begin{verbatim}
[Raw PDF File]
      |
      v  (load_pdf)
List[Dict[str, Any]]  -->  {"text": str, "source": str, "page": int}
      |
      v  (clean_documents)
List[Dict[str, Any]]  -->  {"text": str (normalized), "source": str, "page": int}
      |
      v  (create_chunks)
List[Dict[str, Any]]  -->  {"text": str, "source": str, "page": int, "chunk_id": int}
      |
      v  (raw_chunks_to_documents)
List[Document]        -->  LangChain Document(page_content, metadata={source, page, chunk_id})
      |
      v  (similarity_search_with_score)
List[Tuple[Document, float]]  -->  (Document, raw_score)
      |
      v  (filter_by_threshold)
List[Dict[str, Any]]  -->  {"text": str, "source": str, "page": int, "chunk_id": int, "score": float}
\end{verbatim}

\subsection{4.4 Data Flow Architecture: From PDF Upload to Grounded
Response}\label{data-flow-architecture-from-pdf-upload-to-grounded-response}

\begin{Shaded}
\begin{Highlighting}[]
\NormalTok{sequenceDiagram}
\NormalTok{    autonumber}
\NormalTok{    actor Student as  Student / User}
\NormalTok{    participant UI as  app.py (Streamlit)}
\NormalTok{    participant Ingest as  Ingestion Pipeline}
\NormalTok{    participant Embed as [Embeddings Engine]  Embedding Manager}
\NormalTok{    participant VStore as  VectorStoreManager}
\NormalTok{    participant Retr as  Retrieval Pipeline}
\NormalTok{    participant LLM as  LLM Manager}

\NormalTok{    \%\% Ingestion Phase}
\NormalTok{    Note over Student, VStore: Phase 1: Ingestion \& Indexing}
\NormalTok{    Student{-}\textgreater{}\textgreater{}UI: Upload syllabus.pdf}
\NormalTok{    UI{-}\textgreater{}\textgreater{}Ingest: process\_pdf(syllabus.pdf, size=500, overlap=50)}
\NormalTok{    Ingest{-}\textgreater{}\textgreater{}Ingest: PyMuPDF extract text page{-}by{-}page}
\NormalTok{    Ingest{-}\textgreater{}\textgreater{}Ingest: Clean whitespace \& remove empty pages}
\NormalTok{    Ingest{-}\textgreater{}\textgreater{}Ingest: RecursiveCharacterTextSplitter (chunks)}
\NormalTok{    Ingest{-}{-}\textgreater{}\textgreater{}UI: Return List of chunk dicts}
\NormalTok{    UI{-}\textgreater{}\textgreater{}Embed: Get HuggingFace embeddings (all{-}MiniLM{-}L6{-}v2)}
\NormalTok{    UI{-}\textgreater{}\textgreater{}VStore: add\_documents(chunks) \& save()}
\NormalTok{    VStore{-}{-}\textgreater{}\textgreater{}UI: Index serialized to disk}

\NormalTok{    \%\% RAG Q\&A Phase}
\NormalTok{    Note over Student, LLM: Phase 2: Grounded Q\&A}
\NormalTok{    Student{-}\textgreater{}\textgreater{}UI: Submit question: "What is an OS thread?"}
\NormalTok{    UI{-}\textgreater{}\textgreater{}Retr: retrieve\_context(query, top\_k=5, threshold=0.20)}
\NormalTok{    Retr{-}\textgreater{}\textgreater{}Retr: process\_query() validation \& whitespace strip}
\NormalTok{    Retr{-}\textgreater{}\textgreater{}VStore: similarity\_search(clean\_query, k=5)}
\NormalTok{    VStore{-}{-}\textgreater{}\textgreater{}Retr: Raw tuples [(Doc, distance)]}
\NormalTok{    Retr{-}\textgreater{}\textgreater{}Retr: normalize\_score() (L2 / Cosine {-}\textgreater{} 0.0{-}1.0)}
\NormalTok{    Retr{-}\textgreater{}\textgreater{}Retr: filter\_by\_threshold(score \textgreater{}= 0.20)}
    
\NormalTok{    alt If relevant chunks found}
\NormalTok{        Retr{-}{-}\textgreater{}\textgreater{}UI: Formatted results + metadata}
\NormalTok{        UI{-}\textgreater{}\textgreater{}LLM: invoke(RAG Prompt + Context + Question)}
\NormalTok{        LLM{-}{-}\textgreater{}\textgreater{}UI: Grounded answer with citations}
\NormalTok{        UI{-}\textgreater{}\textgreater{}Student: Display answer + expander with Page references}
\NormalTok{    else If score \textless{} 0.20 (Out of syllabus)}
\NormalTok{        Retr{-}{-}\textgreater{}\textgreater{}UI: No relevant results found}
\NormalTok{        UI{-}\textgreater{}\textgreater{}Student: Fallback: "I could not find enough information in uploaded documents."}
\NormalTok{    end}
\end{Highlighting}
\end{Shaded}

\subsection{4.5 Request and Execution
Lifecycles}\label{request-and-execution-lifecycles}

\subsubsection{A. Interactive Q\&A
Lifecycle}\label{a.-interactive-qa-lifecycle}

\begin{enumerate}
\def\labelenumi{\arabic{enumi}.}
\tightlist
\item
  User enters natural language query into \texttt{st.\allowbreak chat\_\allowbreak input}.
\item
  Query validation checks string length and whitespace.
\item
  Semantic retrieval fetches top-K candidates from FAISS/Chroma.
\item
  Score normalization evaluates relevance against tau = 0.20 threshold.
\item
  Context formatter synthesizes chunk texts with
  \texttt{Source\ \allowbreak X\ \allowbreak (\allowbreak File:\ \allowbreak .\allowbreak .\allowbreak .,\allowbreak \ \allowbreak Page:\ \allowbreak .\allowbreak .\allowbreak .)\allowbreak } headers.
\item
  LangChain LCEL pipeline invokes LLM with deterministic temperature
  (temperature = 0.0).
\item
  UI updates chat history in \texttt{st.\allowbreak session\_\allowbreak state} and appends an
  audit record to \texttt{data/\allowbreak history.\allowbreak json}.
\end{enumerate}

\subsubsection{B. Asynchronous Practice Quiz
Lifecycle}\label{b.-asynchronous-practice-quiz-lifecycle}

\begin{enumerate}
\def\labelenumi{\arabic{enumi}.}
\tightlist
\item
  User specifies a topic (e.g., ``Virtual Memory''), number of questions
  (\(N=5\)), and document scope.
\item
  Background thread (\texttt{worker\_\allowbreak generate\_\allowbreak quiz}) spawns without
  blocking the main UI thread.
\item
  Worker retrieves top-5 context chunks relevant to the topic.
\item
  Prompt template enforces strict JSON array formatting.
\item
  Worker cleans potential markdown formatting
  (\texttt{\textasciigrave{}\textasciigrave{}\textasciigrave{}json}),
  parses JSON, and stores results in \texttt{GLOBAL\_\allowbreak BACKGROUND\_\allowbreak JOBS}.
\item
  History record is created. User can review and take the quiz at any
  time with instant automated grading.
\end{enumerate}

\subsection{4.6 Asynchronous Threading Model and Non-Blocking UI
Execution}\label{asynchronous-threading-model-and-non-blocking-ui-execution}

Streamlit runs on a linear script execution model: when user input
occurs, the entire script reruns from top to bottom. In traditional
implementations, generating a 10-question quiz or comprehensive summary
blocks the Python execution thread for 10-30 seconds, causing browser
freezes.

StudyRAG circumvents this limitation through an asynchronous background
worker architecture: - \textbf{Global Job Registry}: A module-level
dictionary \texttt{GLOBAL\_\allowbreak BACKGROUND\_\allowbreak JOBS} tracks job status
(\texttt{running}, \texttt{completed}, \texttt{failed}), title, results,
and error logs across script re-runs. - \textbf{Execution Context
Binding}: Background threads are initialized using
\texttt{threading.\allowbreak Thread} and wrapped with Streamlit's internal
\texttt{add\_\allowbreak script\_\allowbreak run\_\allowbreak ctx(\allowbreak t)\allowbreak } helper, allowing background threads to
access server resources safely without threading violations. -
\textbf{Asynchronous Polling UI}: When jobs are active, the sidebar
displays a non-blocking indicator: \emph{'' 1 task(s) generating in
background\ldots{} You can switch pages freely while your content
builds!{\kern0pt}``}

\subsection{4.7 Error Propagation and Fault Tolerance
Strategy}\label{error-propagation-and-fault-tolerance-strategy}

\begin{itemize}
\tightlist
\item
  \textbf{Ingestion Failures}: If a user uploads an encrypted,
  password-protected, or corrupted PDF, \texttt{pdf\_\allowbreak loader.\allowbreak py} catches
  the PyMuPDF exception, logs an error via \texttt{logger.\allowbreak error(\allowbreak )\allowbreak }, and
  raises a user-friendly \texttt{ValueError} displayed as a Streamlit
  warning card.
\item
  \textbf{Missing API Keys}: If a user selects Groq or OpenAI without
  defining environment keys, \texttt{LLMManager} raises a
  \texttt{ValueError} indicating the missing key. The UI catches this
  and directs the user to configure \texttt{.\allowbreak env}.
\item
  \textbf{JSON Parsing Errors}: If an LLM returns malformed JSON for a
  quiz or flashcard deck, \texttt{study\_\allowbreak generator.\allowbreak py} logs the raw text
  and gracefully raises a clean exception rather than crashing the
  application. \#\# Chapter 5: Complete Folder and File Structure
\end{itemize}

\subsection{5.1 Verified Repository
Tree}\label{verified-repository-tree}

The complete, active directory layout of the repository is documented
below:

\begin{verbatim}
rag-study-assistant/
    .env                                # Local secrets & API keys (Git-ignored)
    .env.example                        # Template for environment variables
    .gitignore                          # Git tracking exclusions
    .streamlit/                         # Streamlit server configurations
    config/
        settings.yaml                   # Core application defaults & parameters
    data/
        history.json                    # Persistent study activity log (quizzes, cards, chat)
        uploaded/                       # Uploaded study materials (e.g. os.pdf)
        vectorstores/                   # Serialized vector indices
            chroma/                     # ChromaDB SQLite & HNSW collection storage
            faiss/                      # FAISS binary index (index.faiss, index.pkl)
    docs/                               # Project documentation artifacts
    evaluation/                         # Evaluation datasets & benchmark scripts
    src/
        embeddings/
            embedding_manager.py        # HuggingFace / SentenceTransformers manager
        evaluation/
            retrieval_metrics.py        # Precision@K, Recall@K, Hit Rate@K metrics
        generation/
            llm_manager.py              # Groq, OpenAI, and Ollama multi-provider factory
            prompt_templates.py         # LCEL system prompts for RAG, Quizzes, Cards, Summaries
            study_generator.py          # Study artifact generation & JSON parsing engine
        ingestion/
            chunker.py                  # RecursiveCharacterTextSplitter wrapper
            document_cleaner.py         # Whitespace normalizer & empty page stripper
            ingestion_pipeline.py       # End-to-end PDF processing pipeline coordinator
            pdf_loader.py               # PyMuPDF page-by-page text extractor
        rag/
            rag_pipeline.py             # End-to-end RAG orchestrator with fallback logic
        retrieval/
            query_processor.py          # Query cleaner, trimmer & validator
            retrieval_pipeline.py       # End-to-end retrieval coordinator with timing metrics
            retriever.py                # Source-filtering vector search executor
            score_filter.py             # L2/Cosine score normalization & threshold filter
        utils/
            file_utils.py               # SHA-256 file hashing & deduplication
            history_manager.py          # JSON history CRUD operations & UUID tracking
            logger.py                   # Centralized Python logging configuration
        vectorstore/
            chroma_store.py             # Persistent ChromaDB wrapper with cosine metric
            faiss_store.py              # In-memory & serialized FAISS CPU wrapper
            vector_store_manager.py     # Unified factory interface for vector DB backends
    tests/
        test_chunking.py                # Unit tests for text splitting & overlap logic
        test_document_cleaner.py        # Unit tests for whitespace & newline cleaning
        test_embeddings.py              # Integration tests for HuggingFace embeddings
        test_llm_manager.py             # Unit tests for multi-provider LLM initializations
        test_pdf_loader.py              # Unit tests for PyMuPDF loading & empty PDF handling
        test_query_processor.py         # Unit tests for query sanitization & length limits
        test_rag.py                     # Integration tests for RAG pipeline & fallback logic
        test_retrieval.py               # Unit tests for score normalization & metrics
        test_vectorstore.py             # Integration tests for FAISS and Chroma CRUD ops
    app.py                              # Main interactive Streamlit application
    pytest.ini                          # Pytest configuration (pythonpath = .)
    requirements.txt                    # Python dependency manifest
    README.md                           # Public repository overview
\end{verbatim}

\subsection{5.2 Exhaustive Directory and File
Catalog}\label{exhaustive-directory-and-file-catalog}

\begingroup\scriptsize
\begin{longtable}[]{@{}
  >{\raggedright\arraybackslash\sloppy}p{(\linewidth - 8\tabcolsep) * \real{0.2600}}
  >{\raggedright\arraybackslash\sloppy}p{(\linewidth - 8\tabcolsep) * \real{0.1400}}
  >{\raggedright\arraybackslash\sloppy}p{(\linewidth - 8\tabcolsep) * \real{0.3000}}
  >{\raggedright\arraybackslash\sloppy}p{(\linewidth - 8\tabcolsep) * \real{0.1600}}
  >{\raggedright\arraybackslash\sloppy}p{(\linewidth - 8\tabcolsep) * \real{0.1000}}@{}}
\toprule\noalign{}
\begin{minipage}[b]{\linewidth}\raggedright
File / Directory Path
\end{minipage} & \begin{minipage}[b]{\linewidth}\raggedright
Role \& Category
\end{minipage} & \begin{minipage}[b]{\linewidth}\raggedright
Primary Responsibility
\end{minipage} & \begin{minipage}[b]{\linewidth}\raggedright
Dependencies
\end{minipage} & \begin{minipage}[b]{\linewidth}\raggedright
Execution Criticality
\end{minipage} \\
\midrule\noalign{}
\endhead
\bottomrule\noalign{}
\endlastfoot
\texttt{app.\allowbreak py} & UI Entry Point & Implements Streamlit UI, routing,
theme CSS, background workers & All \texttt{src.*} modules &
\textbf{CRITICAL} (Required to run app) \\
\texttt{config/\allowbreak settings.\allowbreak yaml} & Configuration & Defines default chunk
sizes, overlap, top\_k, and thresholds & PyYAML & \textbf{CRITICAL}
(Loaded at startup) \\
\texttt{.\allowbreak env} & Environment Config & Contains API keys
(\texttt{GROQ\_\allowbreak API\_\allowbreak KEY}, \texttt{OPENAI\_\allowbreak API\_\allowbreak KEY}) & python-dotenv &
\textbf{HIGH} (Required for cloud LLMs) \\
\texttt{requirements.\allowbreak txt} & Package Manifest & Declares required Python
package versions & pip & \textbf{CRITICAL} (Required for
installation) \\
\texttt{pytest.\allowbreak ini} & Test Configuration & Defines pytest root directory
and pythonpath & pytest & \textbf{HIGH} (Required to run tests
cleanly) \\
\texttt{src/\allowbreak ingestion/\allowbreak pdf\_\allowbreak loader.\allowbreak py} & Document Processing & Extracts
text and 1-indexed page metadata via PyMuPDF & pymupdf, pathlib &
\textbf{CRITICAL} (Required for document loading) \\
\texttt{src/\allowbreak ingestion/\allowbreak document\_\allowbreak cleaner.\allowbreak py} & Text Cleaning & Normalizes
spaces and removes empty pages & re & \textbf{CRITICAL} (Prepares text
for chunking) \\
\texttt{src/\allowbreak ingestion/\allowbreak chunker.\allowbreak py} & Text Splitting & Splits text
recursively preserving metadata and chunk IDs & langchain-text-splitters
& \textbf{CRITICAL} (Creates vector chunks) \\
\texttt{src/\allowbreak ingestion/\allowbreak ingestion\_\allowbreak pipeline.\allowbreak py} & Pipeline Coordinator &
Orchestrates loading, cleaning, and chunking & \texttt{pdf\_\allowbreak loader},
\texttt{cleaner}, \texttt{chunker} & \textbf{CRITICAL} (Main ingestion
API) \\
\texttt{src/\allowbreak embeddings/\allowbreak embedding\_\allowbreak manager.\allowbreak py} & Vector Embedding & Loads
and caches SentenceTransformer models & langchain-huggingface &
\textbf{CRITICAL} (Vectorization engine) \\
\texttt{src/\allowbreak vectorstore/\allowbreak vector\_\allowbreak store\_\allowbreak manager.\allowbreak py} & Storage Factory &
Provides unified API for FAISS and ChromaDB & \texttt{faiss\_\allowbreak store},
\texttt{chroma\_\allowbreak store} & \textbf{CRITICAL} (Database abstraction) \\
\texttt{src/\allowbreak vectorstore/\allowbreak faiss\_\allowbreak store.\allowbreak py} & Vector Index & Manages
in-memory and disk FAISS L2 index & faiss-cpu & \textbf{CRITICAL}
(Default vector database) \\
\texttt{src/\allowbreak vectorstore/\allowbreak chroma\_\allowbreak store.\allowbreak py} & Vector Index & Manages
persistent SQLite Chroma collection & chromadb & \textbf{HIGH}
(Alternative vector database) \\
\texttt{src/\allowbreak retrieval/\allowbreak query\_\allowbreak processor.\allowbreak py} & Query Sanitizer & Trims and
validates user input strings & re & \textbf{HIGH} (Guards retrieval
pipeline) \\
\texttt{src/\allowbreak retrieval/\allowbreak score\_\allowbreak filter.\allowbreak py} & Mathematics & Normalizes
distance metrics to \([0, 1]\) and applies \$ au\$ & math &
\textbf{CRITICAL} (Determines answer relevance) \\
\texttt{src/\allowbreak retrieval/\allowbreak retriever.\allowbreak py} & Vector Search & Executes
similarity search with source document filters &
\texttt{VectorStoreManager} & \textbf{CRITICAL} (Executes search
queries) \\
\texttt{src/\allowbreak retrieval/\allowbreak retrieval\_\allowbreak pipeline.\allowbreak py} & Pipeline Coordinator &
Combines query processing, search, filtering, and timers &
\texttt{query\_\allowbreak proc}, \texttt{retriever}, \texttt{filter} &
\textbf{CRITICAL} (Main retrieval API) \\
\texttt{src/\allowbreak generation/\allowbreak prompt\_\allowbreak templates.\allowbreak py} & Prompt Engineering &
Declares ChatPromptTemplates and context formatters & langchain-core &
\textbf{CRITICAL} (Formats LLM inputs) \\
\texttt{src/\allowbreak generation/\allowbreak llm\_\allowbreak manager.\allowbreak py} & Model Factory & Initializes
Groq, OpenAI, or Ollama model instances & langchain-openai,
langchain-ollama & \textbf{CRITICAL} (LLM communication) \\
\texttt{src/\allowbreak generation/\allowbreak study\_\allowbreak generator.\allowbreak py} & Generator Engine &
Generates quizzes, flashcards, summaries; parses JSON &
\texttt{llm\_\allowbreak manager}, \texttt{prompts}, \texttt{retrieval} &
\textbf{CRITICAL} (Study tool generation) \\
\texttt{src/\allowbreak rag/\allowbreak rag\_\allowbreak pipeline.\allowbreak py} & RAG Orchestrator & Connects
retrieval, prompt formatting, and LLM output & \texttt{retrieval\_\allowbreak pipe},
\texttt{prompts}, LCEL & \textbf{CRITICAL} (Core Q\&A pipeline) \\
\texttt{src/\allowbreak evaluation/\allowbreak retrieval\_\allowbreak metrics.\allowbreak py} & Evaluation Engine &
Calculates Precision@K, Recall@K, and Hit Rate@K & typing &
\textbf{MEDIUM} (Benchmarking \& evaluation) \\
\texttt{src/\allowbreak utils/\allowbreak file\_\allowbreak utils.\allowbreak py} & Utilities & Computes SHA-256 hashes
for document deduplication & hashlib & \textbf{HIGH} (Prevents duplicate
processing) \\
\texttt{src/\allowbreak utils/\allowbreak history\_\allowbreak manager.\allowbreak py} & Persistence & Manages JSON
read/write/delete operations for history & json, uuid & \textbf{HIGH}
(Study session history) \\
\texttt{src/\allowbreak utils/\allowbreak logger.\allowbreak py} & System Diagnostics & Configures console
logging across all modules & logging, sys & \textbf{HIGH} (Debugging \&
system audit) \\
\texttt{data/\allowbreak history.\allowbreak json} & Data Store & Persistent JSON log of all
study artifacts & Built-in JSON & \textbf{HIGH} (Runtime persistence
datastore) \\
\texttt{data/\allowbreak uploaded/\allowbreak } & Storage Directory & Directory storing
user-uploaded PDF files & OS File system & \textbf{HIGH} (Document
archive) \\
\texttt{data/\allowbreak vectorstores/\allowbreak } & Storage Directory & Directory storing
serialized FAISS/Chroma indices & OS File system & \textbf{HIGH} (Vector
index storage) \\
\end{longtable}
\endgroup

\section{Chapter 6: Module-by-Module
Implementation}\label{chapter-6-module-by-module-implementation}

\subsection{\texorpdfstring{6.1 Module: Document Ingestion and
Preprocessing
(\texttt{src/\allowbreak ingestion})}{6.1 Module: Document Ingestion and Preprocessing (src/ingestion)}}\label{module-document-ingestion-and-preprocessing-srcingestion}

The ingestion subsystem converts raw PDF documents into semantically
coherent, metadata-rich text chunks ready for vectorization.

\subsubsection{\texorpdfstring{1.
\texttt{pdf\_\allowbreak loader.\allowbreak py}}{1. pdf\_loader.py}}\label{pdf_loader.py}

\begin{itemize}
\tightlist
\item
  \textbf{Function:}
  \texttt{load\_\allowbreak pdf(\allowbreak pdf\_\allowbreak path:\ \allowbreak str\ \allowbreak \textbar{}\ \allowbreak Path)\allowbreak \ \allowbreak -\textgreater{}\allowbreak \ \allowbreak List{[}Dict{[}str,\allowbreak \ \allowbreak Any{]}{]}}
\item
  \textbf{Implementation Logic:}

  \begin{itemize}
  \tightlist
  \item
    Validates file existence using \texttt{pathlib.\allowbreak Path.\allowbreak exists(\allowbreak )\allowbreak }.
  \item
    Enforces \texttt{.\allowbreak pdf} extension check, raising \texttt{ValueError}
    if non-PDF.
  \item
    Opens document using \texttt{pymupdf.\allowbreak open(\allowbreak path)\allowbreak }.
  \item
    Iterates through \texttt{doc} page-by-page. For each page, executes
    \texttt{page.\allowbreak get\_\allowbreak text(\allowbreak "text")\allowbreak }.
  \item
    Tracks boolean flag \texttt{text\_\allowbreak found} across all pages. If no
    extractable text is detected, logs a warning indicating a potential
    scanned image.
  \item
    Stores each page as a dictionary:
    \texttt{\{"text":\ \allowbreak text,\allowbreak \ \allowbreak "source":\ \allowbreak path.\allowbreak name,\allowbreak \ \allowbreak "page":\ \allowbreak page\_\allowbreak num\ \allowbreak +\ \allowbreak 1\}}
    (1-indexed for human readability).
  \end{itemize}
\item
  \textbf{Error Handling:} Catches corrupt/encrypted files, logs via
  \texttt{logger.\allowbreak error}, and raises
  \texttt{ValueError(\allowbreak "Unable\ \allowbreak to\ \allowbreak open\ \allowbreak PDF.\allowbreak .\allowbreak .\ \allowbreak It\ \allowbreak may\ \allowbreak be\ \allowbreak corrupted\ \allowbreak or\ \allowbreak encrypted.")\allowbreak }.
\end{itemize}

\subsubsection{\texorpdfstring{2.
\texttt{document\_\allowbreak cleaner.\allowbreak py}}{2. document\_cleaner.py}}\label{document_cleaner.py}

\begin{itemize}
\tightlist
\item
  \textbf{Functions:}

  \begin{itemize}
  \tightlist
  \item
    \texttt{clean\_\allowbreak text(\allowbreak text:\ \allowbreak str)\allowbreak \ \allowbreak -\textgreater{}\allowbreak \ \allowbreak str}: Uses regex
    \texttt{re.\allowbreak sub(\allowbreak r\textquotesingle{}{[}\ \allowbreak \textbackslash{}\textbackslash{}t{]}+\textquotesingle{},\allowbreak \ \allowbreak \textquotesingle{}\ \allowbreak \textquotesingle{},\allowbreak \ \allowbreak text)\allowbreak }
    to collapse inline whitespace. Uses
    \texttt{re.\allowbreak sub(\allowbreak r\textquotesingle{}\textbackslash{}\textbackslash{}n\{3,\allowbreak \}\textquotesingle{},\allowbreak \ \allowbreak \textquotesingle{}\textbackslash{}\textbackslash{}n\textbackslash{}\textbackslash{}n\textquotesingle{},\allowbreak \ \allowbreak text)\allowbreak }
    to compress excess blank lines into standard two-newline paragraph
    separators. Strips leading and trailing whitespace.
  \item
    \texttt{clean\_\allowbreak documents(\allowbreak documents:\ \allowbreak List{[}Dict{[}str,\allowbreak \ \allowbreak Any{]}{]})\allowbreak \ \allowbreak -\textgreater{}\allowbreak \ \allowbreak List{[}Dict{[}str,\allowbreak \ \allowbreak Any{]}{]}}:
    Applies \texttt{clean\_\allowbreak text} to all documents. Filters out any
    document whose resulting text is empty
    (\texttt{if\ \allowbreak cleaned\_\allowbreak text:}), removing blank pages.
  \end{itemize}
\end{itemize}

\subsubsection{\texorpdfstring{3.
\texttt{chunker.\allowbreak py}}{3. chunker.py}}\label{chunker.py}

\begin{itemize}
\tightlist
\item
  \textbf{Function:}
  \texttt{create\_\allowbreak chunks(\allowbreak documents:\ \allowbreak List{[}Dict{[}str,\allowbreak \ \allowbreak Any{]}{]},\allowbreak \ \allowbreak chunk\_\allowbreak size:\ \allowbreak int\ \allowbreak =\allowbreak \ \allowbreak 500,\allowbreak \ \allowbreak chunk\_\allowbreak overlap:\ \allowbreak int\ \allowbreak =\allowbreak \ \allowbreak 50)\allowbreak \ \allowbreak -\textgreater{}\allowbreak \ \allowbreak List{[}Dict{[}str,\allowbreak \ \allowbreak Any{]}{]}}
\item
  \textbf{Validation Rules:}

  \begin{itemize}
  \tightlist
  \item
    Enforces \texttt{chunk\_\allowbreak size\ \allowbreak \textgreater{}\ \allowbreak 0}.
  \item
    Enforces \texttt{chunk\_\allowbreak overlap\ \allowbreak \textgreater{}=\allowbreak \ \allowbreak 0}.
  \item
    Enforces \texttt{chunk\_\allowbreak overlap\ \allowbreak \textless{}\ \allowbreak chunk\_\allowbreak size} (strictly
    less than chunk size).
  \end{itemize}
\item
  \textbf{Splitting Logic:}

  \begin{itemize}
  \tightlist
  \item
    Initializes
    \texttt{RecursiveCharacterTextSplitter(\allowbreak chunk\_\allowbreak size=\allowbreak chunk\_\allowbreak size,\allowbreak \ \allowbreak chunk\_\allowbreak overlap=\allowbreak chunk\_\allowbreak overlap,\allowbreak \ \allowbreak separators=\allowbreak {[}"\textbackslash{}\textbackslash{}n\textbackslash{}\textbackslash{}n",\allowbreak \ \allowbreak "\textbackslash{}\textbackslash{}n",\allowbreak \ \allowbreak ".\ \allowbreak ",\allowbreak \ \allowbreak "\ \allowbreak ",\allowbreak \ \allowbreak ""{]})\allowbreak }.
  \item
    Iterates through documents, executing
    \texttt{text\_\allowbreak splitter.\allowbreak split\_\allowbreak text(\allowbreak doc{[}"text"{]})\allowbreak }.
  \item
    For each generated chunk, appends metadata: \texttt{source},
    \texttt{page}, and an incrementing integer \texttt{chunk\_\allowbreak id}.
  \end{itemize}
\end{itemize}

\subsubsection{\texorpdfstring{4.
\texttt{ingestion\_\allowbreak pipeline.\allowbreak py}}{4. ingestion\_pipeline.py}}\label{ingestion_pipeline.py}

\begin{itemize}
\tightlist
\item
  \textbf{Function:}
  \texttt{process\_\allowbreak pdf(\allowbreak pdf\_\allowbreak path:\ \allowbreak str\ \allowbreak \textbar{}\ \allowbreak Path,\allowbreak \ \allowbreak chunk\_\allowbreak size:\ \allowbreak int\ \allowbreak =\allowbreak \ \allowbreak 500,\allowbreak \ \allowbreak chunk\_\allowbreak overlap:\ \allowbreak int\ \allowbreak =\allowbreak \ \allowbreak 50)\allowbreak \ \allowbreak -\textgreater{}\allowbreak \ \allowbreak List{[}Dict{[}str,\allowbreak \ \allowbreak Any{]}{]}}
\item
  \textbf{Workflow:} Combines \texttt{load\_\allowbreak pdf} \(\rightarrow\)
  \texttt{clean\_\allowbreak documents} \(\rightarrow\) \texttt{create\_\allowbreak chunks}.
  Logs chunk count and processing status.
\end{itemize}

\begin{center}\rule{0.5\linewidth}{0.5pt}\end{center}

\subsection{\texorpdfstring{6.2 Module: Semantic Embedding Management
(\texttt{src/\allowbreak embeddings})}{6.2 Module: Semantic Embedding Management (src/embeddings)}}\label{module-semantic-embedding-management-srcembeddings}

\subsubsection{\texorpdfstring{\texttt{embedding\_\allowbreak manager.\allowbreak py}}{embedding\_manager.py}}\label{embedding_manager.py}

\begin{itemize}
\tightlist
\item
  \textbf{Class:}
  \texttt{EmbeddingManager(\allowbreak model\_\allowbreak key:\ \allowbreak str\ \allowbreak =\allowbreak \ \allowbreak "minilm")\allowbreak }
\item
  \textbf{Model Registry:}

  \begin{itemize}
  \tightlist
  \item
    \texttt{"minilm"}: maps to
    \texttt{sentence-transformers/\allowbreak all-MiniLM-L6-v2}
  \item
    \texttt{"bge"}: maps to \texttt{BAAI/\allowbreak bge-small-en-v1.5}
  \end{itemize}
\item
  \textbf{Methods:}

  \begin{itemize}
  \tightlist
  \item
    \texttt{get\_\allowbreak embeddings(\allowbreak )\allowbreak \ \allowbreak -\textgreater{}\allowbreak \ \allowbreak HuggingFaceEmbeddings}:
    Returns the underlying LangChain embedding object.
  \item
    \texttt{get\_\allowbreak model\_\allowbreak name(\allowbreak )\allowbreak \ \allowbreak -\textgreater{}\allowbreak \ \allowbreak str}: Returns the
    active HuggingFace repository string.
  \item
    \texttt{get\_\allowbreak embedding\_\allowbreak dimension(\allowbreak )\allowbreak \ \allowbreak -\textgreater{}\allowbreak \ \allowbreak int}: Executes
    a single test query embedding
    (\texttt{self.\allowbreak embeddings.\allowbreak embed\_\allowbreak query(\allowbreak "test")\allowbreak }), computes
    \texttt{len(\allowbreak vector)\allowbreak }, and caches \texttt{self.\_\allowbreak dimension} (384
    dimensions).
  \end{itemize}
\end{itemize}

\begin{center}\rule{0.5\linewidth}{0.5pt}\end{center}

\subsection{\texorpdfstring{6.3 Module: Vector Store Subsystem
(\texttt{src/\allowbreak vectorstore})}{6.3 Module: Vector Store Subsystem (src/vectorstore)}}\label{module-vector-store-subsystem-srcvectorstore}

\subsubsection{\texorpdfstring{1.
\texttt{faiss\_\allowbreak store.\allowbreak py}}{1. faiss\_store.py}}\label{faiss_store.py}

\begin{itemize}
\tightlist
\item
  \textbf{Class:}
  \texttt{FAISSStore(\allowbreak embeddings:\ \allowbreak Embeddings,\allowbreak \ \allowbreak persist\_\allowbreak directory:\ \allowbreak str\ \allowbreak =\allowbreak \ \allowbreak "data/\allowbreak vectorstores/\allowbreak faiss/\allowbreak ")\allowbreak }
\item
  \textbf{Methods:}

  \begin{itemize}
  \tightlist
  \item
    \texttt{add\_\allowbreak documents(\allowbreak documents:\ \allowbreak List{[}Document{]})\allowbreak }: If
    \texttt{self.\allowbreak vector\_\allowbreak store} is \texttt{None}, initializes
    \texttt{FAISS.\allowbreak from\_\allowbreak documents(\allowbreak documents,\allowbreak \ \allowbreak self.\allowbreak embeddings)\allowbreak }. If
    already initialized, executes
    \texttt{self.\allowbreak vector\_\allowbreak store.\allowbreak add\_\allowbreak documents(\allowbreak documents)\allowbreak }.
  \item
    \texttt{save(\allowbreak )\allowbreak }: Creates directory and calls
    \texttt{self.\allowbreak vector\_\allowbreak store.\allowbreak save\_\allowbreak local(\allowbreak str(\allowbreak self.\allowbreak persist\_\allowbreak directory)\allowbreak )\allowbreak }.
  \item
    \texttt{load(\allowbreak )\allowbreak \ \allowbreak -\textgreater{}\allowbreak \ \allowbreak bool}: Verifies
    \texttt{index.\allowbreak faiss} exists. Calls
    \texttt{FAISS.\allowbreak load\_\allowbreak local(\allowbreak .\allowbreak .\allowbreak .,\allowbreak \ \allowbreak allow\_\allowbreak dangerous\_\allowbreak deserialization=\allowbreak True)\allowbreak }.
    Returns \texttt{True} on success.
  \item
    \texttt{similarity\_\allowbreak search(\allowbreak query:\ \allowbreak str,\allowbreak \ \allowbreak k:\ \allowbreak int\ \allowbreak =\allowbreak \ \allowbreak 5)\allowbreak \ \allowbreak -\textgreater{}\allowbreak \ \allowbreak List{[}tuple{[}Document,\allowbreak \ \allowbreak float{]}{]}}:
    Calls \texttt{similarity\_\allowbreak search\_\allowbreak with\_\allowbreak score(\allowbreak query,\allowbreak \ \allowbreak k=\allowbreak k)\allowbreak }. Returns
    raw L2 distances.
  \item
    \texttt{reset(\allowbreak )\allowbreak }: Sets \texttt{vector\_\allowbreak store\ \allowbreak =\allowbreak \ \allowbreak None} and removes
    directory using \texttt{shutil.\allowbreak rmtree(\allowbreak )\allowbreak }.
  \end{itemize}
\end{itemize}

\subsubsection{\texorpdfstring{2.
\texttt{chroma\_\allowbreak store.\allowbreak py}}{2. chroma\_store.py}}\label{chroma_store.py}

\begin{itemize}
\tightlist
\item
  \textbf{Class:}
  \texttt{ChromaStore(\allowbreak embeddings:\ \allowbreak Embeddings,\allowbreak \ \allowbreak persist\_\allowbreak directory:\ \allowbreak str\ \allowbreak =\allowbreak \ \allowbreak "data/\allowbreak vectorstores/\allowbreak chroma/\allowbreak ")\allowbreak }
\item
  \textbf{Methods:}

  \begin{itemize}
  \tightlist
  \item
    \texttt{add\_\allowbreak documents(\allowbreak documents:\ \allowbreak List{[}Document{]})\allowbreak }: Initializes
    Chroma with collection metadata
    \texttt{collection\_\allowbreak metadata=\allowbreak \{"hnsw:space":\ \allowbreak "cosine"\}}.
  \item
    \texttt{load(\allowbreak )\allowbreak \ \allowbreak -\textgreater{}\allowbreak \ \allowbreak bool}: Verifies
    \texttt{chroma.\allowbreak sqlite3} exists and connects.
  \item
    \texttt{similarity\_\allowbreak search(\allowbreak query:\ \allowbreak str,\allowbreak \ \allowbreak k:\ \allowbreak int\ \allowbreak =\allowbreak \ \allowbreak 5)\allowbreak \ \allowbreak -\textgreater{}\allowbreak \ \allowbreak List{[}tuple{[}Document,\allowbreak \ \allowbreak float{]}{]}}:
    Executes similarity search returning cosine distances.
  \item
    \texttt{reset(\allowbreak )\allowbreak }: Calls \texttt{delete\_\allowbreak collection(\allowbreak )\allowbreak } and deletes
    the persist directory.
  \end{itemize}
\end{itemize}

\subsubsection{\texorpdfstring{3.
\texttt{vector\_\allowbreak store\_\allowbreak manager.\allowbreak py}}{3. vector\_store\_manager.py}}\label{vector_store_manager.py}

\begin{itemize}
\tightlist
\item
  \textbf{Class:}
  \texttt{VectorStoreManager(\allowbreak provider:\ \allowbreak str,\allowbreak \ \allowbreak embeddings:\ \allowbreak Embeddings,\allowbreak \ \allowbreak persist\_\allowbreak directory\_\allowbreak base:\ \allowbreak str\ \allowbreak =\allowbreak \ \allowbreak "data/\allowbreak vectorstores")\allowbreak }
\item
  \textbf{Role:} Factory class providing unified routing between FAISS
  and Chroma. Provides \texttt{raw\_\allowbreak chunks\_\allowbreak to\_\allowbreak documents(\allowbreak chunks)\allowbreak } to
  convert raw chunk dictionaries into LangChain \texttt{Document}
  objects with metadata preservation.
\end{itemize}

\begin{center}\rule{0.5\linewidth}{0.5pt}\end{center}

\subsection{\texorpdfstring{6.4 Module: Context Retrieval Engine
(\texttt{src/\allowbreak retrieval})}{6.4 Module: Context Retrieval Engine (src/retrieval)}}\label{module-context-retrieval-engine-srcretrieval}

\subsubsection{\texorpdfstring{1.
\texttt{query\_\allowbreak processor.\allowbreak py}}{1. query\_processor.py}}\label{query_processor.py}

\begin{itemize}
\tightlist
\item
  \textbf{Function:}
  \texttt{process\_\allowbreak query(\allowbreak query:\ \allowbreak str,\allowbreak \ \allowbreak max\_\allowbreak query\_\allowbreak length:\ \allowbreak int\ \allowbreak =\allowbreak \ \allowbreak 1000)\allowbreak \ \allowbreak -\textgreater{}\allowbreak \ \allowbreak str}
\item
  \textbf{Validations:} Rejects empty or whitespace-only queries
  (\texttt{raise\ \allowbreak ValueError(\allowbreak "Please\ \allowbreak enter\ \allowbreak a\ \allowbreak question.")\allowbreak }). Uses
  regex
  \texttt{re.\allowbreak sub(\allowbreak r\textquotesingle{}\textbackslash{}\textbackslash{}s+\textquotesingle{},\allowbreak \ \allowbreak \textquotesingle{}\ \allowbreak \textquotesingle{},\allowbreak \ \allowbreak query)\allowbreak .\allowbreak strip(\allowbreak )\allowbreak }
  to collapse internal whitespace. Enforces length \(\le 1000\)
  characters.
\end{itemize}

\subsubsection{\texorpdfstring{2.
\texttt{score\_\allowbreak filter.\allowbreak py}}{2. score\_filter.py}}\label{score_filter.py}

\begin{itemize}
\tightlist
\item
  \textbf{Function:}
  \texttt{normalize\_\allowbreak score(\allowbreak raw\_\allowbreak score:\ \allowbreak float,\allowbreak \ \allowbreak provider:\ \allowbreak str)\allowbreak \ \allowbreak -\textgreater{}\allowbreak \ \allowbreak float}

  \begin{itemize}
  \tightlist
  \item
    \textbf{FAISS (L2 Distance):} Returns
    \(\frac{1.0}{1.0 + \text{raw\_score}}\) (\(0\) distance
    \(\rightarrow 1.0\) relevance).
  \item
    \textbf{Chroma (Cosine Distance):} Returns
    \(\max(0.0, \min(1.0, 1.0 - \text{raw\_score}))\).
  \end{itemize}
\item
  \textbf{Function:}
  \texttt{filter\_\allowbreak by\_\allowbreak threshold(\allowbreak results,\allowbreak \ \allowbreak provider,\allowbreak \ \allowbreak threshold=\allowbreak None)\allowbreak \ \allowbreak -\textgreater{}\allowbreak \ \allowbreak List{[}Dict{[}str,\allowbreak \ \allowbreak Any{]}{]}}:
  Normalizes scores, filters out items below \texttt{threshold}, and
  sorts candidates in descending order of relevance.
\end{itemize}

\subsubsection{\texorpdfstring{3.
\texttt{retriever.\allowbreak py}}{3. retriever.py}}\label{retriever.py}

\begin{itemize}
\tightlist
\item
  \textbf{Class:}
  \texttt{Retriever(\allowbreak vector\_\allowbreak store:\ \allowbreak VectorStoreManager,\allowbreak \ \allowbreak top\_\allowbreak k:\ \allowbreak int\ \allowbreak =\allowbreak \ \allowbreak 5)\allowbreak }
\item
  \textbf{Method:}
  \texttt{retrieve\_\allowbreak with\_\allowbreak scores(\allowbreak query:\ \allowbreak str,\allowbreak \ \allowbreak source\_\allowbreak filter:\ \allowbreak Optional{[}str{]}\ \allowbreak =\allowbreak \ \allowbreak None)\allowbreak }:
  Executes vector similarity search. If a \texttt{source\_\allowbreak filter} is
  specified (and not ``All Documents''), passes native metadata filters
  to the vector store or applies post-retrieval filtering.
\end{itemize}

\subsubsection{\texorpdfstring{4.
\texttt{retrieval\_\allowbreak pipeline.\allowbreak py}}{4. retrieval\_pipeline.py}}\label{retrieval_pipeline.py}

\begin{itemize}
\tightlist
\item
  \textbf{Function:}
  \texttt{retrieve\_\allowbreak context(\allowbreak query,\allowbreak \ \allowbreak vector\_\allowbreak store,\allowbreak \ \allowbreak top\_\allowbreak k=\allowbreak 5,\allowbreak \ \allowbreak score\_\allowbreak threshold=\allowbreak None,\allowbreak \ \allowbreak source\_\allowbreak filter=\allowbreak None)\allowbreak \ \allowbreak -\textgreater{}\allowbreak \ \allowbreak Dict{[}str,\allowbreak \ \allowbreak Any{]}}
\item
  \textbf{Telemetry:} Measures search latency (\texttt{search\_\allowbreak time}),
  filtering latency (\texttt{filter\_\allowbreak time}), and total retrieval time
  (\texttt{total\_\allowbreak time}). Returns structured dictionary with status flag
  \texttt{found:\ \allowbreak bool}.
\end{itemize}

\begin{center}\rule{0.5\linewidth}{0.5pt}\end{center}

\subsection{\texorpdfstring{6.5 Module: LLM Management and Generation
Pipelines
(\texttt{src/\allowbreak generation})}{6.5 Module: LLM Management and Generation Pipelines (src/generation)}}\label{module-llm-management-and-generation-pipelines-srcgeneration}

\subsubsection{\texorpdfstring{1.
\texttt{prompt\_\allowbreak templates.\allowbreak py}}{1. prompt\_templates.py}}\label{prompt_templates.py}

Contains all ChatPromptTemplates utilized across the application: -
\texttt{RAG\_\allowbreak SYSTEM\_\allowbreak PROMPT}: Instructs the AI tutor to answer questions
using ONLY provided context, to cite sources and page numbers, and to
explicitly state \emph{``I could not find enough information in the
uploaded documents to answer this question.''} if context is
insufficient. -
\texttt{format\_\allowbreak context\_\allowbreak for\_\allowbreak prompt(\allowbreak retrieved\_\allowbreak results)\allowbreak }: Formats
chunks into clean textual context blocks:
\texttt{-\/\allowbreak -\/\allowbreak -\ \allowbreak Source\ \allowbreak \{idx\}\ \allowbreak (\allowbreak File:\ \allowbreak \{source\},\allowbreak \ \allowbreak Page:\ \allowbreak \{page\})\allowbreak \ \allowbreak -\/\allowbreak -\/\allowbreak -\ \allowbreak \ \allowbreak \ \allowbreak \{text\}}
- \texttt{get\_\allowbreak quiz\_\allowbreak prompt\_\allowbreak template(\allowbreak )\allowbreak }: Instructs LLM to generate
\(N\) MCQs formatted strictly as a JSON array of
\texttt{\{"question":\ \allowbreak ".\allowbreak .\allowbreak .",\allowbreak \ \allowbreak "options":\ \allowbreak {[}"A.\ \allowbreak .\allowbreak .\allowbreak .",\allowbreak \ \allowbreak "B.\ \allowbreak .\allowbreak .\allowbreak .",\allowbreak \ \allowbreak "C.\ \allowbreak .\allowbreak .\allowbreak .",\allowbreak \ \allowbreak "D.\ \allowbreak .\allowbreak .\allowbreak ."{]},\allowbreak \ \allowbreak "answer":\ \allowbreak "A.\ \allowbreak .\allowbreak .\allowbreak ."\}}.
- \texttt{get\_\allowbreak flashcard\_\allowbreak prompt\_\allowbreak template(\allowbreak )\allowbreak }: Instructs LLM to generate
\(N\) flashcards formatted strictly as a JSON array of
\texttt{\{"front":\ \allowbreak ".\allowbreak .\allowbreak .",\allowbreak \ \allowbreak "back":\ \allowbreak ".\allowbreak .\allowbreak ."\}}. -
\texttt{get\_\allowbreak summary\_\allowbreak prompt\_\allowbreak template(\allowbreak )\allowbreak }: Instructs LLM to generate
structured Markdown with \texttt{\#\#\#\ \allowbreak Summary} and
\texttt{\#\#\#\ \allowbreak Key\ \allowbreak Points}.

\subsubsection{\texorpdfstring{2.
\texttt{llm\_\allowbreak manager.\allowbreak py}}{2. llm\_manager.py}}\label{llm_manager.py}

\begin{itemize}
\tightlist
\item
  \textbf{Class:}
  \texttt{LLMManager(\allowbreak provider:\ \allowbreak str\ \allowbreak =\allowbreak \ \allowbreak "ollama",\allowbreak \ \allowbreak model\_\allowbreak name:\ \allowbreak str\ \allowbreak =\allowbreak \ \allowbreak "llama3",\allowbreak \ \allowbreak temperature:\ \allowbreak float\ \allowbreak =\allowbreak \ \allowbreak 0.0)\allowbreak }
\item
  \textbf{Provider Routing:}

  \begin{itemize}
  \tightlist
  \item
    \texttt{"ollama"}: Connects to local \texttt{ChatOllama} at
    \texttt{OLLAMA\_\allowbreak BASE\_\allowbreak URL} (default
    \texttt{http:/\allowbreak /\allowbreak localhost:11434}).
  \item
    \texttt{"openai"}: Reads \texttt{OPENAI\_\allowbreak API\_\allowbreak KEY} and initializes
    \texttt{ChatOpenAI}.
  \item
    \texttt{"groq"}: Reads \texttt{GROQ\_\allowbreak API\_\allowbreak KEY}. Reuses
    \texttt{ChatOpenAI} pointed to
    \texttt{base\_\allowbreak url=\allowbreak "https:/\allowbreak /\allowbreak api.\allowbreak groq.\allowbreak com/\allowbreak openai/\allowbreak v1"}. Maps deprecated
    model identifiers (e.g., \texttt{llama3-8b-8192}) to active Groq
    endpoints (\texttt{openai/\allowbreak gpt-oss-20b}).
  \end{itemize}
\end{itemize}

\subsubsection{\texorpdfstring{3.
\texttt{study\_\allowbreak generator.\allowbreak py}}{3. study\_generator.py}}\label{study_generator.py}

\begin{itemize}
\tightlist
\item
  Helper function
  \texttt{\_\allowbreak clean\_\allowbreak json\_\allowbreak output(\allowbreak text:\ \allowbreak str)\allowbreak \ \allowbreak -\textgreater{}\allowbreak \ \allowbreak str}:
  Strips markdown wrappers (\texttt{json\ \allowbreak .\allowbreak .\allowbreak .}) from LLM output.
\item
  \texttt{generate\_\allowbreak quiz(\allowbreak .\allowbreak .\allowbreak .)\allowbreak }: Retrieves context for topic, executes
  LLM chain, cleans output, parses JSON, and returns quiz list.
\item
  \texttt{generate\_\allowbreak flashcards(\allowbreak .\allowbreak .\allowbreak .)\allowbreak }: Retrieves context, executes LLM
  chain, cleans output, parses JSON, and returns card list.
\item
  \texttt{generate\_\allowbreak summary(\allowbreak .\allowbreak .\allowbreak .)\allowbreak }: Retrieves context (top-8 chunks),
  executes LLM chain, and returns Markdown summary string.
\end{itemize}

\begin{center}\rule{0.5\linewidth}{0.5pt}\end{center}

\subsection{\texorpdfstring{6.6 Module: Core RAG Pipeline
(\texttt{src/\allowbreak rag})}{6.6 Module: Core RAG Pipeline (src/rag)}}\label{module-core-rag-pipeline-srcrag}

\subsubsection{\texorpdfstring{\texttt{rag\_\allowbreak pipeline.\allowbreak py}}{rag\_pipeline.py}}\label{rag_pipeline.py}

\begin{itemize}
\tightlist
\item
  \textbf{Function:}
  \texttt{generate\_\allowbreak rag\_\allowbreak response(\allowbreak query,\allowbreak \ \allowbreak vector\_\allowbreak store,\allowbreak \ \allowbreak llm\_\allowbreak manager,\allowbreak \ \allowbreak top\_\allowbreak k=\allowbreak 5,\allowbreak \ \allowbreak score\_\allowbreak threshold=\allowbreak None,\allowbreak \ \allowbreak source\_\allowbreak filter=\allowbreak None)\allowbreak \ \allowbreak -\textgreater{}\allowbreak \ \allowbreak Dict{[}str,\allowbreak \ \allowbreak Any{]}}
\item
  \textbf{Anti-Hallucination Guard:} Checks
  \texttt{if\ \allowbreak not\ \allowbreak retrieval\_\allowbreak response{[}"found"{]}:}. If false,
  bypasses the LLM entirely and immediately returns the fallback refusal
  with zero generation latency.
\item
  \textbf{LCEL Chain:}
  \texttt{chain\ \allowbreak =\allowbreak \ \allowbreak prompt\ \allowbreak \textbar{}\ \allowbreak llm\ \allowbreak \textbar{}\ \allowbreak StrOutputParser(\allowbreak )\allowbreak }.
  Invokes chain with formatted context and question. Measures
  \texttt{generation\_\allowbreak time} and \texttt{total\_\allowbreak rag\_\allowbreak time}. Returns
  answer, sources, and metrics dictionary.
\end{itemize}

\begin{center}\rule{0.5\linewidth}{0.5pt}\end{center}

\subsection{\texorpdfstring{6.7 Module: Retrieval Evaluation Metrics
(\texttt{src/\allowbreak evaluation})}{6.7 Module: Retrieval Evaluation Metrics (src/evaluation)}}\label{module-retrieval-evaluation-metrics-srcevaluation}

\subsubsection{\texorpdfstring{\texttt{retrieval\_\allowbreak metrics.\allowbreak py}}{retrieval\_metrics.py}}\label{retrieval_metrics.py}

Contains mathematical evaluation metrics for offline benchmarking: -
\texttt{calculate\_\allowbreak precision\_\allowbreak at\_\allowbreak k(\allowbreak retrieved\_\allowbreak sources,\allowbreak \ \allowbreak expected\_\allowbreak source,\allowbreak \ \allowbreak k)\allowbreak \ \allowbreak -\textgreater{}\allowbreak \ \allowbreak float}:
Evaluates the fraction of top-\(K\) retrieved sources that match the
ground-truth document:
\[\text{Precision@K} = \frac{\sum_{i=1}^K \mathbb{I}(\text{source}_i = \text{expected})}{K}\]
-
\texttt{calculate\_\allowbreak recall\_\allowbreak at\_\allowbreak k(\allowbreak retrieved\_\allowbreak sources,\allowbreak \ \allowbreak expected\_\allowbreak source,\allowbreak \ \allowbreak k)\allowbreak \ \allowbreak -\textgreater{}\allowbreak \ \allowbreak float}:
Evaluates if at least one relevant document was retrieved:
\[\text{Recall@K} = \min\left(1.0, \sum_{i=1}^K \mathbb{I}(\text{source}_i = \text{expected})\right)\]
-
\texttt{calculate\_\allowbreak hit\_\allowbreak rate\_\allowbreak at\_\allowbreak k(\allowbreak retrieved\_\allowbreak sources,\allowbreak \ \allowbreak expected\_\allowbreak source,\allowbreak \ \allowbreak k)\allowbreak \ \allowbreak -\textgreater{}\allowbreak \ \allowbreak int}:
Returns \(1\) if expected source is present in top-\(K\), else \(0\).

\begin{center}\rule{0.5\linewidth}{0.5pt}\end{center}

\subsection{\texorpdfstring{6.8 Module: Utilities, Logging, and
Persistence
(\texttt{src/\allowbreak utils})}{6.8 Module: Utilities, Logging, and Persistence (src/utils)}}\label{module-utilities-logging-and-persistence-srcutils}

\subsubsection{\texorpdfstring{1.
\texttt{file\_\allowbreak utils.\allowbreak py}}{1. file\_utils.py}}\label{file_utils.py}

\begin{itemize}
\tightlist
\item
  \texttt{calculate\_\allowbreak file\_\allowbreak hash(\allowbreak file\_\allowbreak content:\ \allowbreak Union{[}bytes,\allowbreak \ \allowbreak Path,\allowbreak \ \allowbreak str{]})\allowbreak \ \allowbreak -\textgreater{}\allowbreak \ \allowbreak str}:
  Streams file in 4096-byte blocks through \texttt{hashlib.\allowbreak sha256(\allowbreak )\allowbreak }.
  Generates deterministic SHA-256 hex digest for document deduplication.
\end{itemize}

\subsubsection{\texorpdfstring{2.
\texttt{history\_\allowbreak manager.\allowbreak py}}{2. history\_manager.py}}\label{history_manager.py}

\begin{itemize}
\tightlist
\item
  \texttt{HISTORY\_\allowbreak FILE\ \allowbreak =\allowbreak \ \allowbreak Path(\allowbreak "data/\allowbreak history.\allowbreak json")\allowbreak }
\item
  \texttt{load\_\allowbreak history(\allowbreak )\allowbreak \ \allowbreak -\textgreater{}\allowbreak \ \allowbreak list}: Loads JSON records
  from disk.
\item
  \texttt{save\_\allowbreak history(\allowbreak items)\allowbreak \ \allowbreak -\textgreater{}\allowbreak \ \allowbreak bool}: Writes formatted
  JSON atomically with \texttt{indent=\allowbreak 2}.
\item
  \texttt{add\_\allowbreak history\_\allowbreak entry(\allowbreak item\_\allowbreak type,\allowbreak \ \allowbreak title,\allowbreak \ \allowbreak details,\allowbreak \ \allowbreak source)\allowbreak }:
  Inserts new entry at index 0 (newest first) with generated
  \texttt{uuid.\allowbreak uuid4(\allowbreak )\allowbreak } and formatted timestamp string.
\item
  \texttt{delete\_\allowbreak history\_\allowbreak entry(\allowbreak entry\_\allowbreak id)\allowbreak }: Removes specified record
  and saves file.
\item
  \texttt{clear\_\allowbreak all\_\allowbreak history(\allowbreak )\allowbreak }: Wipes all records.
\end{itemize}

\subsubsection{\texorpdfstring{3.
\texttt{logger.\allowbreak py}}{3. logger.py}}\label{logger.py}

\begin{itemize}
\tightlist
\item
  \texttt{setup\_\allowbreak logger(\allowbreak name=\allowbreak "studyrag")\allowbreak \ \allowbreak -\textgreater{}\allowbreak \ \allowbreak logging.\allowbreak Logger}:
  Configures structured \texttt{StreamHandler(\allowbreak sys.\allowbreak stdout)\allowbreak } with
  timestamp formatting
  (\texttt{\%(\allowbreak asctime)\allowbreak s\ \allowbreak -\ \allowbreak \%(\allowbreak name)\allowbreak s\ \allowbreak -\ \allowbreak \%(\allowbreak levelname)\allowbreak s\ \allowbreak -\ \allowbreak \%(\allowbreak message)\allowbreak s}).
  Prevents handler duplication across Streamlit reruns.
\end{itemize}

\begin{center}\rule{0.5\linewidth}{0.5pt}\end{center}

\subsection{\texorpdfstring{6.9 Module: Interactive Presentation Layer
(\texttt{app.\allowbreak py})}{6.9 Module: Interactive Presentation Layer (app.py)}}\label{module-interactive-presentation-layer-app.py}

\begin{itemize}
\tightlist
\item
  \textbf{State Management:} Manages Streamlit session state keys:

  \begin{itemize}
  \tightlist
  \item
    \texttt{st.\allowbreak session\_\allowbreak state.\allowbreak processed\_\allowbreak documents}: Maps file hashes to
    chunk lists.
  \item
    \texttt{st.\allowbreak session\_\allowbreak state.\allowbreak vector\_\allowbreak store\_\allowbreak manager}: Cached vector
    store instance.
  \item
    \texttt{st.\allowbreak session\_\allowbreak state.\allowbreak document\_\allowbreak sources}: Dynamic list of
    uploaded documents.
  \item
    \texttt{st.\allowbreak session\_\allowbreak state.\allowbreak chat\_\allowbreak messages}: Q\&A chat message history
    with sources.
  \item
    \texttt{st.\allowbreak session\_\allowbreak state.\allowbreak app\_\allowbreak theme}: Active UI theme string.
  \item
    \texttt{st.\allowbreak session\_\allowbreak state.\allowbreak background\_\allowbreak jobs}: Active and completed
    worker jobs.
  \end{itemize}
\item
  \textbf{Navigation Router:} Routes between 6 full views:

  \begin{enumerate}
  \def\labelenumi{\arabic{enumi}.}
  \tightlist
  \item
    \texttt{{[}Study\ \allowbreak Library{]}\ \allowbreak \ \allowbreak Study\ \allowbreak Library}: Upload, chunking,
    indexing, and deletion of study materials.
  \item
    \texttt{{[}Ask\ \allowbreak Questions{]}\ \allowbreak \ \allowbreak Ask\ \allowbreak Questions}: Interactive chat
    with page-level citation expanders.
  \item
    \texttt{{[}Practice\ \allowbreak Quiz{]}\ \allowbreak \ \allowbreak Practice\ \allowbreak Quiz}: Interactive MCQ
    testing with automated grading and balloons.
  \item
    \texttt{{[}Flashcards{]}\ \allowbreak \ \allowbreak Flashcards}: Expandable study cards for
    definition and term memorization.
  \item
    \texttt{{[}Topic\ \allowbreak Summaries{]}\ \allowbreak \ \allowbreak Topic\ \allowbreak Summaries}: High-level
    chapter overviews and key takeaway generation.
  \item
    \texttt{{[}Study\ \allowbreak History{]}\ \allowbreak \ \allowbreak Study\ \allowbreak History}: Filterable audit
    log with item-by-item deletion.
  \end{enumerate}
\end{itemize}

\section{Chapter 7: Methodology and Development
Approach}\label{chapter-7-methodology-and-development-approach}

\subsection{7.1 Development Lifecycle and Implementation
Phases}\label{development-lifecycle-and-implementation-phases}

StudyRAG was developed using an \textbf{Iterative, Test-Driven Modular
Engineering Methodology} divided into distinct development phases:

\begin{verbatim}
[Phase 1: Ingestion & Text Normalization]
      |  (Extract text page-by-page via PyMuPDF, normalize whitespace, regex cleaning)
      v
[Phase 2: Semantic Chunking & Vectorization]
      |  (Recursive character splitting, SentenceTransformers 384d dense embeddings)
      v
[Phase 3: Dual Vector Store Indexing]
      |  (FAISS in-memory L2 index & ChromaDB persistent SQLite cosine collections)
      v
[Phase 4: Retrieval Engine & Relevance Scoring]
      |  (Query validation, metric normalization, threshold filtering, timing telemetry)
      v
[Phase 5: LLM Orchestration & Prompt Engineering]
      |  (Groq, OpenAI, Ollama multi-provider factory, strict JSON prompt templates)
      v
[Phase 6: Multi-Format Study Tool Generation]
      |  (Practice Quiz MCQ generator, Flashcard decks, Chapter revision summaries)
      v
[Phase 7: Presentation, Background Threading & History Persistence]
      |  (Streamlit web app, async daemon threads, Dark/Light modes, JSON audit logs)
      v
[Phase 8: Automated Verification & Documentation]
      |  (28 pytest unit/integration tests, test matrix, comprehensive system handover)
\end{verbatim}

\subsection{7.2 Requirements Gathering and Pedagogical
Design}\label{requirements-gathering-and-pedagogical-design}

The system was designed from the perspective of \textbf{Cognitive
Learning Science}: 1. \textbf{Active Recall}: The system discourages
passive reading by converting static PDFs into dynamic testing tools. 2.
\textbf{Pedagogical Alignment}: Study materials uploaded by the student
form the sole source of truth, guaranteeing alignment with classroom
lectures and exams. 3. \textbf{Citation Verifiability}: Every answer
references exact pages so students can cross-reference textbook figures,
diagrams, and proof steps.

\subsection{7.3 Architectural Decisions and Trade-Off
Analysis}\label{architectural-decisions-and-trade-off-analysis}

\begingroup\footnotesize
\begin{longtable}[]{@{}
  >{\raggedright\arraybackslash\sloppy}p{(\linewidth - 6\tabcolsep) * \real{0.1800}}
  >{\raggedright\arraybackslash\sloppy}p{(\linewidth - 6\tabcolsep) * \real{0.2300}}
  >{\raggedright\arraybackslash\sloppy}p{(\linewidth - 6\tabcolsep) * \real{0.1900}}
  >{\raggedright\arraybackslash\sloppy}p{(\linewidth - 6\tabcolsep) * \real{0.3600}}@{}}
\toprule\noalign{}
\begin{minipage}[b]{\linewidth}\raggedright
Architectural Decision
\end{minipage} & \begin{minipage}[b]{\linewidth}\raggedright
Option Selected
\end{minipage} & \begin{minipage}[b]{\linewidth}\raggedright
Alternatives Considered
\end{minipage} & \begin{minipage}[b]{\linewidth}\raggedright
Engineering Rationale \& Trade-Off Analysis
\end{minipage} \\
\midrule\noalign{}
\endhead
\bottomrule\noalign{}
\endlastfoot
\textbf{PDF Extraction Engine} & PyMuPDF (\texttt{fitz}) &
\texttt{pypdf}, \texttt{pdfplumber}, \texttt{PyPDF2} & \textbf{Chosen:}
Written in C/C++; extracts complex pages 10x-20x faster than pure Python
while tracking exact page numbers. \newline **Trade-off:** Requires
binary wheel dependencies on host platform. \\
\textbf{Chunking Strategy} & Recursive Character Splitting
(\texttt{size=\allowbreak 500,\allowbreak \ \allowbreak overlap=\allowbreak 50}) & Fixed character window, Semantic
splitting & \textbf{Chosen:} Tries paragraph breaks
(\texttt{\textbackslash{}n\textbackslash{}n}), then newlines
(\texttt{\textbackslash{}n}), then sentences (\texttt{.}) before
resorting to spaces. Keeps semantic units intact.
\newline **Trade-off:** Chunks vary slightly in character length. \\
\textbf{Embedding Model} & \texttt{all-MiniLM-L6-v2} & OpenAI
\texttt{text-embedding-3-small}, Cohere Embed & \textbf{Chosen:} Free,
open-source, runs locally on CPU in \(<50\,\text{ms}\), 384 dimensions,
zero data privacy leaks. \newline **Trade-off:** English-centric;
slightly lower capacity than massive 1536d models. \\
\textbf{Default Vector Store} & FAISS CPU (\texttt{IndexFlatL2}) &
ChromaDB, Qdrant, Pinecone & \textbf{Chosen:} Extremely lightweight,
requires no external database server, zero network latency, serializes
cleanly to disk. \newline **Trade-off:** In-memory index recreated or
deserialized upon startup. \\
\textbf{LLM Inference Provider} & Multi-Provider (Default: Groq Cloud) &
OpenAI Only, Local Only & \textbf{Chosen:} Groq provides ultra-fast
inference (500+ tokens/sec) for quizzes/flashcards, with Ollama
providing 100\% offline fallback and OpenAI for enterprise needs.
\newline **Trade-off:** Requires API key management for cloud modes. \\
\textbf{Concurrency Model} & Python \texttt{threading.\allowbreak Thread} +
Streamlit context & Celery + Redis, Ray & \textbf{Chosen:} Zero external
infrastructure required; lightweight daemon threads run inside the same
process without Redis overhead. \newline **Trade-off:** Jobs do not
survive hard application server restarts. \\
\end{longtable}
\endgroup

\section{Chapter 8: Algorithms, Formulas, and Core
Logic}\label{chapter-8-algorithms-formulas-and-core-logic}

\subsection{8.1 SHA-256 Content Deduplication
Algorithm}\label{sha-256-content-deduplication-algorithm}

To prevent redundant processing and vector index bloat when users upload
duplicate or already-indexed documents, StudyRAG computes the
cryptographic SHA-256 hash across the raw byte stream of each uploaded
file:

\[\text{Hash} = \text{SHA-256}(B) = \sum_{i=0}^{N-1} \text{Compression}(H_{i-1}, M_i)\]

Where the file is processed in 4,096-byte blocks (\(M_i\)). If the
resulting 64-character hexadecimal digest matches a key in
\texttt{st.\allowbreak session\_\allowbreak state.\allowbreak processed\_\allowbreak documents}, re-indexing is
completely bypassed.

\subsection{8.2 Recursive Character Text Splitting
Algorithm}\label{recursive-character-text-splitting-algorithm}

The recursive chunking algorithm implements a hierarchical separator
cascade:
\[S = [\texttt{"\textbackslash n\textbackslash n"}, \texttt{"\textbackslash n"}, \texttt{". "}, \texttt{" "}, \texttt{""}]\]

Given a page text \(T\), the algorithm executes as follows: 1. Tests the
primary separator \(S_0 = \texttt{"\textbackslash n\textbackslash n"}\).
If splitting by \(S_0\) produces segments \(\le \text{chunk\_size}\)
(500 chars), the segments are retained. 2. If any segment
\(> \text{chunk\_size}\), the algorithm recursively falls back to
\(S_1 = \texttt{"\textbackslash n"}\), then \(S_2 = \text{". "}\), and
finally single spaces. 3. Adjacent segments are merged with a sliding
window that preserves \(\text{chunk\_overlap}\) (50 characters),
guaranteeing that sentences spanning chunk boundaries maintain syntactic
continuity.

\subsection{8.3 Vector Similarity Mathematics and Metric
Normalization}\label{vector-similarity-mathematics-and-metric-normalization}

Vector databases employ different mathematical metrics to calculate the
proximity between a query vector \(\mathbf{q} \in \mathbb{R}^D\) and a
document chunk vector \(\mathbf{d}_i \in \mathbb{R}^D\) (\(D = 384\)).

\subsubsection{1. FAISS Euclidean (L2)
Distance}\label{faiss-euclidean-l2-distance}

FAISS computes the squared Euclidean distance:
\[d_{\text{L2}}(\mathbf{q}, \mathbf{d}_i) = \\|\mathbf{q} - \mathbf{d}_i\\|_2 = \sqrt{\sum_{j=1}^D (q_j - d_{i,j})^2}\]
Where \(d_{\text{L2}} \in [0, \infty)\). A distance of \(0\) indicates
identity, while larger values indicate dissimilarity.

\textbf{StudyRAG Normalization Formula
(\texttt{src/\allowbreak retrieval/\allowbreak score\_\allowbreak filter.\allowbreak py}):} To map the unbounded L2
distance into a normalized relevance score
\(S_{\text{FAISS}} \in (0.0, 1.0]\):
\[S_{\text{FAISS}} = \frac{1.0}{1.0 + d_{\text{L2}}(\mathbf{q}, \mathbf{d}_i)}\]
- If \(d_{\text{L2}} = 0.0\) (identical vectors)
\(\implies S_{\text{FAISS}} = 1.0\). - If
\(d_{\text{L2}} = 1.0 \implies S_{\text{FAISS}} = 0.5\). - As
\(d_{\text{L2}} \to \infty \implies S_{\text{FAISS}} \to 0.0\).

\subsubsection{2. ChromaDB Cosine
Distance}\label{chromadb-cosine-distance}

ChromaDB collections in StudyRAG are initialized with
\texttt{\{"hnsw:space":\ \allowbreak "cosine"\}}. The cosine distance is defined as:
\[d_{\text{cosine}}(\mathbf{q}, \mathbf{d}_i) = 1.0 - \cos(\theta) = 1.0 - \frac{\mathbf{q} \cdot \mathbf{d}_i}{\\|\mathbf{q}\\|_2 \\|\mathbf{d}_i\\|_2}\]
Where \(d_{\text{cosine}} \in [0.0, 2.0]\).

\textbf{StudyRAG Normalization Formula
(\texttt{src/\allowbreak retrieval/\allowbreak score\_\allowbreak filter.\allowbreak py}):}
\[S_{\text{Chroma}} = \max\left(0.0, \min\left(1.0, 1.0 - d_{\text{cosine}}\right)\right)\]
- If \(d_{\text{cosine}} = 0.0\) (collinear vectors)
\(\implies S_{\text{Chroma}} = 1.0\). - If
\(d_{\text{cosine}} = 0.2 \implies S_{\text{Chroma}} = 0.8\). - If
\(d_{\text{cosine}} \ge 1.0\) (orthogonal/opposite vectors)
\(\implies S_{\text{Chroma}} = 0.0\).

\subsection{8.4 Relevance Threshold Filtering
Logic}\label{relevance-threshold-filtering-logic}

Given normalized scores \(S_i \in [0, 1]\) and a user-configured
similarity threshold \(\tau\) (default \(\tau = 0.20\)):
\[\text{FilteredResults} = \\{ (\mathbf{d}_i, S_i) \mid S_i \ge \tau \\}\]

The filtered set is sorted in descending order of relevance:
\[\text{SortedResults} = \operatorname{sort}_{\downarrow}( \text{FilteredResults}, \text{key}=S_i )\]

If \(|\text{FilteredResults}| = 0\), the retrieval pipeline sets
\texttt{found\ \allowbreak =\allowbreak \ \allowbreak False}, activating the anti-hallucination fallback.

\subsection{8.5 Offline Retrieval Evaluation
Metrics}\label{offline-retrieval-evaluation-metrics}

The evaluation module (\texttt{src/\allowbreak evaluation/\allowbreak retrieval\_\allowbreak metrics.\allowbreak py})
provides quantitative metrics for measuring information retrieval
performance against gold-standard evaluation datasets:

\subsubsection{1. Precision@K}\label{precisionk}

Measures the proportion of top-\(K\) retrieved document chunks that
originate from the expected source document:
\[\text{Precision@K} = \frac{\sum_{i=1}^K \mathbb{I}(\text{source}_i = \text{expected\_source})}{K}\]

\subsubsection{2. Recall@K}\label{recallk}

Measures whether the relevant document was successfully recovered in the
top-\(K\) candidates:
\[\text{Recall@K} = \min\left(1.0, \sum_{i=1}^K \mathbb{I}(\text{source}_i = \text{expected\_source})\right)\]

\subsubsection{3. Hit Rate@K}\label{hit-ratek}

Binary metric indicating whether at least one chunk in top-\(K\) matches
the target source:
\[\text{Hit Rate@K} = \begin{cases} 1 & \text{if } \exists i \in [1, K] \text{ such that } \text{source}_i = \text{expected\_source} \\ 0 & \text{otherwise} \end{cases}\]

\section{Chapter 9: Database and Data
Management}\label{chapter-9-database-and-data-management}

\subsection{9.1 Persistent Storage
Architecture}\label{persistent-storage-architecture}

StudyRAG operates a hybrid data management strategy comprising file
storage, vector serialization, and structured JSON record storage:

\begin{verbatim}
data/
    uploaded/                   # Physical PDF storage
        os.pdf                  # Preserved for re-indexing & inspection
    vectorstores/               # Vector database storage
        faiss/
            index.faiss         # Binary serialized FAISS L2 index
            index.pkl           # Pickled LangChain Document & ID mapping
        chroma/
            chroma.sqlite3      # Embedded SQLite database for collections & metadata
            <uuid>/             # Segment directory containing binary HNSW graphs
    history.json                # Structured JSON audit log of all study artifacts
\end{verbatim}

\subsection{\texorpdfstring{9.2 Complete Schema: Study History
(\texttt{data/\allowbreak history.\allowbreak json})}{9.2 Complete Schema: Study History (data/history.json)}}\label{complete-schema-study-history-datahistory.json}

The study history file stores an array of structured JSON objects. Every
record conforms to the following schema:

\begin{Shaded}
\begin{Highlighting}[]
\FunctionTok{\{}
  \DataTypeTok{"id"}\FunctionTok{:} \StringTok{"b203dc1f{-}7fd9{-}468a{-}aa2d{-}31169f359203"}\FunctionTok{,}
  \DataTypeTok{"type"}\FunctionTok{:} \StringTok{"quiz"}\FunctionTok{,}
  \DataTypeTok{"title"}\FunctionTok{:} \StringTok{"Practice Quiz: Operating System Deadlocks"}\FunctionTok{,}
  \DataTypeTok{"source"}\FunctionTok{:} \StringTok{"os.pdf"}\FunctionTok{,}
  \DataTypeTok{"timestamp"}\FunctionTok{:} \StringTok{"Oct 14, 2026 {-} 11:45 AM"}\FunctionTok{,}
  \DataTypeTok{"details"}\FunctionTok{:} \FunctionTok{\{}
    \DataTypeTok{"topic"}\FunctionTok{:} \StringTok{"Operating System Deadlocks"}\FunctionTok{,}
    \DataTypeTok{"num\_questions"}\FunctionTok{:} \DecValTok{5}\FunctionTok{,}
    \DataTypeTok{"questions"}\FunctionTok{:} \OtherTok{[}
      \FunctionTok{\{}
        \DataTypeTok{"question"}\FunctionTok{:} \StringTok{"Which condition is NOT required for a deadlock to occur?"}\FunctionTok{,}
        \DataTypeTok{"options"}\FunctionTok{:} \OtherTok{[}
          \StringTok{"A. Mutual Exclusion"}\OtherTok{,}
          \StringTok{"B. Hold and Wait"}\OtherTok{,}
          \StringTok{"C. Preemption Allowed"}\OtherTok{,}
          \StringTok{"D. Circular Wait"}
        \OtherTok{]}\FunctionTok{,}
        \DataTypeTok{"answer"}\FunctionTok{:} \StringTok{"C. Preemption Allowed"}
      \FunctionTok{\}}
    \OtherTok{]}
  \FunctionTok{\}}
\FunctionTok{\}}
\end{Highlighting}
\end{Shaded}

\subsubsection{Field Specifications:}\label{field-specifications}

\begin{itemize}
\tightlist
\item
  \texttt{id} (string, required): RFC 4122 compliant UUID v4 string.
\item
  \texttt{type} (string, required): One of \texttt{"quiz"},
  \texttt{"flashcard"}, \texttt{"summary"}, \texttt{"chat"},
  \texttt{"material"}.
\item
  \texttt{title} (string, required): Human-readable descriptive header.
\item
  \texttt{source} (string, required): Target document filename
  (e.g.~\texttt{"os.\allowbreak pdf"}) or \texttt{"All\ \allowbreak Documents"}.
\item
  \texttt{timestamp} (string, required): Formatted timestamp string
  (\texttt{"\%b\ \allowbreak \%d,\allowbreak \ \allowbreak \%Y\ \allowbreak -\ \allowbreak \%I:\%M\ \allowbreak \%p"}).
\item
  \texttt{details} (object, required): Payload schema varies by
  \texttt{type}:

  \begin{itemize}
  \tightlist
  \item
    For \texttt{"quiz"}: contains \texttt{topic},
    \texttt{num\_\allowbreak questions}, and \texttt{questions} array.
  \item
    For \texttt{"flashcard"}: contains \texttt{topic},
    \texttt{num\_\allowbreak cards}, and \texttt{cards} array
    (\texttt{{[}\{"front":\ \allowbreak ".\allowbreak .\allowbreak .",\allowbreak \ \allowbreak "back":\ \allowbreak ".\allowbreak .\allowbreak ."\}{]}}).
  \item
    For \texttt{"summary"}: contains \texttt{topic} and
    \texttt{summary\_\allowbreak text} (Markdown string).
  \item
    For \texttt{"chat"}: contains \texttt{query}, \texttt{answer}, and
    \texttt{sources} list.
  \item
    For \texttt{"material"}: contains \texttt{doc\_\allowbreak name} and
    \texttt{action} (\texttt{"upload"} or \texttt{"delete"}).
  \end{itemize}
\end{itemize}

\subsection{9.3 Entity-Relationship
Diagram}\label{entity-relationship-diagram}

\begin{Shaded}
\begin{Highlighting}[]
\NormalTok{erDiagram}
\NormalTok{    UPLOADED\_DOCUMENT ||{-}{-}o\{ DOCUMENT\_CHUNK : "split into"}
\NormalTok{    DOCUMENT\_CHUNK ||{-}{-}o| EMBEDDING\_VECTOR : "represented by"}
\NormalTok{    UPLOADED\_DOCUMENT ||{-}{-}o\{ STUDY\_HISTORY\_ENTRY : "scoped to"}
    
\NormalTok{    UPLOADED\_DOCUMENT \{}
\NormalTok{        string filename PK}
\NormalTok{        string sha256\_hash}
\NormalTok{        int total\_pages}
\NormalTok{        datetime upload\_time}
\NormalTok{    \}}
    
\NormalTok{    DOCUMENT\_CHUNK \{}
\NormalTok{        int chunk\_id PK}
\NormalTok{        string source\_doc FK}
\NormalTok{        int page\_number}
\NormalTok{        string text\_content}
\NormalTok{    \}}
    
\NormalTok{    EMBEDDING\_VECTOR \{}
\NormalTok{        int chunk\_id FK}
\NormalTok{        float32\_array vector\_384d}
\NormalTok{        string store\_backend}
\NormalTok{    \}}
    
\NormalTok{    STUDY\_HISTORY\_ENTRY \{}
\NormalTok{        string uuid PK}
\NormalTok{        string entry\_type}
\NormalTok{        string title}
\NormalTok{        string source\_filter FK}
\NormalTok{        string timestamp}
\NormalTok{        json payload\_details}
\NormalTok{    \}}
\end{Highlighting}
\end{Shaded}

\subsection{9.4 CRUD Operations \& Data
Lifecycle}\label{crud-operations-data-lifecycle}

\begin{itemize}
\tightlist
\item
  \textbf{Create}:

  \begin{itemize}
  \tightlist
  \item
    Uploading a PDF writes the binary file to \texttt{data/\allowbreak uploaded/\allowbreak },
    splits chunks, computes embeddings, saves the vector store, and
    writes an \texttt{"upload"} event to \texttt{data/\allowbreak history.\allowbreak json}.
  \item
    Generating quizzes, flashcards, or summaries appends a record at
    index 0 in \texttt{data/\allowbreak history.\allowbreak json}.
  \end{itemize}
\item
  \textbf{Read}:

  \begin{itemize}
  \tightlist
  \item
    \texttt{load\_\allowbreak history(\allowbreak )\allowbreak } reads and parses
    \texttt{data/\allowbreak history.\allowbreak json}.
  \item
    \texttt{load(\allowbreak )\allowbreak } in \texttt{FAISSStore} deserializes
    \texttt{index.\allowbreak faiss}.
  \end{itemize}
\item
  \textbf{Delete}:

  \begin{itemize}
  \tightlist
  \item
    Deleting a document removes the physical file from
    \texttt{data/\allowbreak uploaded/\allowbreak }, drops matching chunks from memory, resets
    the vector index, rebuilds the index from remaining documents, and
    records a \texttt{"delete"} event in history.
  \item
    Deleting a history entry calls
    \texttt{delete\_\allowbreak history\_\allowbreak entry(\allowbreak entry\_\allowbreak id)\allowbreak }, filtering the array and
    overwriting \texttt{history.\allowbreak json}.
  \item
    \texttt{clear\_\allowbreak all\_\allowbreak history(\allowbreak )\allowbreak } atomically overwrites
    \texttt{history.\allowbreak json} with \texttt{{[}{]}}.
  \end{itemize}
\end{itemize}

\section{Chapter 10: API and Integration
Documentation}\label{chapter-10-api-and-integration-documentation}

\subsection{10.1 Internal Component API
Contracts}\label{internal-component-api-contracts}

\subsubsection{Ingestion Pipeline API}\label{ingestion-pipeline-api}

\begin{itemize}
\tightlist
\item
  \texttt{process\_\allowbreak pdf(\allowbreak pdf\_\allowbreak path:\ \allowbreak str\ \allowbreak \textbar{}\ \allowbreak Path,\allowbreak \ \allowbreak chunk\_\allowbreak size:\ \allowbreak int\ \allowbreak =\allowbreak \ \allowbreak 500,\allowbreak \ \allowbreak chunk\_\allowbreak overlap:\ \allowbreak int\ \allowbreak =\allowbreak \ \allowbreak 50)\allowbreak \ \allowbreak -\textgreater{}\allowbreak \ \allowbreak List{[}Dict{[}str,\allowbreak \ \allowbreak Any{]}{]}}

  \begin{itemize}
  \tightlist
  \item
    Input: Valid path to \texttt{.\allowbreak pdf} file.
  \item
    Output:
    \texttt{{[}\{"text":\ \allowbreak ".\allowbreak .\allowbreak .",\allowbreak \ \allowbreak "source":\ \allowbreak "os.\allowbreak pdf",\allowbreak \ \allowbreak "page":\ \allowbreak 1,\allowbreak \ \allowbreak "chunk\_\allowbreak id":\ \allowbreak 0\},\allowbreak \ \allowbreak .\allowbreak .\allowbreak .{]}}
  \item
    Exceptions: \texttt{FileNotFoundError}, \texttt{ValueError} (invalid
    extension or corrupted file).
  \end{itemize}
\end{itemize}

\subsubsection{Retrieval Pipeline API}\label{retrieval-pipeline-api}

\begin{itemize}
\tightlist
\item
  \texttt{retrieve\_\allowbreak context(\allowbreak query:\ \allowbreak str,\allowbreak \ \allowbreak vector\_\allowbreak store:\ \allowbreak VectorStoreManager,\allowbreak \ \allowbreak top\_\allowbreak k:\ \allowbreak int\ \allowbreak =\allowbreak \ \allowbreak 5,\allowbreak \ \allowbreak score\_\allowbreak threshold:\ \allowbreak Optional{[}float{]}\ \allowbreak =\allowbreak \ \allowbreak None,\allowbreak \ \allowbreak source\_\allowbreak filter:\ \allowbreak Optional{[}str{]}\ \allowbreak =\allowbreak \ \allowbreak None)\allowbreak \ \allowbreak -\textgreater{}\allowbreak \ \allowbreak Dict{[}str,\allowbreak \ \allowbreak Any{]}}

  \begin{itemize}
  \item
    Output Schema:

\begin{Shaded}
\begin{Highlighting}[]
\FunctionTok{\{}
  \DataTypeTok{"query"}\FunctionTok{:} \StringTok{"string (cleaned)"}\FunctionTok{,}
  \DataTypeTok{"found"}\FunctionTok{:} \KeywordTok{true}\FunctionTok{,}
  \DataTypeTok{"results"}\FunctionTok{:} \OtherTok{[}
    \FunctionTok{\{}
      \DataTypeTok{"text"}\FunctionTok{:} \StringTok{"chunk content"}\FunctionTok{,}
      \DataTypeTok{"source"}\FunctionTok{:} \StringTok{"os.pdf"}\FunctionTok{,}
      \DataTypeTok{"page"}\FunctionTok{:} \DecValTok{4}\FunctionTok{,}
      \DataTypeTok{"chunk\_id"}\FunctionTok{:} \DecValTok{12}\FunctionTok{,}
      \DataTypeTok{"score"}\FunctionTok{:} \FloatTok{0.8421}
    \FunctionTok{\}}
  \OtherTok{]}\FunctionTok{,}
  \DataTypeTok{"message"}\FunctionTok{:} \StringTok{"Found 1 relevant results."}\FunctionTok{,}
  \DataTypeTok{"metrics"}\FunctionTok{:} \FunctionTok{\{}
    \DataTypeTok{"search\_time"}\FunctionTok{:} \FloatTok{0.012}\FunctionTok{,}
    \DataTypeTok{"filter\_time"}\FunctionTok{:} \FloatTok{0.001}\FunctionTok{,}
    \DataTypeTok{"total\_time"}\FunctionTok{:} \FloatTok{0.013}
  \FunctionTok{\}}
\FunctionTok{\}}
\end{Highlighting}
\end{Shaded}
  \end{itemize}
\end{itemize}

\subsubsection{RAG Pipeline API}\label{rag-pipeline-api}

\begin{itemize}
\tightlist
\item
  \texttt{generate\_\allowbreak rag\_\allowbreak response(\allowbreak query:\ \allowbreak str,\allowbreak \ \allowbreak vector\_\allowbreak store:\ \allowbreak VectorStoreManager,\allowbreak \ \allowbreak llm\_\allowbreak manager:\ \allowbreak LLMManager,\allowbreak \ \allowbreak top\_\allowbreak k:\ \allowbreak int\ \allowbreak =\allowbreak \ \allowbreak 5,\allowbreak \ \allowbreak score\_\allowbreak threshold:\ \allowbreak Optional{[}float{]}\ \allowbreak =\allowbreak \ \allowbreak None,\allowbreak \ \allowbreak source\_\allowbreak filter:\ \allowbreak Optional{[}str{]}\ \allowbreak =\allowbreak \ \allowbreak None)\allowbreak \ \allowbreak -\textgreater{}\allowbreak \ \allowbreak Dict{[}str,\allowbreak \ \allowbreak Any{]}}

  \begin{itemize}
  \item
    Output Schema:

\begin{Shaded}
\begin{Highlighting}[]
\FunctionTok{\{}
  \DataTypeTok{"answer"}\FunctionTok{:} \StringTok{"string (grounded response)"}\FunctionTok{,}
  \DataTypeTok{"sources"}\FunctionTok{:} \OtherTok{[}\FunctionTok{\{}\DataTypeTok{"text"}\FunctionTok{:} \StringTok{"..."}\FunctionTok{,} \DataTypeTok{"source"}\FunctionTok{:} \StringTok{"..."}\FunctionTok{,} \DataTypeTok{"page"}\FunctionTok{:} \DecValTok{1}\FunctionTok{\}}\OtherTok{]}\FunctionTok{,}
  \DataTypeTok{"metrics"}\FunctionTok{:} \FunctionTok{\{}
    \DataTypeTok{"search\_time"}\FunctionTok{:} \FloatTok{0.012}\FunctionTok{,}
    \DataTypeTok{"filter\_time"}\FunctionTok{:} \FloatTok{0.001}\FunctionTok{,}
    \DataTypeTok{"total\_time"}\FunctionTok{:} \FloatTok{0.013}\FunctionTok{,}
    \DataTypeTok{"generation\_time"}\FunctionTok{:} \FloatTok{0.450}\FunctionTok{,}
    \DataTypeTok{"total\_rag\_time"}\FunctionTok{:} \FloatTok{0.463}
  \FunctionTok{\}}
\FunctionTok{\}}
\end{Highlighting}
\end{Shaded}
  \end{itemize}
\end{itemize}

\subsection{10.2 External Service
Integrations}\label{external-service-integrations}

\subsubsection{1. Groq Cloud API}\label{groq-cloud-api}

\begin{itemize}
\tightlist
\item
  \textbf{Endpoint:}
  \texttt{https:/\allowbreak /\allowbreak api.\allowbreak groq.\allowbreak com/\allowbreak openai/\allowbreak v1/\allowbreak chat/\allowbreak completions}
\item
  \textbf{Protocol:} HTTPS REST (OpenAI-compatible)
\item
  \textbf{Authentication:}
  \texttt{Authorization:\ \allowbreak Bearer\ \allowbreak \$GROQ\_\allowbreak API\_\allowbreak KEY}
\item
  \textbf{Supported Models:}

  \begin{itemize}
  \tightlist
  \item
    \texttt{openai/\allowbreak gpt-oss-20b} (Default fast model)
  \item
    \texttt{openai/\allowbreak gpt-oss-120b} (High-capacity reasoning)
  \item
    \texttt{qwen/\allowbreak qwen3.8-27b}
  \end{itemize}
\item
  \textbf{Parameters:} \texttt{temperature:\ \allowbreak 0.0} (deterministic),
  \texttt{stream:\ \allowbreak false}.
\end{itemize}

\subsubsection{2. OpenAI API}\label{openai-api}

\begin{itemize}
\tightlist
\item
  \textbf{Endpoint:} \texttt{https:/\allowbreak /\allowbreak api.\allowbreak openai.\allowbreak com/\allowbreak v1/\allowbreak chat/\allowbreak completions}
\item
  \textbf{Protocol:} HTTPS REST
\item
  \textbf{Authentication:}
  \texttt{Authorization:\ \allowbreak Bearer\ \allowbreak \$OPENAI\_\allowbreak API\_\allowbreak KEY}
\item
  \textbf{Supported Models:} \texttt{gpt-4o}, \texttt{gpt-3.5-turbo}.
\end{itemize}

\subsubsection{3. Local Ollama Server}\label{local-ollama-server}

\begin{itemize}
\tightlist
\item
  \textbf{Endpoint:} \texttt{http:/\allowbreak /\allowbreak localhost:11434/\allowbreak api/\allowbreak chat}
\item
  \textbf{Protocol:} HTTP REST
\item
  \textbf{Authentication:} None (Local loopback)
\item
  \textbf{Supported Models:} \texttt{llama3}, \texttt{mistral},
  \texttt{gemma}.
\item
  \textbf{Health Check:}
  \texttt{curl\ \allowbreak http:/\allowbreak /\allowbreak localhost:11434/\allowbreak api/\allowbreak version}
\end{itemize}

\section{Chapter 11: User Interface and User
Workflows}\label{chapter-11-user-interface-and-user-workflows}

\subsection{11.1 Implemented Screens and
Views}\label{implemented-screens-and-views}

The Streamlit interface is structured into six dedicated views
accessible via the sidebar navigation radio:

\begin{Shaded}
\begin{Highlighting}[]
\NormalTok{graph LR}
\NormalTok{    NAV[Sidebar Navigation] {-}{-}\textgreater{} V1[[Study Library]  Study Library]}
\NormalTok{    NAV {-}{-}\textgreater{} V2[[Ask Questions]  Ask Questions]}
\NormalTok{    NAV {-}{-}\textgreater{} V3[[Practice Quiz]  Practice Quiz]}
\NormalTok{    NAV {-}{-}\textgreater{} V4[[Flashcards]  Flashcards]}
\NormalTok{    NAV {-}{-}\textgreater{} V5[[Topic Summaries]  Topic Summaries]}
\NormalTok{    NAV {-}{-}\textgreater{} V6[[Study History]  Study History]}
\end{Highlighting}
\end{Shaded}

\subsubsection{\texorpdfstring{1.
\texttt{{[}Study\ \allowbreak Library{]}\ \allowbreak \ \allowbreak Study\ \allowbreak Library}}{1. {[}Study Library{]}  Study Library}}\label{study-library-study-library}

\begin{itemize}
\tightlist
\item
  \textbf{Purpose:} Document ingestion, indexing, and library status
  tracking.
\item
  \textbf{Controls:} Multi-file drag-and-drop
  \texttt{st.\allowbreak file\_\allowbreak uploader}, status placeholders, active materials list
  with individual {[}Delete{]} Delete buttons.
\item
  \textbf{Feedback:} Success banners upon chunking and vector index
  serialization.
\end{itemize}

\subsubsection{\texorpdfstring{2.
\texttt{{[}Ask\ \allowbreak Questions{]}\ \allowbreak \ \allowbreak Ask\ \allowbreak Questions}}{2. {[}Ask Questions{]}  Ask Questions}}\label{ask-questions-ask-questions}

\begin{itemize}
\tightlist
\item
  \textbf{Purpose:} Grounded interactive Q\&A.
\item
  \textbf{Controls:} Document focus selectbox (\texttt{st.\allowbreak selectbox}),
  clear chat button, chat input field (\texttt{st.\allowbreak chat\_\allowbreak input}).
\item
  \textbf{Features:} Assistant messages display expandable citations
  (\texttt{st.\allowbreak expander(\allowbreak "\ \allowbreak View\ \allowbreak Referenced\ \allowbreak Sources")\allowbreak }) showing exact
  source filename, page number, and raw matching text snippet.
\end{itemize}

\subsubsection{\texorpdfstring{3.
\texttt{{[}Practice\ \allowbreak Quiz{]}\ \allowbreak \ \allowbreak Practice\ \allowbreak Quiz}}{3. {[}Practice Quiz{]}  Practice Quiz}}\label{practice-quiz-practice-quiz}

\begin{itemize}
\tightlist
\item
  \textbf{Purpose:} Self-assessment through active recall.
\item
  \textbf{Controls:} Topic input text box, number of questions slider
  (1-10), study source filter, ``{[}Execute{]} Generate Practice Quiz''
  button.
\item
  \textbf{Interactive Form:} Renders questions with radio buttons
  (\texttt{st.\allowbreak radio}). Clicking ``Grade My Quiz'' checks answers, marks
  correct/incorrect selections with explanations, displays the final
  score (e.g.~\texttt{{[}Success{]}\ \allowbreak \ \allowbreak Your\ \allowbreak Score:\ \allowbreak 5\ \allowbreak /\allowbreak \ \allowbreak 5}), and
  fires celebration balloons (\texttt{st.\allowbreak balloons(\allowbreak )\allowbreak }).
\end{itemize}

\subsubsection{\texorpdfstring{4.
\texttt{{[}Flashcards{]}\ \allowbreak \ \allowbreak Flashcards}}{4. {[}Flashcards{]}  Flashcards}}\label{flashcards-flashcards}

\begin{itemize}
\tightlist
\item
  \textbf{Purpose:} Term and definition memorization.
\item
  \textbf{Controls:} Topic input, number of cards (1-15), source filter,
  ``{[}Flashcards{]} Create Flashcards'' button.
\item
  \textbf{Interactive Deck:} Renders cards inside expandable sections
  with front prompts and hidden/revealed back definitions.
\end{itemize}

\subsubsection{\texorpdfstring{5.
\texttt{{[}Topic\ \allowbreak Summaries{]}\ \allowbreak \ \allowbreak Topic\ \allowbreak Summaries}}{5. {[}Topic Summaries{]}  Topic Summaries}}\label{topic-summaries-topic-summaries}

\begin{itemize}
\tightlist
\item
  \textbf{Purpose:} Fast revision guide generation.
\item
  \textbf{Controls:} Topic input, source filter, ``{[}Topic Summaries{]}
  Generate Revision Summary'' button.
\item
  \textbf{Output:} Renders structured Markdown with executive overview
  and bulleted key points.
\end{itemize}

\subsubsection{\texorpdfstring{6.
\texttt{{[}Study\ \allowbreak History{]}\ \allowbreak \ \allowbreak Study\ \allowbreak History}}{6. {[}Study History{]}  Study History}}\label{study-history-study-history}

\begin{itemize}
\tightlist
\item
  \textbf{Purpose:} Historical audit log and session persistence.
\item
  \textbf{Controls:} Activity type filter dropdown, ``{[}Delete{]} Clear
  All History'' button, individual ``{[}Delete{]} Delete Entry''
  buttons.
\item
  \textbf{Contents:} Expander cards for each past quiz, flashcard deck,
  revision summary, chat question, or library upload.
\end{itemize}

\subsection{11.2 Theme Toggle and Student-Centric
Aesthetics}\label{theme-toggle-and-student-centric-aesthetics}

\begin{itemize}
\tightlist
\item
  \textbf{Dark Mode (Default):} Deep slate \texttt{\#0F172A} background,
  \texttt{\#1E293B} cards, \texttt{\#818CF8} gradient titles,
  high-contrast readable text.
\item
  \textbf{Light Mode:} Crisp \texttt{\#F8FAFC} background, pure white
  \texttt{\#FFFFFF} cards, \texttt{\#4F46E5} indigo accents.
\item
  \textbf{Typography:} Modern Google Fonts (\texttt{Inter} for body
  text, \texttt{Outfit} for headers).
\item
  \textbf{Responsive Animations:} Hover transforms on buttons and
  interactive radio items.
\end{itemize}

\section{Chapter 12: Installation and Environment
Setup}\label{chapter-12-installation-and-environment-setup}

\subsection{12.1 Prerequisites Checklist}\label{prerequisites-checklist}

Before installing StudyRAG, ensure your host system satisfies the
following prerequisites: - \textbf{Operating System:} macOS (Apple
Silicon or Intel), Linux (Ubuntu 20.04+), or Windows 10/11. -
\textbf{Python Version:} Python 3.10, 3.11, or 3.12 (Verified on Python
3.12.2). - \textbf{Package Manager:} \texttt{pip} (Python package
manager). - \textbf{Git:} Installed and configured.

\subsection{12.2 Step-by-Step Installation
Instructions}\label{step-by-step-installation-instructions}

\subsubsection{Step 1: Clone the
Repository}\label{step-1-clone-the-repository}

\begin{Shaded}
\begin{Highlighting}[]
\FunctionTok{git}\NormalTok{ clone https://github.com/kirit{-}somani/Rag{-}Based{-}Study{-}Application{-}.git}
\BuiltInTok{cd}\NormalTok{ Rag{-}Based{-}Study{-}Application{-}}
\end{Highlighting}
\end{Shaded}

\subsubsection{Step 2: Create and Activate a Dedicated Virtual
Environment}\label{step-2-create-and-activate-a-dedicated-virtual-environment}

It is strongly recommended to use a clean virtual environment to prevent
package version conflicts:

\textbf{On macOS / Linux:}

\begin{Shaded}
\begin{Highlighting}[]
\ExtensionTok{python3} \AttributeTok{{-}m}\NormalTok{ venv venv}
\BuiltInTok{source}\NormalTok{ venv/bin/activate}
\end{Highlighting}
\end{Shaded}

\textbf{On Windows (PowerShell):}

\begin{Shaded}
\begin{Highlighting}[]
\NormalTok{python }\OperatorTok{{-}}\NormalTok{m venv venv}
\OperatorTok{.}\NormalTok{\textbackslash{}venv\textbackslash{}Scripts\textbackslash{}Activate}\OperatorTok{.}\FunctionTok{ps1}
\end{Highlighting}
\end{Shaded}

\subsubsection{Step 3: Configure Environment
Variables}\label{step-3-configure-environment-variables}

Copy the template configuration file to \texttt{.\allowbreak env}:

\begin{Shaded}
\begin{Highlighting}[]
\FunctionTok{cp}\NormalTok{ .env.example .env}
\end{Highlighting}
\end{Shaded}

Open \texttt{.\allowbreak env} in any text editor and configure your preferred LLM
provider:

\begin{Shaded}
\begin{Highlighting}[]
\NormalTok{\# LLM Provider Configuration (options: groq, openai, ollama)}
\NormalTok{LLM\_PROVIDER=groq}

\NormalTok{\# Groq Cloud API Key (Free tier at https://console.groq.com/keys)}
\NormalTok{GROQ\_API\_KEY=your\_groq\_api\_key\_here}

\NormalTok{\# OpenAI API Key (Optional, required only if LLM\_PROVIDER=openai)}
\NormalTok{OPENAI\_API\_KEY=your\_openai\_api\_key\_here}

\NormalTok{\# Ollama Configuration (Optional, required only if running local offline LLMs)}
\NormalTok{OLLAMA\_BASE\_URL=http://localhost:11434}
\NormalTok{MODEL\_NAME=llama3}

\NormalTok{\# Embedding and Vector Database Configuration}
\NormalTok{EMBEDDING\_MODEL=minilm}
\NormalTok{VECTOR\_STORE=faiss}
\end{Highlighting}
\end{Shaded}

\begin{quote}
{[}!IMPORTANT{]} Never commit your \texttt{.\allowbreak env} file to version
control. The repository includes \texttt{.\allowbreak env} in \texttt{.\allowbreak gitignore} by
default.
\end{quote}

\subsubsection{Step 4: Install Python
Dependencies}\label{step-4-install-python-dependencies}

\begin{Shaded}
\begin{Highlighting}[]
\ExtensionTok{pip}\NormalTok{ install }\AttributeTok{{-}{-}upgrade}\NormalTok{ pip}
\ExtensionTok{pip}\NormalTok{ install }\AttributeTok{{-}r}\NormalTok{ requirements.txt}
\end{Highlighting}
\end{Shaded}

\subsubsection{Step 5: Apple Silicon / macOS Special
Consideration}\label{step-5-apple-silicon-macos-special-consideration}

In Anaconda environments on macOS where TensorFlow is installed,
TensorFlow's C-extension initialization can cause a segmentation fault
when Transformers loads. To prevent this, ensure \texttt{USE\_\allowbreak TF=\allowbreak 0} and
\texttt{USE\_\allowbreak TORCH=\allowbreak 1} are set:

\begin{Shaded}
\begin{Highlighting}[]
\BuiltInTok{export} \VariableTok{USE\_TF}\OperatorTok{=}\NormalTok{0}
\BuiltInTok{export} \VariableTok{USE\_TORCH}\OperatorTok{=}\NormalTok{1}
\end{Highlighting}
\end{Shaded}

\emph{(Note: \texttt{app.\allowbreak py} sets these environment flags automatically
at the top of the file.)}

\section{Chapter 13: Complete Execution and Usage
Guide}\label{chapter-13-complete-execution-and-usage-guide}

\subsection{13.1 Launching the
Application}\label{launching-the-application}

With your virtual environment activated, launch the Streamlit server:

\begin{Shaded}
\begin{Highlighting}[]
\ExtensionTok{streamlit}\NormalTok{ run app.py}
\end{Highlighting}
\end{Shaded}

Expected terminal output:

\begin{verbatim}
  You can now view your Streamlit app in your browser.

  Local URL: http://localhost:8501
  Network URL: http://192.168.1.X:8501
\end{verbatim}

Open your web browser and navigate to \texttt{http:/\allowbreak /\allowbreak localhost:8501}.

\subsection{13.2 End-to-End Walkthrough
Tutorial}\label{end-to-end-walkthrough-tutorial}

\subsubsection{Step 1: Ingest Study
Materials}\label{step-1-ingest-study-materials}

\begin{enumerate}
\def\labelenumi{\arabic{enumi}.}
\tightlist
\item
  In the sidebar, click on \textbf{{[}Study Library{]} Study Library}.
\item
  Drag and drop a course PDF (for example, \texttt{data/\allowbreak uploaded/\allowbreak os.\allowbreak pdf}
  or your own lecture notes) into the file uploader.
\item
  Observe the progress status: \emph{'' Indexing \{filename\} for
  instant study\ldots``} followed by \emph{``{[}Success{]} \{filename\}
  successfully indexed!''}.
\item
  Your document will now appear under ** Currently Active Study
  Materials**.
\end{enumerate}

\subsubsection{Step 2: Interactive Cited
Q\&A}\label{step-2-interactive-cited-qa}

\begin{enumerate}
\def\labelenumi{\arabic{enumi}.}
\tightlist
\item
  Navigate to \textbf{{[}Ask Questions{]} Ask Questions}.
\item
  Select your document scope (or leave as \texttt{"All\ \allowbreak Materials"}).
\item
  Type a specific question: \emph{``What is the difference between a
  user thread and a kernel thread?''}
\item
  The assistant generates a grounded answer.
\item
  Click ** View Referenced Sources** to inspect the exact page number
  and text excerpt.
\end{enumerate}

\subsubsection{Step 3: Generate a Practice
Quiz}\label{step-3-generate-a-practice-quiz}

\begin{enumerate}
\def\labelenumi{\arabic{enumi}.}
\tightlist
\item
  Navigate to \textbf{{[}Practice Quiz{]} Practice Quiz}.
\item
  Enter a topic: \texttt{"Threads\ \allowbreak and\ \allowbreak Concurrency"}.
\item
  Choose 5 questions and click \textbf{{[}Execute{]} Generate Practice
  Quiz}.
\item
  A notification indicates background generation has started.
\item
  Select your answers using the radio buttons and click \textbf{Grade My
  Quiz }.
\item
  Review your score and detailed rationale for each question.
\end{enumerate}

\subsubsection{Step 4: Create Flashcards and
Summaries}\label{step-4-create-flashcards-and-summaries}

\begin{enumerate}
\def\labelenumi{\arabic{enumi}.}
\tightlist
\item
  Navigate to \textbf{{[}Flashcards{]} Flashcards}, enter
  \texttt{"Process\ \allowbreak Synchronization"}, and click \textbf{Create
  Flashcards}.
\item
  Expand each card to test your definition recall.
\item
  Navigate to \textbf{{[}Topic Summaries{]} Topic Summaries}, enter a
  chapter name, and review the structured key takeaway notes.
\end{enumerate}

\subsubsection{Step 5: Review and Manage Study
History}\label{step-5-review-and-manage-study-history}

\begin{enumerate}
\def\labelenumi{\arabic{enumi}.}
\tightlist
\item
  Navigate to \textbf{{[}Study History{]} Study History}.
\item
  Filter activities by \emph{``{[}Practice Quiz{]} Practice Quizzes''}
  or \emph{``{[}Ask Questions{]} Chat Questions''}.
\item
  Click {[}Delete{]} on any entry to remove it, or use \textbf{Clear All
  History} to reset your logs.
\end{enumerate}

\subsection{13.3 Stopping and Restarting the
Application}\label{stopping-and-restarting-the-application}

\begin{itemize}
\tightlist
\item
  \textbf{To Stop:} In the terminal running Streamlit, press
  \texttt{Ctrl\ \allowbreak +\ \allowbreak C}.
\item
  \textbf{To Restart:} Execute \texttt{streamlit\ \allowbreak run\ \allowbreak app.\allowbreak py}. All
  previously uploaded documents in \texttt{data/\allowbreak uploaded/\allowbreak } and history
  in \texttt{data/\allowbreak history.\allowbreak json} persist automatically across restarts.
\end{itemize}

\section{Chapter 14: Testing and
Validation}\label{chapter-14-testing-and-validation}

\subsection{14.1 Automated Testing
Architecture}\label{automated-testing-architecture}

The StudyRAG test suite is located in \texttt{tests/\allowbreak } and configured via
\texttt{pytest.\allowbreak ini}. It covers all critical subsystems: - Document
extraction and cleaning. - Recursive chunking and metadata preservation.
- SentenceTransformer embeddings. - Vector database operations (FAISS \&
ChromaDB). - Query processing and validation. - Relevance score
normalization and threshold filtering. - LLM initialization and provider
routing. - RAG pipeline orchestration and anti-hallucination fallback
logic.

\subsection{14.2 Running the Automated Test
Suite}\label{running-the-automated-test-suite}

To run all automated tests:

\begin{Shaded}
\begin{Highlighting}[]
\VariableTok{USE\_TF}\OperatorTok{=}\NormalTok{0 }\VariableTok{USE\_TORCH}\OperatorTok{=}\NormalTok{1 }\ExtensionTok{pytest}
\end{Highlighting}
\end{Shaded}

Expected test session summary:

\begin{verbatim}
============================= test session starts ==============================
platform darwin -- Python 3.12.2, pytest-8.1.1, pluggy-1.5.0
rootdir: /Users/kiritsomani/Desktop/rag-study-assistant
configfile: pytest.ini

tests/test_chunking.py ...                                               [ 10%]
tests/test_document_cleaner.py ....                                      [ 25%]
tests/test_embeddings.py ..                                              [ 32%]
tests/test_llm_manager.py ......                                         [ 53%]
tests/test_pdf_loader.py ..                                              [ 60%]
tests/test_query_processor.py ...                                        [ 71%]
tests/test_rag.py ..                                                     [ 78%]
tests/test_retrieval.py ...                                              [ 89%]
tests/test_vectorstore.py ...                                            [100%]

======================= 28 passed, 9 warnings in 23.21s ========================
\end{verbatim}

\subsection{14.3 Complete 28-Test Verification
Matrix}\label{complete-28-test-verification-matrix}

\begin{landscape}
\begingroup\scriptsize
\begin{longtable}[]{@{}
  >{\raggedright\arraybackslash\sloppy}p{(\linewidth - 14\tabcolsep) * \real{0.0800}}
  >{\raggedright\arraybackslash\sloppy}p{(\linewidth - 14\tabcolsep) * \real{0.1400}}
  >{\raggedright\arraybackslash\sloppy}p{(\linewidth - 14\tabcolsep) * \real{0.1900}}
  >{\raggedright\arraybackslash\sloppy}p{(\linewidth - 14\tabcolsep) * \real{0.1600}}
  >{\raggedright\arraybackslash\sloppy}p{(\linewidth - 14\tabcolsep) * \real{0.1100}}
  >{\raggedright\arraybackslash\sloppy}p{(\linewidth - 14\tabcolsep) * \real{0.1300}}
  >{\raggedright\arraybackslash\sloppy}p{(\linewidth - 14\tabcolsep) * \real{0.1300}}
  >{\raggedright\arraybackslash\sloppy}p{(\linewidth - 14\tabcolsep) * \real{0.0500}}@{}}
\toprule\noalign{}
\begin{minipage}[b]{\linewidth}\raggedright
Test ID
\end{minipage} & \begin{minipage}[b]{\linewidth}\raggedright
Test File
\end{minipage} & \begin{minipage}[b]{\linewidth}\raggedright
Test Function Name
\end{minipage} & \begin{minipage}[b]{\linewidth}\raggedright
Tested Feature / Logic
\end{minipage} & \begin{minipage}[b]{\linewidth}\raggedright
Preconditions \& Inputs
\end{minipage} & \begin{minipage}[b]{\linewidth}\raggedright
Expected Result
\end{minipage} & \begin{minipage}[b]{\linewidth}\raggedright
Actual Result
\end{minipage} & \begin{minipage}[b]{\linewidth}\raggedright
Status
\end{minipage} \\
\midrule\noalign{}
\endhead
\bottomrule\noalign{}
\endlastfoot
\textbf{TC-01} & \texttt{test\_\allowbreak chunking.\allowbreak py} &
\texttt{test\_\allowbreak create\_\allowbreak chunks} & Text chunking \& metadata & Sample
sentence, \texttt{size=\allowbreak 20,\allowbreak \ \allowbreak overlap=\allowbreak 5} & Multiple chunks with unique IDs
\& metadata & Unique IDs verified, metadata preserved &
\textbf{PASSED} \\
\textbf{TC-02} & \texttt{test\_\allowbreak chunking.\allowbreak py} &
\texttt{test\_\allowbreak invalid\_\allowbreak chunk\_\allowbreak size} & Parameter validation &
\texttt{chunk\_\allowbreak size=\allowbreak 0,\allowbreak \ \allowbreak overlap=\allowbreak 0} & Raises \texttt{ValueError} &
\texttt{ValueError} raised & \textbf{PASSED} \\
\textbf{TC-03} & \texttt{test\_\allowbreak chunking.\allowbreak py} &
\texttt{test\_\allowbreak invalid\_\allowbreak overlap} & Parameter validation &
\texttt{overlap=\allowbreak -1} and \texttt{overlap\ \allowbreak \textgreater{}=\allowbreak \ \allowbreak size} & Raises
\texttt{ValueError} & \texttt{ValueError} raised & \textbf{PASSED} \\
\textbf{TC-04} & \texttt{test\_\allowbreak document\_\allowbreak cleaner.\allowbreak py} &
\texttt{test\_\allowbreak clean\_\allowbreak extra\_\allowbreak spaces} & Whitespace normalization &
\texttt{"This\ \allowbreak \ \allowbreak \ \allowbreak has\ \allowbreak \ \allowbreak \ \allowbreak \ \allowbreak way\ \allowbreak \ \allowbreak too\ \allowbreak many\ \allowbreak spaces."} &
\texttt{"This\ \allowbreak has\ \allowbreak way\ \allowbreak too\ \allowbreak many\ \allowbreak spaces."} & String matched exactly &
\textbf{PASSED} \\
\textbf{TC-05} & \texttt{test\_\allowbreak document\_\allowbreak cleaner.\allowbreak py} &
\texttt{test\_\allowbreak clean\_\allowbreak newlines} & Paragraph preserving &
\texttt{"Paragraph\ \allowbreak 1\textbackslash{}n\textbackslash{}n\textbackslash{}n\textbackslash{}nParagraph\ \allowbreak 2"}
& \texttt{"Paragraph\ \allowbreak 1\textbackslash{}n\textbackslash{}nParagraph\ \allowbreak 2"}
& Exactly two newlines & \textbf{PASSED} \\
\textbf{TC-06} & \texttt{test\_\allowbreak document\_\allowbreak cleaner.\allowbreak py} &
\texttt{test\_\allowbreak empty\_\allowbreak text} & Edge case cleaning & Empty strings \&
whitespace strings & Returns \texttt{""} & Empty string returned &
\textbf{PASSED} \\
\textbf{TC-07} & \texttt{test\_\allowbreak document\_\allowbreak cleaner.\allowbreak py} &
\texttt{test\_\allowbreak clean\_\allowbreak documents} & Empty page removal & 3-page doc list
with 1 empty page & Len is 2, blank page dropped & 2 valid pages
returned & \textbf{PASSED} \\
\textbf{TC-08} & \texttt{test\_\allowbreak embeddings.\allowbreak py} &
\texttt{test\_\allowbreak embedding\_\allowbreak manager\_\allowbreak initialization} & Model initialization
& \texttt{EmbeddingManager(\allowbreak "minilm")\allowbreak } & Non-null HuggingFace embeddings
& Embedding object instantiated & \textbf{PASSED} \\
\textbf{TC-09} & \texttt{test\_\allowbreak embeddings.\allowbreak py} &
\texttt{test\_\allowbreak embedding\_\allowbreak dimension} & Dimension calculation & Embedding
query on ``test'' & Dimension equals 384 & 384 dimensions confirmed &
\textbf{PASSED} \\
\textbf{TC-10} & \texttt{test\_\allowbreak llm\_\allowbreak manager.\allowbreak py} &
\texttt{test\_\allowbreak ollama\_\allowbreak initialization} & Ollama provider setup &
\texttt{LLMManager(\allowbreak provider=\allowbreak "ollama")\allowbreak } & \texttt{ChatOllama} instance
created & Instantiated with base URL & \textbf{PASSED} \\
\textbf{TC-11} & \texttt{test\_\allowbreak llm\_\allowbreak manager.\allowbreak py} &
\texttt{test\_\allowbreak openai\_\allowbreak missing\_\allowbreak key} & OpenAI error handling & Missing
\texttt{OPENAI\_\allowbreak API\_\allowbreak KEY} & Raises \texttt{ValueError} &
\texttt{ValueError} raised & \textbf{PASSED} \\
\textbf{TC-12} & \texttt{test\_\allowbreak llm\_\allowbreak manager.\allowbreak py} &
\texttt{test\_\allowbreak openai\_\allowbreak initialization} & OpenAI provider setup & Mocked
\texttt{OPENAI\_\allowbreak API\_\allowbreak KEY=\allowbreak "test"} & \texttt{ChatOpenAI} instance created
& \texttt{ChatOpenAI} returned & \textbf{PASSED} \\
\textbf{TC-13} & \texttt{test\_\allowbreak llm\_\allowbreak manager.\allowbreak py} &
\texttt{test\_\allowbreak groq\_\allowbreak missing\_\allowbreak key} & Groq error handling & Missing
\texttt{GROQ\_\allowbreak API\_\allowbreak KEY} & Raises \texttt{ValueError} &
\texttt{ValueError} raised & \textbf{PASSED} \\
\textbf{TC-14} & \texttt{test\_\allowbreak llm\_\allowbreak manager.\allowbreak py} &
\texttt{test\_\allowbreak groq\_\allowbreak initialization} & Groq provider setup & Mocked
\texttt{GROQ\_\allowbreak API\_\allowbreak KEY=\allowbreak "test"} & \texttt{ChatOpenAI} with Groq base URL
& Correct base URL \& model mapping & \textbf{PASSED} \\
\textbf{TC-15} & \texttt{test\_\allowbreak llm\_\allowbreak manager.\allowbreak py} &
\texttt{test\_\allowbreak unsupported\_\allowbreak provider} & Invalid provider validation &
\texttt{LLMManager(\allowbreak provider=\allowbreak "unsupported")\allowbreak } & Raises \texttt{ValueError}
& \texttt{ValueError} raised & \textbf{PASSED} \\
\textbf{TC-16} & \texttt{test\_\allowbreak pdf\_\allowbreak loader.\allowbreak py} &
\texttt{test\_\allowbreak load\_\allowbreak pdf\_\allowbreak missing\_\allowbreak file} & File existence guard &
Non-existent file path & Raises \texttt{FileNotFoundError} &
\texttt{FileNotFoundError} raised & \textbf{PASSED} \\
\textbf{TC-17} & \texttt{test\_\allowbreak pdf\_\allowbreak loader.\allowbreak py} &
\texttt{test\_\allowbreak load\_\allowbreak pdf\_\allowbreak invalid\_\allowbreak extension} & File extension guard &
Path pointing to \texttt{.\allowbreak txt} file & Raises \texttt{ValueError} &
\texttt{ValueError} raised & \textbf{PASSED} \\
\textbf{TC-18} & \texttt{test\_\allowbreak query\_\allowbreak processor.\allowbreak py} &
\texttt{test\_\allowbreak process\_\allowbreak valid\_\allowbreak query} & Query sanitization &
\texttt{"\ \allowbreak \ \allowbreak \ \allowbreak What\ \allowbreak is\ \allowbreak \ \allowbreak RAG?\ \allowbreak \ \allowbreak \ \allowbreak "} & \texttt{"What\ \allowbreak is\ \allowbreak RAG?"} &
Sanitized query returned & \textbf{PASSED} \\
\textbf{TC-19} & \texttt{test\_\allowbreak query\_\allowbreak processor.\allowbreak py} &
\texttt{test\_\allowbreak empty\_\allowbreak query} & Empty query rejection & Empty and
whitespace queries & Raises \texttt{ValueError} & \texttt{ValueError}
raised & \textbf{PASSED} \\
\textbf{TC-20} & \texttt{test\_\allowbreak query\_\allowbreak processor.\allowbreak py} &
\texttt{test\_\allowbreak query\_\allowbreak too\_\allowbreak long} & Max length guard & String of 1,001
characters & Raises \texttt{ValueError} & \texttt{ValueError} raised &
\textbf{PASSED} \\
\textbf{TC-21} & \texttt{test\_\allowbreak rag.\allowbreak py} &
\texttt{test\_\allowbreak generate\_\allowbreak rag\_\allowbreak response\_\allowbreak no\_\allowbreak context} & Anti-hallucination
guard & Mocked retrieval returning \texttt{found=\allowbreak False} & Fallback
message, LLM NOT called & Zero LLM calls, refusal returned &
\textbf{PASSED} \\
\textbf{TC-22} & \texttt{test\_\allowbreak rag.\allowbreak py} &
\texttt{test\_\allowbreak generate\_\allowbreak rag\_\allowbreak response\_\allowbreak with\_\allowbreak context} & Grounded
generation & Mocked context with \texttt{math.\allowbreak pdf} & Answer returned
with citations & Answer returned with \texttt{math.\allowbreak pdf} &
\textbf{PASSED} \\
\textbf{TC-23} & \texttt{test\_\allowbreak retrieval.\allowbreak py} &
\texttt{test\_\allowbreak normalize\_\allowbreak score} & Distance normalization & L2=0.0 \&
1.0; Cosine=0.0 \& 0.2 & Maps to 1.0, 0.5; 1.0, 0.8 & Normalized values
accurate & \textbf{PASSED} \\
\textbf{TC-24} & \texttt{test\_\allowbreak retrieval.\allowbreak py} &
\texttt{test\_\allowbreak filter\_\allowbreak by\_\allowbreak threshold} & Thresholding logic & 3 docs with
L2 scores, \(\tau=0.5\) & 2 docs kept, sorted descending & 2 docs kept,
1 doc filtered & \textbf{PASSED} \\
\textbf{TC-25} & \texttt{test\_\allowbreak retrieval.\allowbreak py} & \texttt{test\_\allowbreak metrics} &
Evaluation metrics & Retrieved list, expected=``doc1'' & Precision=0.33,
Hit=1, Recall=1.0 & Calculations verified & \textbf{PASSED} \\
\textbf{TC-26} & \texttt{test\_\allowbreak vectorstore.\allowbreak py} &
\texttt{test\_\allowbreak faiss\_\allowbreak vector\_\allowbreak store} & FAISS lifecycle & Temp dir, sample
documents & Add, save, load, search verified & Stack LIFO document
recovered & \textbf{PASSED} \\
\textbf{TC-27} & \texttt{test\_\allowbreak vectorstore.\allowbreak py} &
\texttt{test\_\allowbreak chroma\_\allowbreak vector\_\allowbreak store} & Chroma lifecycle & Temp dir,
sample documents & Add, save, load, search verified & Stack LIFO
document recovered & \textbf{PASSED} \\
\textbf{TC-28} & \texttt{test\_\allowbreak vectorstore.\allowbreak py} &
\texttt{test\_\allowbreak raw\_\allowbreak chunks\_\allowbreak conversion} & Chunk conversion & Raw chunk
dictionary list & LangChain \texttt{Document} list & Page, source,
chunk\_id preserved & \textbf{PASSED} \\
\end{longtable}
\endgroup
\end{landscape}

\section{Chapter 15: Current Progress and Development
Status}\label{chapter-15-current-progress-and-development-status}

\subsection{15.1 Completed Features \& Verified
Capabilities}\label{completed-features-verified-capabilities}

The following features are 100\% implemented, integrated, and verified
by automated testing:

\begin{enumerate}
\def\labelenumi{\arabic{enumi}.}
\tightlist
\item
  \textbf{Multi-PDF Ingestion Pipeline}: Full support for single and
  batch PDF uploads, SHA-256 deduplication, PyMuPDF page-by-page
  extraction, and error handling for corrupted/non-PDF files
  (\texttt{src/\allowbreak ingestion/\allowbreak }).
\item
  \textbf{Text Normalization Engine}: Whitespace and newline collapsing
  with paragraph integrity preservation
  (\texttt{src/\allowbreak ingestion/\allowbreak document\_\allowbreak cleaner.\allowbreak py}).
\item
  \textbf{Recursive Character Chunking}: LangChain recursive splitting
  with 5-tier separator fallback and metadata propagation
  (\texttt{src/\allowbreak ingestion/\allowbreak chunker.\allowbreak py}).
\item
  \textbf{SentenceTransformer Embeddings}: Local 384-dimensional dense
  vector embeddings with dynamic dimension caching
  (\texttt{src/\allowbreak embeddings/\allowbreak embedding\_\allowbreak manager.\allowbreak py}).
\item
  \textbf{Dual Vector Store Backends}: Fully operational FAISS
  (in-memory flat L2 index with disk serialization) and ChromaDB
  (SQLite-backed with cosine distance collection configuration)
  (\texttt{src/\allowbreak vectorstore/\allowbreak }).
\item
  \textbf{Query Processing \& Sanitization}: Input validation,
  whitespace collapse, and 1,000-character safety guard
  (\texttt{src/\allowbreak retrieval/\allowbreak query\_\allowbreak processor.\allowbreak py}).
\item
  \textbf{Relevance Normalization \& Threshold Filtering}: Mathematical
  score mapping for L2 and Cosine metrics with anti-hallucination
  threshold cutoff (\texttt{src/\allowbreak retrieval/\allowbreak score\_\allowbreak filter.\allowbreak py}).
\item
  \textbf{Source-Specific Retrieval}: Selective filtering to single
  documents or across all uploaded materials
  (\texttt{src/\allowbreak retrieval/\allowbreak retriever.\allowbreak py}).
\item
  \textbf{Multi-Provider LLM Integration}: Dynamic switching between
  Groq Cloud (ultra-fast LPU), OpenAI (GPT-4o), and local Ollama
  (\texttt{src/\allowbreak generation/\allowbreak llm\_\allowbreak manager.\allowbreak py}).
\item
  \textbf{4 Dedicated Study Artifact Generators}:

  \begin{itemize}
  \tightlist
  \item
    Grounded Q\&A with verifiable page-level references.
  \item
    Multiple Choice Practice Quizzes with automated form-based grading.
  \item
    Digital Flashcard decks with front prompt / back definition
    expansion.
  \item
    Chapter Revision Summaries with key takeaway bullet points.
  \end{itemize}
\item
  \textbf{Non-Blocking Background Threading}: Daemon threads with
  Streamlit script execution context tracking generation jobs
  asynchronously without UI lockup (\texttt{app.\allowbreak py}).
\item
  \textbf{Persistent Study Activity Store}: Atomic JSON persistence
  (\texttt{data/\allowbreak history.\allowbreak json}) with UUID tracking, activity filtering,
  and entry deletion.
\item
  \textbf{Modern Student-Centric Interface}: Fully responsive Streamlit
  UI with {[}Dark Mode{]} Dark Mode and {[}Light Mode{]} Light Mode
  themes and custom CSS.
\item
  \textbf{Automated Test Suite}: 28 passing unit and integration tests
  (\texttt{tests/\allowbreak }).
\end{enumerate}

\subsection{15.2 Partially Completed Features \&
Nuances}\label{partially-completed-features-nuances}

\begin{enumerate}
\def\labelenumi{\arabic{enumi}.}
\tightlist
\item
  \textbf{ChromaDB LangChain Deprecation Warning}:

  \begin{itemize}
  \tightlist
  \item
    \emph{Status:} Chroma functions properly, but LangChain emits a
    deprecation warning recommending migration from
    \texttt{langchain\_\allowbreak community.\allowbreak vectorstores.\allowbreak Chroma} to
    \texttt{langchain\_\allowbreak chroma.\allowbreak Chroma}.
  \item
    \emph{Impact:} Harmless in LangChain 0.1/0.2, but should be updated
    before LangChain 1.0.
  \end{itemize}
\item
  \textbf{FAISS Deserialization Flag}:

  \begin{itemize}
  \tightlist
  \item
    \emph{Status:} FAISS loading requires
    \texttt{allow\_\allowbreak dangerous\_\allowbreak deserialization=\allowbreak True} due to pickle
    security checks in newer LangChain versions.
  \item
    \emph{Impact:} Safe for local user files, but multi-user
    environments should transition to trusted formats or isolated
    storage.
  \end{itemize}
\end{enumerate}

\subsection{15.3 Work Currently in Progress / Immediate Next
Steps}\label{work-currently-in-progress-immediate-next-steps}

\begin{enumerate}
\def\labelenumi{\arabic{enumi}.}
\tightlist
\item
  \textbf{OCR Integration for Scanned PDFs}: Integrating Tesseract OCR
  or \texttt{pdf2image} for documents containing scanned photocopies
  without extractable text streams.
\item
  \textbf{Flashcard Spaced Repetition (SM-2) Scheduling}: Adding
  interval tracking (1 day, 3 days, 7 days) to the flashcard interface.
\end{enumerate}

\subsection{15.4 Technical Debt \& Known
Limitations}\label{technical-debt-known-limitations}

\begin{enumerate}
\def\labelenumi{\arabic{enumi}.}
\tightlist
\item
  \textbf{In-Process Background Jobs}: Background generation threads run
  inside the Streamlit Python process. If the server is terminated
  during generation, in-flight jobs are interrupted.
\item
  \textbf{Single-User Session Isolation}: The current Streamlit
  application shares the file system and vector store across connected
  browser tabs. Multi-user enterprise deployments require user
  authentication and partitioned vector collections.
\end{enumerate}

\section{Chapter 16: Configuration and
Customization}\label{chapter-16-configuration-and-customization}

\subsection{\texorpdfstring{16.1 Application Configuration Reference
(\texttt{config/\allowbreak settings.\allowbreak yaml})}{16.1 Application Configuration Reference (config/settings.yaml)}}\label{application-configuration-reference-configsettings.yaml}

The primary system defaults are declared in
\texttt{config/\allowbreak settings.\allowbreak yaml}:

\begin{Shaded}
\begin{Highlighting}[]
\FunctionTok{app}\KeywordTok{:}
\AttributeTok{  }\FunctionTok{title}\KeywordTok{:}\AttributeTok{ }\StringTok{"StudyRAG: AI{-}Powered Study Assistant"}
\AttributeTok{  }\FunctionTok{version}\KeywordTok{:}\AttributeTok{ }\StringTok{"1.0.0"}

\FunctionTok{chunking}\KeywordTok{:}
\AttributeTok{  }\FunctionTok{default\_strategy}\KeywordTok{:}\AttributeTok{ }\StringTok{"recursive"}
\AttributeTok{  }\FunctionTok{chunk\_size}\KeywordTok{:}\AttributeTok{ }\DecValTok{500}\CommentTok{              \# Maximum character length per chunk}
\AttributeTok{  }\FunctionTok{chunk\_overlap}\KeywordTok{:}\AttributeTok{ }\DecValTok{50}\CommentTok{            \# Overlap in characters between adjacent chunks}

\FunctionTok{retrieval}\KeywordTok{:}
\AttributeTok{  }\FunctionTok{top\_k}\KeywordTok{:}\AttributeTok{ }\DecValTok{5}\CommentTok{                     \# Number of chunks retrieved per query}
\AttributeTok{  }\FunctionTok{similarity\_threshold}\KeywordTok{:}\AttributeTok{ }\FloatTok{0.2}\CommentTok{    \# Minimum normalized relevance score (0.0 to 1.0)}
\end{Highlighting}
\end{Shaded}

\subsection{\texorpdfstring{16.2 Environment Variables Reference
(\texttt{.\allowbreak env})}{16.2 Environment Variables Reference (.env)}}\label{environment-variables-reference-.env}

\begingroup\footnotesize
\begin{longtable}[]{@{}
  >{\raggedright\arraybackslash\sloppy}p{(\linewidth - 8\tabcolsep) * \real{0.2200}}
  >{\raggedright\arraybackslash\sloppy}p{(\linewidth - 8\tabcolsep) * \real{0.0800}}
  >{\raggedright\arraybackslash\sloppy}p{(\linewidth - 8\tabcolsep) * \real{0.1500}}
  >{\raggedright\arraybackslash\sloppy}p{(\linewidth - 8\tabcolsep) * \real{0.1500}}
  >{\raggedright\arraybackslash\sloppy}p{(\linewidth - 8\tabcolsep) * \real{0.3600}}@{}}
\toprule\noalign{}
\begin{minipage}[b]{\linewidth}\raggedright
Variable Name
\end{minipage} & \begin{minipage}[b]{\linewidth}\raggedright
Required
\end{minipage} & \begin{minipage}[b]{\linewidth}\raggedright
Default Value
\end{minipage} & \begin{minipage}[b]{\linewidth}\raggedright
Accepted Values
\end{minipage} & \begin{minipage}[b]{\linewidth}\raggedright
Purpose \& Security Implications
\end{minipage} \\
\midrule\noalign{}
\endhead
\bottomrule\noalign{}
\endlastfoot
\texttt{LLM\_\allowbreak PROVIDER} & No & \texttt{"groq"} & \texttt{"groq"},
\texttt{"openai"}, \texttt{"ollama"} & Selects default inference
provider. \\
\texttt{GROQ\_\allowbreak API\_\allowbreak KEY} & Conditional & None & String
(\texttt{gsk\_\allowbreak .\allowbreak .\allowbreak .}) & API key for Groq Cloud. Required if
\texttt{LLM\_\allowbreak PROVIDER=\allowbreak groq}. \\
\texttt{OPENAI\_\allowbreak API\_\allowbreak KEY} & Conditional & None & String
(\texttt{sk-.\allowbreak .\allowbreak .}) & API key for OpenAI. Required if
\texttt{LLM\_\allowbreak PROVIDER=\allowbreak openai}. \\
\texttt{OLLAMA\_\allowbreak BASE\_\allowbreak URL} & No & \texttt{"http:/\allowbreak /\allowbreak localhost:11434"} &
URL string & Address of local Ollama HTTP server. \\
\texttt{MODEL\_\allowbreak NAME} & No & \texttt{"llama3"} & String (model tag) &
Model name for local Ollama inference. \\
\texttt{EMBEDDING\_\allowbreak MODEL} & No & \texttt{"minilm"} & \texttt{"minilm"},
\texttt{"bge"} & Selects embedding model architecture. \\
\texttt{VECTOR\_\allowbreak STORE} & No & \texttt{"faiss"} & \texttt{"faiss"},
\texttt{"chroma"} & Selects vector storage backend. \\
\texttt{USE\_\allowbreak TF} & Recommended & \texttt{"0"} & \texttt{"0"},
\texttt{"1"} & Disables TensorFlow in Transformers to prevent macOS
segfaults. \\
\texttt{USE\_\allowbreak TORCH} & Recommended & \texttt{"1"} & \texttt{"0"},
\texttt{"1"} & Forces PyTorch backend for SentenceTransformers. \\
\end{longtable}
\endgroup

\subsection{16.3 Customizing UI Themes and Visual
Design}\label{customizing-ui-themes-and-visual-design}

Themes are controlled via the \texttt{apply\_\allowbreak theme\_\allowbreak css(\allowbreak )\allowbreak } function in
\texttt{app.\allowbreak py}. Colors can be tuned by editing CSS variables: - Dark
Mode Card Background: \texttt{\#1E293B} - Dark Mode Text Gradient:
\texttt{linear-gradient(\allowbreak 90deg,\allowbreak \ \allowbreak \#818CF8,\allowbreak \ \allowbreak \#38BDF8)\allowbreak } - Primary Accent
Button:
\texttt{linear-gradient(\allowbreak 135deg,\allowbreak \ \allowbreak \#6366F1\ \allowbreak 0\%,\allowbreak \ \allowbreak \#0284C7\ \allowbreak 100\%)\allowbreak }

\section{Chapter 17: Error Handling and
Troubleshooting}\label{chapter-17-error-handling-and-troubleshooting}

\subsection{17.1 Practical Diagnostic
Matrix}\label{practical-diagnostic-matrix}

\begingroup\footnotesize
\begin{longtable}[]{@{}
  >{\raggedright\arraybackslash\sloppy}p{(\linewidth - 6\tabcolsep) * \real{0.2400}}
  >{\raggedright\arraybackslash\sloppy}p{(\linewidth - 6\tabcolsep) * \real{0.2200}}
  >{\raggedright\arraybackslash\sloppy}p{(\linewidth - 6\tabcolsep) * \real{0.3200}}
  >{\raggedright\arraybackslash\sloppy}p{(\linewidth - 6\tabcolsep) * \real{0.1800}}@{}}
\toprule\noalign{}
\begin{minipage}[b]{\linewidth}\raggedright
Symptom / Error Message
\end{minipage} & \begin{minipage}[b]{\linewidth}\raggedright
Likely Cause
\end{minipage} & \begin{minipage}[b]{\linewidth}\raggedright
Step-by-Step Fix
\end{minipage} & \begin{minipage}[b]{\linewidth}\raggedright
Command to Verify
\end{minipage} \\
\midrule\noalign{}
\endhead
\bottomrule\noalign{}
\endlastfoot
\texttt{Fatal\ \allowbreak Python\ \allowbreak error:\ \allowbreak Segmentation\ \allowbreak fault} during
\texttt{import\ \allowbreak transformers} on macOS & Anaconda environment on Apple
Silicon where TensorFlow attempts to execute C-extension preload check.
& Set \texttt{USE\_\allowbreak TF=\allowbreak 0} and \texttt{USE\_\allowbreak TORCH=\allowbreak 1} in the shell or at
top of Python scripts before importing transformers. &
\texttt{USE\_\allowbreak TF=\allowbreak 0\ \allowbreak python3\ \allowbreak -c\ \allowbreak "import\ \allowbreak transformers;\ \allowbreak print(\allowbreak \textquotesingle{}OK\textquotesingle{})\allowbreak "} \\
\texttt{ValueError:\ \allowbreak OPENAI\_\allowbreak API\_\allowbreak KEY\ \allowbreak environment\ \allowbreak variable\ \allowbreak is\ \allowbreak not\ \allowbreak set.}
& OpenAI selected as provider without defining \texttt{OPENAI\_\allowbreak API\_\allowbreak KEY}
in \texttt{.\allowbreak env}. & Add \texttt{OPENAI\_\allowbreak API\_\allowbreak KEY=\allowbreak sk-.\allowbreak .\allowbreak .} to
\texttt{.\allowbreak env} or switch provider to \texttt{groq} or \texttt{ollama}. &
\texttt{python3\ \allowbreak -c\ \allowbreak "from\ \allowbreak dotenv\ \allowbreak import\ \allowbreak load\_\allowbreak dotenv;\ \allowbreak import\ \allowbreak os;\ \allowbreak load\_\allowbreak dotenv(\allowbreak )\allowbreak ;\ \allowbreak print(\allowbreak bool(\allowbreak os.\allowbreak getenv(\allowbreak \textquotesingle{}OPENAI\_\allowbreak API\_\allowbreak KEY\textquotesingle{})\allowbreak )\allowbreak )\allowbreak "} \\
\texttt{ValueError:\ \allowbreak GROQ\_\allowbreak API\_\allowbreak KEY\ \allowbreak environment\ \allowbreak variable\ \allowbreak is\ \allowbreak not\ \allowbreak set.}
& Groq selected as provider without defining \texttt{GROQ\_\allowbreak API\_\allowbreak KEY} in
\texttt{.\allowbreak env}. & Add \texttt{GROQ\_\allowbreak API\_\allowbreak KEY=\allowbreak gsk\_\allowbreak .\allowbreak .\allowbreak .} to \texttt{.\allowbreak env}
(obtain free key at console.groq.com). &
\texttt{python3\ \allowbreak -c\ \allowbreak "from\ \allowbreak dotenv\ \allowbreak import\ \allowbreak load\_\allowbreak dotenv;\ \allowbreak import\ \allowbreak os;\ \allowbreak load\_\allowbreak dotenv(\allowbreak )\allowbreak ;\ \allowbreak print(\allowbreak bool(\allowbreak os.\allowbreak getenv(\allowbreak \textquotesingle{}GROQ\_\allowbreak API\_\allowbreak KEY\textquotesingle{})\allowbreak )\allowbreak )\allowbreak "} \\
\texttt{ConnectionError} /
\texttt{Failed\ \allowbreak to\ \allowbreak connect\ \allowbreak to\ \allowbreak localhost:11434} & Ollama daemon is not
running on host machine. & Start the Ollama service: execute
\texttt{ollama\ \allowbreak serve} in a separate terminal. Verify with
\texttt{ollama\ \allowbreak run\ \allowbreak llama3}. &
\texttt{curl\ \allowbreak http:/\allowbreak /\allowbreak localhost:11434/\allowbreak api/\allowbreak version} \\
\texttt{Warning:\ \allowbreak No\ \allowbreak extractable\ \allowbreak text\ \allowbreak found\ \allowbreak in\ \allowbreak \{file.\allowbreak pdf\}.\ \allowbreak Might\ \allowbreak be\ \allowbreak a\ \allowbreak scanned\ \allowbreak image.}
& The uploaded PDF consists of scanned bitmap images without embedded
text fonts. & Use an OCR tool (e.g.~Adobe Acrobat OCR, ocrmypdf) to
generate a searchable PDF before uploading. & Open PDF in viewer and
verify if text can be highlighted and copied. \\
\texttt{I\ \allowbreak could\ \allowbreak not\ \allowbreak find\ \allowbreak enough\ \allowbreak information\ \allowbreak in\ \allowbreak the\ \allowbreak uploaded\ \allowbreak documents.\allowbreak .\allowbreak .}
& Query relevance score fell below \(\tau=0.20\) or document does not
cover the topic. & Verify that relevant study notes are uploaded in
Study Library, or lower the similarity threshold in
\texttt{settings.\allowbreak yaml}. & Check retrieval scores in logger output. \\
\texttt{LangChainDeprecationWarning:\ \allowbreak The\ \allowbreak class\ \allowbreak Chroma\ \allowbreak was\ \allowbreak deprecated.\allowbreak .\allowbreak .}
& Using \texttt{langchain\_\allowbreak community.\allowbreak vectorstores.\allowbreak Chroma} instead of
standalone \texttt{langchain-chroma}. & Warning only; does not affect
functionality. Can be resolved by running
\texttt{pip\ \allowbreak install\ \allowbreak -U\ \allowbreak langchain-chroma}. &
\texttt{pytest\ \allowbreak tests/\allowbreak test\_\allowbreak vectorstore.\allowbreak py} \\
Port 8501 already in use & Another Streamlit instance is running. &
Specify alternative port:
\texttt{streamlit\ \allowbreak run\ \allowbreak app.\allowbreak py\ \allowbreak -\/\allowbreak -server.\allowbreak port\ \allowbreak 8502}. &
\texttt{lsof\ \allowbreak -i\ \allowbreak :8501} \\
\end{longtable}
\endgroup

\section{Chapter 18: Security, Privacy, and
Reliability}\label{chapter-18-security-privacy-and-reliability}

\subsection{18.1 Secrets Management \& Data
Hygiene}\label{secrets-management-data-hygiene}

\begin{itemize}
\tightlist
\item
  \textbf{API Keys}: All external API keys (\texttt{GROQ\_\allowbreak API\_\allowbreak KEY},
  \texttt{OPENAI\_\allowbreak API\_\allowbreak KEY}) are loaded via \texttt{python-dotenv} from
  the root \texttt{.\allowbreak env} file. They are never written to source code,
  logged in terminal outputs, or exposed to the frontend browser client.
\item
  \textbf{Repository Isolation}: \texttt{.\allowbreak gitignore} explicitly excludes
  \texttt{.\allowbreak env}, \texttt{data/\allowbreak uploaded/\allowbreak }, \texttt{data/\allowbreak vectorstores/\allowbreak },
  and \texttt{data/\allowbreak history.\allowbreak json}, ensuring student notes and API
  credentials cannot be accidentally committed to public repositories.
\end{itemize}

\subsection{18.2 Local-First Privacy
Guarantee}\label{local-first-privacy-guarantee}

For students, researchers, and institutions handling confidential course
exams, proprietary research data, or copyrighted textbooks: - StudyRAG
can operate in a \textbf{100\% Air-Gapped Local Mode}. - Embeddings are
generated entirely on the local CPU via SentenceTransformers
(\texttt{all-MiniLM-L6-v2}). - Vector indices are stored locally in
\texttt{data/\allowbreak vectorstores/\allowbreak }. - LLM inference can be routed to a local
\textbf{Ollama} daemon. - In this configuration, \textbf{zero bytes of
document text or queries leave the host machine}.

\subsection{18.3 Input Validation and Defense in
Depth}\label{input-validation-and-defense-in-depth}

\begin{itemize}
\tightlist
\item
  \textbf{File Validation}: \texttt{pdf\_\allowbreak loader.\allowbreak py} validates file
  existence, enforces \texttt{.\allowbreak pdf} extension checks, and wraps PyMuPDF
  in defensive try-except blocks to reject corrupt or malformed file
  payloads.
\item
  \textbf{Query Bounding}: \texttt{query\_\allowbreak processor.\allowbreak py} strips control
  characters, normalizes whitespace, and rejects queries exceeding 1,000
  characters, mitigating buffer overflow and context saturation attacks.
\item
  \textbf{Strict Anti-Hallucination Guardrails}: Prompts explicitly
  forbid the model from making unsupported assumptions and mandate
  page-level citations.
\end{itemize}

\section{Chapter 19: Deployment and Production
Readiness}\label{chapter-19-deployment-and-production-readiness}

\subsection{19.1 Current Deployment
Profile}\label{current-deployment-profile}

StudyRAG is optimized for \textbf{Single-User Desktop \& Local
Workstation Deployment}: - Web Server: Streamlit embedded Tornado
server. - Database: Embedded FAISS / SQLite ChromaDB. - Concurrency:
In-process Python daemon threads.

\subsection{19.2 Recommended Production Architecture (Multi-User
Cloud)}\label{recommended-production-architecture-multi-user-cloud}

To scale StudyRAG into an enterprise campus-wide SaaS application, the
following architecture is recommended:

\begin{Shaded}
\begin{Highlighting}[]
\NormalTok{graph TD}
\NormalTok{    CLIENT[Web Browser / Mobile Client] {-}{-}\textgreater{} CDN[Cloudflare CDN / SSL]}
\NormalTok{    CDN {-}{-}\textgreater{} LB[Application Load Balancer]}
    
\NormalTok{    subgraph WEB\_CLUSTER ["Containerized App Tier"]}
\NormalTok{        LB {-}{-}\textgreater{} POD1[Streamlit / React Container 1]}
\NormalTok{        LB {-}{-}\textgreater{} POD2[Streamlit / React Container 2]}
\NormalTok{    end}
    
\NormalTok{    subgraph ASYNC\_TIER ["Distributed Task Queue"]}
\NormalTok{        POD1 \& POD2 {-}{-}\textgreater{} REDIS[(Redis Queue)]}
\NormalTok{        REDIS {-}{-}\textgreater{} WORKER1[Celery Worker 1]}
\NormalTok{        WORKER1 {-}{-}\textgreater{} GROQ[Groq Cloud / vLLM Cluster]}
\NormalTok{    end}
    
\NormalTok{    subgraph DATA\_TIER ["Enterprise Data Tier"]}
\NormalTok{        POD1 \& POD2 {-}{-}\textgreater{} PG[(PostgreSQL + pgvector)]}
\NormalTok{        POD1 \& POD2 {-}{-}\textgreater{} S3[(AWS S3 / GCS Document Bucket)]}
\NormalTok{    end}
\end{Highlighting}
\end{Shaded}

\subsubsection{Production Transition
Checklist:}\label{production-transition-checklist}

\begin{enumerate}
\def\labelenumi{\arabic{enumi}.}
\tightlist
\item
  \textbf{Containerization}: Deploy via a multi-stage Docker container
  (\texttt{python:3.12-slim}).
\item
  \textbf{Vector Database}: Migrate from local FAISS/Chroma files to a
  managed vector database (e.g.~\textbf{pgvector} or \textbf{Qdrant
  Cloud}) with user-isolated tenant namespaces.
\item
  \textbf{Task Queue}: Replace in-process \texttt{threading.\allowbreak Thread} with
  \textbf{Celery + Redis} to guarantee that long-running quiz/flashcard
  jobs survive server restarts.
\item
  \textbf{Authentication}: Implement OAuth2 / OIDC login (Google,
  GitHub, Canvas LMS SSO).
\end{enumerate}

\section{Chapter 20: Limitations and Future
Scope}\label{chapter-20-limitations-and-future-scope}

\subsection{20.1 Current Functional and Technical
Limitations}\label{current-functional-and-technical-limitations}

\begingroup\footnotesize
\begin{longtable}[]{@{}
  >{\raggedright\arraybackslash\sloppy}p{(\linewidth - 6\tabcolsep) * \real{0.2200}}
  >{\raggedright\arraybackslash\sloppy}p{(\linewidth - 6\tabcolsep) * \real{0.2000}}
  >{\raggedright\arraybackslash\sloppy}p{(\linewidth - 6\tabcolsep) * \real{0.2600}}
  >{\raggedright\arraybackslash\sloppy}p{(\linewidth - 6\tabcolsep) * \real{0.2800}}@{}}
\toprule\noalign{}
\begin{minipage}[b]{\linewidth}\raggedright
Limitation
\end{minipage} & \begin{minipage}[b]{\linewidth}\raggedright
Impact
\end{minipage} & \begin{minipage}[b]{\linewidth}\raggedright
Root Cause
\end{minipage} & \begin{minipage}[b]{\linewidth}\raggedright
Workaround / Recommendation
\end{minipage} \\
\midrule\noalign{}
\endhead
\bottomrule\noalign{}
\endlastfoot
\textbf{No Native OCR for Scanned PDFs} & Scanned document images yield
empty text chunks. & PyMuPDF relies on embedded text glyphs and font
dictionaries. & Pre-process scanned PDFs using OCR software prior to
upload. \\
\textbf{Volatile In-Memory Thread Jobs} & Generating jobs are
interrupted if the Streamlit server restarts. & Uses lightweight
\texttt{threading.\allowbreak Thread} in-process workers. & Avoid restarting the
server while background generation is active. \\
\textbf{English-Biased Embeddings} & Retrieval performance drops on
non-Latin scripts (e.g., Arabic, Chinese). & \texttt{all-MiniLM-L6-v2}
is primarily trained on English corpora. & Switch to multilingual
embedding models such as
\texttt{paraphrase-multilingual-MiniLM-L12-v2}. \\
\textbf{Single-User Data Namespace} & Multiple browser sessions share
the same active library. & Streamlit session state shares underlying
\texttt{data/\allowbreak } directories. & Run separate project directories for
different study subjects or users. \\
\end{longtable}
\endgroup

\subsection{20.2 Prioritized Development
Roadmap}\label{prioritized-development-roadmap}

\begin{verbatim}
+-----------------------------------------------------------------------------------+
|                        STUDYRAG FUTURE DEVELOPMENT ROADMAP                        |
+-----------------------------------------------------------------------------------+
|  Phase 1 (Immediate - Q4 2026):                                                   |
|  - Tesseract OCR integration for automatic image-to-text fallback.                |
|  - Anki Deck (.apkg) export functionality for mobile flashcard review.            |
|                                                                                   |
|  Phase 2 (Medium Term - Q1 2027):                                                 |
|  - BGE-Reranker cross-encoder integration for two-stage retrieval refinement.     |
|  - Multi-modal image and equation extraction from lecture slides.                 |
|  - PDF export for generated practice quizzes and revision study guides.           |
|                                                                                   |
|  Phase 3 (Long Term - Q2 2027):                                                   |
|  - Canvas / Blackboard / Moodle LMS direct course syllabus sync.                  |
|  - Multi-user authentication, study groups, and collaborative shared decks.       |
+-----------------------------------------------------------------------------------+
\end{verbatim}

\section{Chapter 21: Developer Handover
Guide}\label{chapter-21-developer-handover-guide}

\subsection{21.1 Recommended Codebase Onboarding Reading
Order}\label{recommended-codebase-onboarding-reading-order}

For a new software engineer or contributor joining the project, study
the codebase in the following sequence:

\begin{enumerate}
\def\labelenumi{\arabic{enumi}.}
\tightlist
\item
  \textbf{\texttt{config/\allowbreak settings.\allowbreak yaml}}: Understand default
  hyperparameters (chunk size, overlap, top\_k, threshold).
\item
  \textbf{\texttt{src/\allowbreak ingestion/\allowbreak }}:

  \begin{itemize}
  \tightlist
  \item
    Start with \texttt{pdf\_\allowbreak loader.\allowbreak py} to see how text and 1-indexed
    pages are extracted.
  \item
    Read \texttt{document\_\allowbreak cleaner.\allowbreak py} and \texttt{chunker.\allowbreak py} to see
    text splitting.
  \item
    Review \texttt{ingestion\_\allowbreak pipeline.\allowbreak py} which unifies these steps.
  \end{itemize}
\item
  \textbf{\texttt{src/\allowbreak embeddings/\allowbreak embedding\_\allowbreak manager.\allowbreak py}}: Learn how
  HuggingFace embeddings are wrapped and cached.
\item
  \textbf{\texttt{src/\allowbreak vectorstore/\allowbreak }}:

  \begin{itemize}
  \tightlist
  \item
    Inspect \texttt{faiss\_\allowbreak store.\allowbreak py} and \texttt{chroma\_\allowbreak store.\allowbreak py} to
    understand the storage and search wrappers.
  \item
    Review \texttt{vector\_\allowbreak store\_\allowbreak manager.\allowbreak py} for document conversion
    and factory routing.
  \end{itemize}
\item
  \textbf{\texttt{src/\allowbreak retrieval/\allowbreak }}:

  \begin{itemize}
  \tightlist
  \item
    Read \texttt{query\_\allowbreak processor.\allowbreak py} for validation.
  \item
    Read \texttt{score\_\allowbreak filter.\allowbreak py} for mathematical score normalizations
    (\(S_{\text{FAISS}}\) and \(S_{\text{Chroma}}\)).
  \item
    Review \texttt{retriever.\allowbreak py} and \texttt{retrieval\_\allowbreak pipeline.\allowbreak py}.
  \end{itemize}
\item
  \textbf{\texttt{src/\allowbreak generation/\allowbreak }}:

  \begin{itemize}
  \tightlist
  \item
    Read \texttt{prompt\_\allowbreak templates.\allowbreak py} to study the system prompts.
  \item
    Read \texttt{llm\_\allowbreak manager.\allowbreak py} to see multi-provider routing (Groq,
    OpenAI, Ollama).
  \item
    Read \texttt{study\_\allowbreak generator.\allowbreak py} to understand JSON parsing and
    study artifact creation.
  \end{itemize}
\item
  \textbf{\texttt{src/\allowbreak rag/\allowbreak rag\_\allowbreak pipeline.\allowbreak py}}: See how retrieval and
  generation combine into the core Q\&A pipeline.
\item
  \textbf{\texttt{app.\allowbreak py}}: Study the Streamlit UI, theme toggles,
  routing, and background thread workers.
\item
  \textbf{\texttt{tests/\allowbreak }}: Execute
  \texttt{USE\_\allowbreak TF=\allowbreak 0\ \allowbreak USE\_\allowbreak TORCH=\allowbreak 1\ \allowbreak pytest} and inspect test cases.
\end{enumerate}

\subsection{21.2 Practical Extension
Recipes}\label{practical-extension-recipes}

\subsubsection{How to Add a New Study Tool (e.g.~``Mindmap
Generator'')}\label{how-to-add-a-new-study-tool-e.g.-mindmap-generator}

\begin{enumerate}
\def\labelenumi{\arabic{enumi}.}
\tightlist
\item
  \textbf{Define Prompt}: In
  \texttt{src/\allowbreak generation/\allowbreak prompt\_\allowbreak templates.\allowbreak py}, add
  \texttt{get\_\allowbreak mindmap\_\allowbreak prompt\_\allowbreak template(\allowbreak )\allowbreak } returning Mermaid syntax
  (\texttt{graph\ \allowbreak TD\ \allowbreak .\allowbreak .\allowbreak .}).
\item
  \textbf{Implement Generator}: In
  \texttt{src/\allowbreak generation/\allowbreak study\_\allowbreak generator.\allowbreak py}, create
  \texttt{generate\_\allowbreak mindmap(\allowbreak topic,\allowbreak \ \allowbreak vector\_\allowbreak store,\allowbreak \ \allowbreak llm\_\allowbreak manager,\allowbreak \ \allowbreak source\_\allowbreak filter)\allowbreak }
  that retrieves context, invokes the chain, and extracts Mermaid
  syntax.
\item
  \textbf{Create Worker}: In \texttt{app.\allowbreak py}, define
  \texttt{worker\_\allowbreak generate\_\allowbreak mindmap(\allowbreak .\allowbreak .\allowbreak .)\allowbreak } that updates
  \texttt{GLOBAL\_\allowbreak BACKGROUND\_\allowbreak JOBS} and calls
  \texttt{add\_\allowbreak history\_\allowbreak entry(\allowbreak "mindmap",\allowbreak \ \allowbreak .\allowbreak .\allowbreak .)\allowbreak }.
\item
  \textbf{Add UI Screen}: In \texttt{app.\allowbreak py}, add
  \texttt{"{[}Embeddings\ \allowbreak Engine{]}\ \allowbreak \ \allowbreak Mindmaps"} to the
  \texttt{st.\allowbreak sidebar.\allowbreak radio} list and implement
  \texttt{render\_\allowbreak mindmap\_\allowbreak page(\allowbreak )\allowbreak }.
\item
  \textbf{Add Unit Test}: In \texttt{tests/\allowbreak }, add a unit test validating
  prompt formatting and mock execution.
\end{enumerate}

\subsubsection{How to Add a New LLM Provider (e.g.~Anthropic
Claude)}\label{how-to-add-a-new-llm-provider-e.g.-anthropic-claude}

\begin{enumerate}
\def\labelenumi{\arabic{enumi}.}
\tightlist
\item
  Add \texttt{langchain-anthropic} to \texttt{requirements.\allowbreak txt}.
\item
  In \texttt{src/\allowbreak generation/\allowbreak llm\_\allowbreak manager.\allowbreak py}, add an
  \texttt{elif\ \allowbreak self.\allowbreak provider\ \allowbreak =\allowbreak =\allowbreak \ \allowbreak "anthropic":} block initializing
  \texttt{ChatAnthropic(\allowbreak model=\allowbreak self.\allowbreak model\_\allowbreak name,\allowbreak \ \allowbreak api\_\allowbreak key=\allowbreak os.\allowbreak getenv(\allowbreak "ANTHROPIC\_\allowbreak API\_\allowbreak KEY")\allowbreak )\allowbreak }.
\item
  In \texttt{app.\allowbreak py}, add \texttt{"anthropic"} to the provider selection
  dropdown.
\item
  In \texttt{tests/\allowbreak test\_\allowbreak llm\_\allowbreak manager.\allowbreak py}, add a unit test verifying
  provider initialization.
\end{enumerate}

\section{Chapter 22: Glossary and
References}\label{chapter-22-glossary-and-references}

\subsection{22.1 Technical Glossary}\label{technical-glossary}

\begin{itemize}
\tightlist
\item
  \textbf{Active Recall}: The pedagogical practice of stimulating memory
  during the learning process by testing oneself on concepts rather than
  passively reading.
\item
  \textbf{Cosine Distance}: A metric measuring the angle between two
  multi-dimensional vectors, invariant to vector magnitude:
  \(1 - \cos(\theta)\).
\item
  \textbf{Dense Vector Embedding}: A low-dimensional numerical vector
  representation of text where semantically similar texts are placed
  close together in mathematical space.
\item
  \textbf{Euclidean (L2) Distance}: The straight-line distance between
  two vectors in Euclidean space: \(\|\mathbf{q} - \mathbf{d}\|_2\).
\item
  \textbf{FAISS (Facebook AI Similarity Search)}: An open-source C++
  library with Python bindings developed by Meta for high-speed
  similarity search on dense vectors.
\item
  \textbf{LCEL (LangChain Expression Language)}: A declarative
  composition framework for chaining LangChain components using
  Unix-pipe operators.
\item
  \textbf{PyMuPDF (\texttt{fitz})}: A high-performance Python binding
  for the MuPDF C library, providing fast PDF text and metadata
  extraction.
\item
  \textbf{RAG (Retrieval-Augmented Generation)}: An AI architecture that
  enhances large language model generation by retrieving authoritative
  factual context from external knowledge bases before prompting.
\item
  \textbf{Recursive Character Text Splitting}: An iterative text
  segmentation method that splits text using a priority cascade of
  separators to preserve semantic paragraph and sentence boundaries.
\end{itemize}

\subsection{22.2 References \& Authoritative
Documentation}\label{references-authoritative-documentation}

\begin{enumerate}
\def\labelenumi{\arabic{enumi}.}
\tightlist
\item
  \textbf{Lewis, P., et al.~(2020).} \emph{Retrieval-Augmented
  Generation for Knowledge-Intensive NLP Tasks.} Advances in Neural
  Information Processing Systems (NeurIPS 2020).
\item
  \textbf{Roediger, H. L., \& Karpicke, J. D. (2006).} \emph{The Power
  of Testing Memory: Basic Research and Implications for Educational
  Practice.} Perspectives on Psychological Science, 1(3), 181-210.
\item
  \textbf{LangChain Documentation:} https://python.langchain.com/
\item
  \textbf{PyMuPDF Documentation:} https://pymupdf.readthedocs.io/
\item
  \textbf{Sentence-Transformers Documentation:} https://www.sbert.net/
\item
  \textbf{FAISS GitHub Repository:}
  https://github.com/facebookresearch/faiss
\item
  \textbf{ChromaDB Documentation:} https://docs.trychroma.com/
\item
  \textbf{Groq Cloud API Documentation:} https://console.groq.com/docs
\item
  \textbf{Streamlit Documentation:} https://docs.streamlit.io/
\end{enumerate}

\section{Chapter 23: Final Project
Summary}\label{chapter-23-final-project-summary}

\subsection{23.1 Executive Wrap-Up}\label{executive-wrap-up}

\textbf{StudyRAG} successfully addresses the fundamental weaknesses of
applying generic AI chatbots to academic studying. By grounding all
interactions in the student's exact uploaded course materials, the
application provides an educational environment that combines: -
\textbf{Verifiable Truth}: Every statement is backed by direct document
and page citations. - \textbf{Cognitive Scaffolding}: Automated
generation of multiple-choice quizzes, flashcards, and chapter revision
guides promotes active recall and spaced retrieval. -
\textbf{Engineering Excellence}: Sub-second vector search, asynchronous
background threading, multi-provider model routing, atomic persistence,
and a modern student-friendly interface. - \textbf{Production Rigor}:
Comprehensive automated test coverage (28 passing tests) guaranteeing
reliability across all critical components.

\subsection{23.2 Current Verification
Status}\label{current-verification-status}

\begin{itemize}
\tightlist
\item
  \textbf{Automated Tests:} 28 passed out of 28 (\texttt{pytest} exit
  code 0).
\item
  \textbf{Core Pipelines:} PDF Ingestion, Chunking, Embeddings,
  FAISS/Chroma Vector Storage, Retrieval Filtering, LLM Generation, and
  Study History verified.
\item
  \textbf{Environment:} Tested and validated on Python 3.12.2 on macOS
  (Darwin arm64).
\end{itemize}

\subsection{23.3 Immediate Recommended Next
Actions}\label{immediate-recommended-next-actions}

\begin{enumerate}
\def\labelenumi{\arabic{enumi}.}
\tightlist
\item
  Launch the application: \texttt{streamlit\ \allowbreak run\ \allowbreak app.\allowbreak py}.
\item
  Ingest your course textbooks or lecture slide PDFs in \textbf{{[}Study
  Library{]} Study Library}.
\item
  Engage with \textbf{{[}Ask Questions{]} Ask Questions},
  \textbf{{[}Practice Quiz{]} Practice Quiz}, and
  \textbf{{[}Flashcards{]} Flashcards} for active recall study sessions.
\item
  Maintain test integrity during feature development by executing
  \texttt{USE\_\allowbreak TF=\allowbreak 0\ \allowbreak USE\_\allowbreak TORCH=\allowbreak 1\ \allowbreak pytest}.
\end{enumerate}

\begin{center}\rule{0.5\linewidth}{0.5pt}\end{center}

\emph{End of Master Documentation --- StudyRAG Version 1.0.0 (Release
Candidate)}

\end{document}
